\chapter{Equations with prescribed singularities}
\label{chap:prescribed-equations}

Prescribing a singularity type asks for a solution whose difference from a given potential is bounded. An existence proof must therefore control this difference as well as the Monge--Amp\`ere measure. A strict lower barrier makes comparison possible; relative estimates keep the approximants in the required type; envelope compactness then identifies their limiting equation.

We study both the exponential equation and its zero-exponent limit. A positive exponent fixes the additive constant and gives uniqueness in the prescribed type. At exponent zero, comparison still determines the scalar multiplying the density, while the potential can fail to be unique. The projective example shows how symmetry can nevertheless select the limit of the exponential solutions. For a general big background, a strict analytic barrier also identifies the open set on which stability leads to local regularity.

For integrable densities, the final part of the chapter replaces an $L^p$ estimate by a subsolution threshold. Bounded truncations and Gauduchon tests determine this threshold and construct a greatest exponential solution when it is positive. The argument distinguishes convergence of the scalar and its measures from convergence of the normalized potentials. The existence theory for all finite non-pluripolar measures under mass hypotheses is developed in \cref{chap:relative-finite-measures,chap:bounded-mass-equations}.

Throughout, $X$ is a connected compact complex manifold of dimension $n\geq1$, $\omega$ is a Hermitian metric, and $\mathrm{d}V$ is a smooth positive probability volume form. Unqualified function norms use this volume. We use the singularity order, envelopes and capacity of \cref{chap:framework}. We abbreviate $\mathrm{MA}^{\mathrm{ext}}_\theta(u)$ to $\mathrm{MA}_\theta(u)$. The background hypotheses are specified in each argument, and the extended measures may have infinite mass. Write $s_+=\max(s,0)$ for a real $s$.

\section{Domination with a smaller background}

Throughout this section $\theta$ is semipositive and big. Let $0<\delta<1$. When a singular potential lies above a $\delta\theta$-plurisubharmonic function, the gap between them can be measured by an envelope in the remaining background $(1-\delta)\theta$. This background is again semipositive and big, so its bounded potentials have Monge--Amp\`ere measures of positive mass by \cref{prop:big-noncollapsing}, and bounded domination applies to them. It bounds the envelope from below. A family of stretched envelopes then converts a strict measure inequality into a pointwise inequality, as in the bounded case of \cref{thm:semipositive-bounded-domination}.

We say that $u\in\PSH(X,\theta)$ \emph{admits a smaller-background envelope} if $P_{\delta\theta}(u)\not\equiv-\infty$ for some $0<\delta<1$. This property depends only on the ordinary singularity type of $u$: if $u'\geq u-C$, then $P_{\delta\theta}(u)-C$ is a competitor for $P_{\delta\theta}(u')$.

\Needspace{16\baselineskip}
\begin{lemma}\label{lem:singular-zero-domination}
Suppose $u,v\in\PSH(X,\theta)$, $v\preceq u$, and
$\mathrm{MA}^{\mathrm{ext}}_\theta(u)(\{u<v\})=0$. If $0<\delta<1$ and $v_\delta=P_{\delta\theta}(v)\not\equiv-\infty$, then $u\geq v_\delta$.
\end{lemma}
\begin{proof}
Choose $C$ with $v-C\leq u$. Since $v_\delta\leq v$, the obstacle $u-v_\delta$ is bounded below by $-C$ off the pole sets, and it is finite almost everywhere. Hence the envelope
\[
 p=P_{(1-\delta)\theta}(u-v_\delta)
\]
is bounded by \cref{lem:bounded-below-envelope}, and by contact support (\cref{thm:envelope-contact}) its measure is carried by $D=\{p+v_\delta=u\}$. Moreover $p+v_\delta\leq u$ everywhere, and
\[
 \theta+\ddc(p+v_\delta)
 =((1-\delta)\theta+\ddc p)+(\delta\theta+\ddc v_\delta)
\]
is a sum of two positive currents. Positivity of mixed products and ordered contact therefore give
\[
 \mathbf1_D\mathrm{MA}_{(1-\delta)\theta}(p)
 \leq\mathbf1_D\mathrm{MA}^{\mathrm{ext}}_\theta(u);
\]
this inequality is first proved on common finite truncation levels and then extended by exhaustion. Off the pole sets, every point of $D\cap\{p<0\}$ satisfies
$u=p+v_\delta<v_\delta\leq v$. The hypothesis therefore gives
$\mathrm{MA}_{(1-\delta)\theta}(p)(\{p<0\})=0$. Since $(1-\delta)\theta$ is semipositive and big, bounded domination (\cref{thm:semipositive-bounded-domination}, with \cref{prop:big-noncollapsing}), applied to $p$ and the constant $0$ with coefficient zero, yields $p\geq0$. Hence $u\geq p+v_\delta\geq v_\delta$.
\end{proof}

\begin{theorem}\label{thm:singular-strict-domination}\index{Domination principle!singular potentials}
Let $u\in\PSH(X,\theta)$ satisfy $P_{\delta\theta}(u)\not\equiv-\infty$ for some $0<\delta<1$. Suppose $v\in\PSH(X,\theta)$, $v\preceq u$, and $0\leq c<1$. If
\[
 \mathbf1_{\{u<v\}}\mathrm{MA}^{\mathrm{ext}}_\theta(u)
 \leq c\,\mathbf1_{\{u<v\}}\mathrm{MA}^{\mathrm{ext}}_\theta(v),
\]
then $u\geq v$.
\end{theorem}
\begin{proof}
Since $\max(u,v)=v$ on the plurifine open set $\{u<v\}$, plurifine locality shows that replacing $v$ by $\max(u,v)$ preserves the hypothesis; it does not change the conclusion. We may therefore assume $u\leq v$, and since $v\preceq u$ there is $C\geq0$ with
$u\leq v\leq u+C$. For $b>1$ consider the stretched envelopes
\[
 p_b=P_\theta(bu-(b-1)v),\qquad
 D_b=\{p_b=bu-(b-1)v\}.
\]
The function $u-(b-1)C$ is a competitor, so $p_b\not\equiv-\infty$; and since $bu-(b-1)v\leq u$, we have
$u-(b-1)C\leq p_b\leq u$. By contact support, the measure of $p_b$ is concentrated on $D_b$. The potential
$b^{-1}p_b+(1-b^{-1})v$ lies below $u$ and agrees with it on $D_b$, so ordered contact and positivity of mixed products give
\[
 \mathbf1_{D_b}\bigl[b^{-n}\mathrm{MA}^{\mathrm{ext}}_\theta(p_b)
 +(1-b^{-1})^n\mathrm{MA}^{\mathrm{ext}}_\theta(v)\bigr]
 \leq\mathbf1_{D_b}\mathrm{MA}^{\mathrm{ext}}_\theta(u).
\]
Choose $b$ so large that $(1-b^{-1})^n>c$. On $D_b\cap\{u<v\}$ the hypothesis then forces $\mathrm{MA}^{\mathrm{ext}}_\theta(p_b)$ to vanish. Since the masses may be infinite, we justify this subtraction as follows. On each set where $u,v,p_b>-k$, every measure involved is finite, being bounded by the mass of the corresponding bounded truncation. We subtract the finite terms on each such set and then take the union over $k$, so no infinite masses are ever subtracted.

It follows that the measure of $p_b$ is carried by $D_b\cap\{u=v\}\subset\{p_b=u\}$. We apply \cref{lem:singular-zero-domination} to the pair $(p_b,u)$: their singularities are equivalent, $\mathrm{MA}^{\mathrm{ext}}_\theta(p_b)(\{p_b<u\})=0$, and $P_{\delta\theta}(u)\not\equiv-\infty$. We obtain
\[
 p_b\geq P_{\delta\theta}(u)
\]
for all sufficiently large $b$. Suppose now that $u\geq v$ fails. By \cref{lem:qpsh-almost-everywhere-order}, the set $\{u<v\}$ has positive volume. Since both potentials are finite almost everywhere, some set $E=\{u<v-a\}$ with $a>0$ has positive volume and both functions are finite on $E$. On $E$ the obstacle inequality gives
\[
 P_{\delta\theta}(u)\leq p_b\leq v-ba
\]
quasi-everywhere. Letting $b$ tend to infinity through the integers shows that the quasi-plurisubharmonic function $P_{\delta\theta}(u)\not\equiv-\infty$ equals $-\infty$ on a set of positive volume, a contradiction. Thus $u\geq v$ off the pole sets, and \cref{lem:qpsh-almost-everywhere-order} gives the inequality everywhere.
\end{proof}

\begin{corollary}\label{cor:singular-exponential-order}
Let $u,v\in\PSH(X,\theta)$ satisfy $v\preceq u$, and suppose that $u$ admits a smaller-background envelope. If $\lambda>0$ and
\[
 \mathrm{e}^{-\lambda u}\mathrm{MA}^{\mathrm{ext}}_\theta(u)
 \leq \mathrm{e}^{-\lambda v}\mathrm{MA}^{\mathrm{ext}}_\theta(v)
\]
off the pole sets, then $u\geq v$. In particular, for every $\lambda>0$ and every positive measure $\mu$, the equation $\mathrm{MA}_\theta(u)=\mathrm{e}^{\lambda u}\mu$ has at most one solution in an ordinary singularity type admitting a smaller-background envelope.
\end{corollary}
\begin{proof}
Fix $t>0$. On $\{u<v-t\}$ the assumed inequality gives
$\mathrm{MA}_\theta^{\mathrm{ext}}(u)\leq \mathrm{e}^{-\lambda t}\mathrm{MA}_\theta^{\mathrm{ext}}(v-t)$.
The theorem, applied to $u$ and $v-t$ with $c=\mathrm{e}^{-\lambda t}$, gives $u\geq v-t$; now let $t\downarrow0$. For two solutions of the same equation in the stated type, apply this result in both directions. The weighted measures are compared off the common pole sets, which these measures do not charge, so no value of an infinite weight is needed.
\end{proof}


\Needspace{12\baselineskip}
\section{Bounded equations in a semipositive big background}

In this section $\theta$ is again semipositive and big, and we solve the exponential and the normalized equations with $L^p$ densities by bounded potentials. The basic tool is a bounded subsolution. A strict barrier, corrected by the solution of a Hermitian equation, is composed with a bounded convex function; the result is a bounded $\theta$-plurisubharmonic function whose Monge--Amp\`ere measure dominates a prescribed density. The derivative of the convex function decays polynomially near the poles of the barrier, and we compensate for this decay by multiplying the auxiliary density by the reciprocal weight.

Fix a strict barrier $\rho\in\PSH(X,\theta)$ and $0<a\leq1$ such that
\[
 \theta+\ddc\rho\geq a\omega,\qquad \rho\leq-1.
\]
No property of $\rho$ other than these two inequalities will be used.

\Needspace{18\baselineskip}
\begin{lemma}\label{lem:big-lp-subsolution}
Fix $p>1$ and $L>0$. There is $b>0$ such that, for every $f\geq0$ with $\|f\|_{L^p(\mathrm{d}V)}\leq L$, one can find $s\in\PSH(X,\theta)$ satisfying
\[
 -1\leq s\leq0,\qquad \mathrm{MA}_\theta(s)\geq b f\,\mathrm{d}V.
\]
The constant $b$ depends only on the fixed geometry and on $\rho,a,p,L$.
\end{lemma}
\begin{proof}
Choose $1<q<p$. Since a fixed quasi-plurisubharmonic function belongs to every finite $L^r$ space, H\"older's inequality gives a uniform $L^q$ bound for
\[
 g=(-\rho)^{2n}f.
\]
By \cref{thm:kn-continuous} there are $c>0$ and a continuous $\omega$-plurisubharmonic $v$, normalized by $\sup_Xv=-1$, such that
\[
 \omega_v^n=c(1+g)\,\mathrm{d}V.
\]
The construction in that theorem gives a uniform upper bound for $c$, because in its Gauduchon trace estimate the root integral of $1+g$ is bounded below by the root integral of $1$. The lower-density estimate of \cref{thm:lp-potential-estimate} gives
$c\|1+g\|_q\geq c_q>0$, so $c$ is also bounded below. The potential estimate of the same theorem then gives $-C\leq v\leq-1$ with a uniform $C$.

Put $h=\rho+av\leq-1$. Then $\theta_h\geq a\omega_v$. On $(-\infty,-1]$ define
\[
 \chi(t)=-t^{-1}-1,\qquad
 \chi'(t)=t^{-2}\in(0,1],\qquad \chi''(t)=-2t^{-3}>0,
\]
and set $s=\chi(h)$, with value $-1$ at the poles of $h$. To justify both the positivity of $\theta+\ddc s$ and the computation of its Monge--Amp\`ere measure, we first work with the bounded functions $h_k=\max(h,-k)$, $k\geq1$. The chain rule for bounded potentials and $\theta\geq0$ give
\[
 \theta+\ddc\chi(h_k)
 =(1-\chi'(h_k))\theta+\chi'(h_k)\theta_{h_k}
       +\chi''(h_k)\,\mathrm{d}h_k\wedge \dc h_k\geq0.
\]
These bounded potentials decrease to $s$, so $s\in\PSH(X,\theta)$ and $-1\leq s\leq0$. On $\{h>-k\}$, locality and positivity of mixed products give
\[
 \mathrm{MA}_\theta(s)
 \geq a^n|h|^{-2n}\omega_v^n
 \geq a^nc|h|^{-2n}(-\rho)^{2n}f\,\mathrm{d}V.
\]
Letting $k\to\infty$, we obtain this inequality on the set where $h$ is finite; its complement is pluripolar and has zero volume. Since
$|h|\leq(-\rho)+aC\leq(1+aC)(-\rho)$, the assertion follows with any common lower bound for $a^nc(1+aC)^{-2n}$.
\end{proof}

The lemma gives a subsolution for each density in an $L^p$-bounded family. It does not give a lower bound for the masses of all bounded $\theta$-plurisubharmonic functions.

\Needspace{22\baselineskip}
\subsection{The exponential equation}

We approximate $\theta$ from above by the Hermitian forms $\theta+\varepsilon\omega$. The bounded subsolution provides a common lower barrier for the corresponding solutions, and exponential comparison shows that they decrease as $\varepsilon$ decreases. Their limit solves the equation in the original background.

\begin{theorem}\label{thm:big-bounded-exponential}
Let $0\leq f\in L^p(\mathrm{d}V)$, $p>1$, be nonzero, and let $\lambda>0$. There is a unique bounded $v\in\PSH(X,\theta)$ such that
\[
 \mathrm{MA}_\theta(v)=\mathrm{e}^{\lambda v}f\,\mathrm{d}V.
\]
For fixed $\lambda$, these solutions are uniformly bounded when $\|f\|_p$ is bounded above and $\int_X f^{1/n}\,\mathrm{d}V$ is bounded away from zero.
\end{theorem}
\begin{proof}
For $0<\varepsilon\leq1$, set $\theta_\varepsilon=\theta+\varepsilon\omega$. By the Hermitian existence theorem \cref{thm:nguyen-exp}, there is a continuous solution $v_\varepsilon$ of
\[
 \mathrm{MA}_{\theta_\varepsilon}(v_\varepsilon)
       =\mathrm{e}^{\lambda v_\varepsilon}f\,\mathrm{d}V.
\]
Take $s,b$ from \cref{lem:big-lp-subsolution}, and put
$w=s+\lambda^{-1}\log b$. Positivity of bounded products gives
\[
 \mathrm{MA}_{\theta_\varepsilon}(w)
 \geq\mathrm{MA}_\theta(s)\geq bf\,\mathrm{d}V
 \geq \mathrm{e}^{\lambda w}f\,\mathrm{d}V,
\]
and bounded exponential comparison (\cref{cor:bounded-exponential-comparison}) yields $w\leq v_\varepsilon$.
If $0<\varepsilon'<\varepsilon$, then $v_{\varepsilon'}$ is also $\theta_\varepsilon$-plurisubharmonic and
\[
 \mathrm{MA}_{\theta_\varepsilon}(v_{\varepsilon'})
 \geq \mathrm{e}^{\lambda v_{\varepsilon'}}f\,\mathrm{d}V,
\]
so the same comparison gives $v_{\varepsilon'}\leq v_\varepsilon$. Hence $v_\varepsilon$ decreases to a bounded $\theta$-plurisubharmonic function $v$ with $w\leq v\leq v_1$.

To pass to the limit in the equation, we use the fixed Hermitian background $\Omega=\theta+\omega$ and the expansion
\[
 (\theta_\varepsilon+\ddc v_\varepsilon)^n
 =\sum_{k=0}^n\binom nk
       (\Omega+\ddc v_\varepsilon)^k
       \wedge\bigl(-(1-\varepsilon)\omega\bigr)^{n-k}.
\]
Continuity of bounded products along monotone sequences applies to every term, and the smooth coefficients converge. Dominated convergence applies to the right-hand sides. This proves the equation. Uniqueness follows from bounded domination in the semipositive big background (\cref{thm:semipositive-bounded-domination,prop:big-noncollapsing}) by the usual exponential comparison argument.

For the uniform assertion, note first that the lower barrier $w$ is uniform by \cref{lem:big-lp-subsolution}. The upper barrier $v_1$ solves an equation in the fixed Hermitian background $\Omega$. The proof of \cref{lem:exponential-gauduchon-mean} bounds its supremum using only an upper bound for the density norm and a positive lower bound for the root integral: these two bounds control, respectively, the H\"older term and the denominator in \eqref{eq:gauduchon-exponential-mean}, and they survive the smooth approximation used there. Thus $v_1$, and hence $v$, has a uniform upper bound.
\end{proof}

Local continuity of this bounded solution will follow from \cref{thm:local-big-bounded-continuity}.

\subsection{The normalized equation}

For the normalized equation the approximating potentials need not be monotone. Upper envelopes and lower envelopes of their tails provide monotone sequences on either side of an $L^1$ cluster limit. Comparison identifies the two limits, and hence the equation, while the constant is controlled by the Gauduchon trace and the bounded subsolution.

\Needspace{14\baselineskip}
\begin{theorem}\label{thm:big-bounded-normalized}
For every nonzero $0\leq f\in L^p(\mathrm{d}V)$, $p>1$, there are a bounded $\varphi\in\PSH(X,\theta)$ and a uniquely determined $c>0$ such that
\[
 \mathrm{MA}_\theta(\varphi)=cf\,\mathrm{d}V,\qquad \sup_X\varphi=0.
\]
The construction gives uniform bounds for $\|\varphi\|_\infty$, $c$ and $c^{-1}$ when $\|f\|_p$ is bounded above and $\int_Xf^{1/n}\,\mathrm{d}V$ is bounded away from zero.
\end{theorem}

Only the constant $c$ is asserted to be unique; the normalized potential need not be unique, as \cref{thm:degenerate-hermitian-nonuniqueness} shows already for a Hermitian background.

\begin{proof}
\emph{Step 1: approximating equations and their constants.} For integers $j\geq1$, put $\omega_j=\theta+2^{-j}\omega$, and use \cref{thm:kn-continuous}, with the Hermitian form $\omega_j$ in place of $\omega$, to choose continuous normalized solutions
\[
 \mathrm{MA}_{\omega_j}(\varphi_j)=c_jf\,\mathrm{d}V,
 \qquad \sup_X\varphi_j=0.
\]
We first bound $c_j$. Fix a Gauduchon metric $\gamma_0$; when $n=1$ any Hermitian form will do, and the trace estimate reduces to Stokes' theorem. The smooth approximation in the proof of \cref{thm:kn-continuous}, together with the Gauduchon trace, gives a common upper bound: the numerator $\int_X\omega_j\wedge\gamma_0^{n-1}$ is bounded above by $\int_X(\theta+\omega)\wedge\gamma_0^{n-1}$, and the root integral of $f$ is positive and fixed. Here the estimate is applied to the smooth approximating equations before passing to their continuous limits.

For a lower bound, take $s,b$ from \cref{lem:big-lp-subsolution}. If $c_j<b$, then
\[
 \mathrm{MA}_{\omega_j}(\varphi_j)
 \leq(c_j/b)\mathrm{MA}_{\omega_j}(s+A)
\]
for every constant $A$, and strict domination in the Hermitian background $\omega_j$ (\cref{thm:bounded-strict-domination}) would give $\varphi_j\geq s+A$ for every $A$, which is absurd. Thus $c_j\geq b$. In fact the constants decrease with $j$: if $k>j$, positivity gives
$\mathrm{MA}_{\omega_j}(\varphi_k)\geq c_kf\,\mathrm{d}V$, and the same argument excludes $c_j<c_k$. Hence $c_j\downarrow c>0$.

\emph{Step 2: a uniform bound for the potentials.} For the Hermitian form $\Omega=\theta+\omega$, all $\varphi_j$ are normalized $\Omega$-plurisubharmonic functions. Using uniform exponential integrability (\cref{lem:uniform-exponential}) and H\"older's inequality, choose $\gamma>0$ small and $1<q<p$ such that
\[
 g_j=\mathrm{e}^{-\gamma\varphi_j}f
\]
is bounded in $L^q$. Since $\varphi_j\leq0$, we have $g_j\geq f$. \Cref{thm:big-bounded-exponential}, applied with the exponent $q$, therefore gives uniformly bounded $w_j\in\PSH(X,\theta)$ with
\[
 \mathrm{MA}_\theta(w_j)=\mathrm{e}^{\gamma w_j}g_j\,\mathrm{d}V.
\]
Choose $C_0\geq\max\{1,\sup_jc_j\}$ and set $z_j=w_j-\gamma^{-1}\log C_0$. Then
\[
 \mathrm{MA}_{\omega_j}(z_j)
 \geq \mathrm{e}^{\gamma z_j}c_jg_j\,\mathrm{d}V,
 \qquad
 \mathrm{MA}_{\omega_j}(\varphi_j)
       =\mathrm{e}^{\gamma\varphi_j}c_jg_j\,\mathrm{d}V,
\]
and exponential comparison gives $z_j\leq\varphi_j\leq0$. Thus all $\varphi_j$ take values in a common bounded interval.

\emph{Step 3: the upper envelopes.} Extract a subsequence converging in $L^1$ and almost everywhere to the canonical $\theta$-plurisubharmonic representative $\varphi$. We keep the indices $j$; the coefficients of $\omega$ in $\omega_j$ still decrease to zero. Define
\[
 \Phi_j^+=\left(\sup_{\ell\geq j}\varphi_\ell\right)^*,
 \qquad h_j=\inf_{\ell\geq j}\varphi_\ell,
 \qquad p_j=P_\theta(h_j).
\]
All these functions take values in a common bounded interval. The upper envelopes are $\omega_j$-plurisubharmonic and decrease to $\varphi$, since the regularized upper envelope agrees almost everywhere with the $L^1$ limit; this also follows from the local Hartogs lemma applied to the bounded tails.

For $\ell\geq j$, positivity gives $\mathrm{MA}_{\omega_j}(\varphi_\ell)\geq c_\ell f\,\mathrm{d}V\geq cf\,\mathrm{d}V$. By ordered contact (\cref{lem:bounded-ordered-contact}), finite maxima of these functions satisfy the same lower bound. Continuity of bounded products along increasing sequences, followed by the decreasing limit in $j$ and the expansion in the smooth background, gives
\begin{equation}\label{eq:big-normalized-upper-tail}
 \mathrm{MA}_{\omega_j}(\Phi_j^+)\geq cf\,\mathrm{d}V,
 \qquad \mathrm{MA}_\theta(\varphi)\geq cf\,\mathrm{d}V.
\end{equation}

\emph{Step 4: the lower envelopes.} For the reverse inequality we use lower envelopes, which carry a useful exponential weight. For $N\geq j$ put
\[
 p_{j,N}=P_\theta\left(\min_{j\leq\ell\leq N}\varphi_\ell\right).
\]
The obstacle is continuous and bounded, so the measure of $p_{j,N}$ is carried by its contact set. Partition this set into disjoint Borel pieces according to the first index $\ell$ at which $p_{j,N}=\varphi_\ell$. On such a piece, positivity and ordered contact in the background $\omega_\ell$ give
\[
 \mathrm{MA}_\theta(p_{j,N})
 \leq\mathrm{MA}_{\omega_\ell}(p_{j,N})
 \leq c_\ell f\,\mathrm{d}V.
\]
Moreover $p_{j,N}=\varphi_\ell\geq h_j$ there, and $c_\ell\leq c_j$. Summing over the pieces therefore yields the global inequality
\[
 \mathrm{MA}_\theta(p_{j,N})
 \leq \mathrm{e}^{p_{j,N}-h_j}c_jf\,\mathrm{d}V.
\]
As $N\to\infty$, the envelopes decrease to $p_j$. Indeed, their limit lies below each tail obstacle quasi-everywhere and is thus a competitor for $P_\theta(h_j)$, while the reverse inequality is immediate. Continuity of bounded products along monotone sequences and dominated convergence give
\[
 \mathrm{MA}_\theta(p_j)\leq \mathrm{e}^{p_j-h_j}c_jf\,\mathrm{d}V.
\]
The functions $p_j$ increase. Let $\varphi'=(\lim_jp_j)^*$. Since $h_j\to\varphi$ almost everywhere, we have $\varphi'\leq\varphi$ for the canonical representatives. Passing to the increasing limit and using \eqref{eq:big-normalized-upper-tail}, we obtain
\[
 \mathrm{MA}_\theta(\varphi')
 \leq \mathrm{e}^{\varphi'-\varphi}cf\,\mathrm{d}V
 \leq \mathrm{e}^{\varphi'-\varphi}\mathrm{MA}_\theta(\varphi).
\]
Exponential comparison for bounded potentials in the semipositive big background $\theta$, which follows from \cref{thm:semipositive-bounded-domination,prop:big-noncollapsing} as in \cref{cor:bounded-exponential-comparison}, forces $\varphi'\geq\varphi$. Thus $\varphi'=\varphi$, and the first inequality of the last display, together with \eqref{eq:big-normalized-upper-tail}, proves the equation. If necessary, subtracting $\sup_X\varphi$ gives the required normalization.

\emph{Step 5: uniformity and uniqueness of $c$.} All bounds used above depend only on the two stated density bounds and the fixed data, which proves the uniform assertion for the constructed solution. Finally, two different positive constants for the same $f$ would contradict strict domination after an arbitrary translation of one of the bounded solutions, exactly as in the lower bound for $c_j$. This proves the uniqueness of $c$.
\end{proof}

For $\lambda>0$, let $v_\lambda$ be the bounded solution of $\mathrm{MA}_\theta(v_\lambda)=\mathrm{e}^{\lambda v_\lambda}f\,\mathrm{d}V$, and let $(\varphi,c)$ be a normalized solution given by \cref{thm:big-bounded-normalized}. The two exponential barriers used in \cref{prop:vanishing-exponent} give
\[
 \|\lambda v_\lambda-\log c\|_\infty
 \leq\lambda\operatorname{osc}_X\varphi
\]
also in this semipositive background.


\section{A quantitative stability estimate}
\label{sec:bounded-stability}

If a $\theta$-plurisubharmonic function exceeds a bounded solution only by a small amount in $L^1$, the excess is also uniformly small. To prove this, we restrict the density of the solution to a strict superlevel set of the difference and construct a subsolution for the restricted density. Domination bounds the height of the superlevel set in terms of the norm of the restricted density, and H\"older's inequality converts its volume into a power of the $L^1$ error.

Throughout this section $\theta$ is semipositive and big, and $p>1$ is fixed.

\Needspace{20\baselineskip}
\begin{theorem}\label{thm:bounded-one-sided-stability}\index{Stability!bounded potentials}
Fix $L,C_0>0$. Let $u\in\PSH(X,\theta)$ satisfy
\[
 -C_0\leq u\leq0,\qquad
 \mathrm{MA}_\theta(u)=f\,\mathrm{d}V,\qquad
 0\leq f\in L^p(\mathrm{d}V),\quad\|f\|_p\leq L.
\]
For every
\[
 0<\eta<\frac{p-1}{np+p-1}
\]
there is a constant $C$, depending only on the fixed data and $\eta$, such that every $v\in\PSH(X,\theta)$ with $v\leq0$ satisfies
\[
 \sup_X(v-u)_+
       \leq C\left(\int_X(v-u)_+\,\mathrm{d}V\right)^\eta.
\]
\end{theorem}

The comparison potential $v$ need not be bounded.

\begin{proof}
Replacing $v$ by $\max(v,u)$ leaves both sides unchanged, so we may assume
$-C_0\leq u\leq v\leq0$. Put $I=\int_X(v-u)\,\mathrm{d}V$. If $I=0$, the two canonical potentials agree everywhere and there is nothing to prove, so we assume $I>0$.

\Needspace{5\baselineskip}
Choose $q\in(1,p)$, to be specified at the end. For $\delta>0$ put
\[
 E_\delta=\{u<v-\delta\},\qquad
 g_\delta=\mathbf1_{E_\delta}f,\qquad
 m_\delta=\|g_\delta\|_q.
\]
If $m_\delta>0$, apply \cref{lem:big-lp-subsolution} to $g_\delta/m_\delta$. It gives $h\in\PSH(X,\theta)$ and $b>0$ such that
\[
 -1\leq h\leq0,\qquad
 \mathrm{MA}_\theta(h)\geq b\,g_\delta/m_\delta\,\mathrm{d}V.
\]
The constant $b$ does not depend on $\delta,f,u,v$, since the auxiliary density has $L^q$ norm one.

First suppose $0<m_\delta<b/2$, and take
\[
 t=(2m_\delta/b)^{1/n}\in(0,1),\qquad
 w=(1-t)v+th-\delta-C_0t.
\]
Because $h-v\leq C_0$, we have $w\leq v-\delta$, and hence
$\{u<w\}\subset E_\delta$. Positivity of bounded mixed products gives
\[
 \mathrm{MA}_\theta(w)\geq t^n\mathrm{MA}_\theta(h)
       \geq2g_\delta\,\mathrm{d}V.
\]
On $\{u<w\}$ the last measure equals $2\mathrm{MA}_\theta(u)$. Bounded strict domination in the semipositive big background (\cref{thm:semipositive-bounded-domination}) therefore implies $u\geq w$. Consequently
\[
 v-u\leq v-w
      =\delta+C_0t+t(v-h)\leq\delta+(C_0+1)t.
\]
If $m_\delta=0$, the measure of $u$ vanishes on $E_\delta$, and the same domination theorem with coefficient zero, applied to $u$ and $v-\delta$, gives $v-u\leq\delta$. If $m_\delta\geq b/2$, the bound $v-u\leq C_0$ gives the following inequality after an increase of its constant. Thus in every case
\begin{equation}\label{eq:stability-small-set}
 \sup_X(v-u)\leq\delta+C_1m_\delta^{1/n}.
\end{equation}

H\"older's inequality and Chebyshev's inequality give
\[
 m_\delta\leq\|f\|_p\,\mathrm{d}V(E_\delta)^{1/q-1/p}
       \leq L(I/\delta)^{(p-q)/(pq)}.
\]
Set $a=(p-q)/(npq)>0$. For $0<I\leq1$, choose
$\delta=I^{a/(1+a)}$. Substituting in \eqref{eq:stability-small-set} yields
\[
 \sup_X(v-u)\leq C_2 I^{\eta_q},
 \qquad \eta_q=\frac{p-q}{npq+p-q}.
\]
As $q\downarrow1$, the exponent $\eta_q$ tends to $(p-1)/(np+p-1)$. For the prescribed $\eta$ we choose $q>1$ so close to $1$ that $\eta_q>\eta$; then $I^{\eta_q}\leq I^\eta$ on $(0,1]$. When $I>1$, the bound $v-u\leq C_0$ gives the same conclusion after increasing $C$. This proves the theorem.
\end{proof}

\Needspace{14\baselineskip}
\begin{corollary}\label{cor:bounded-stability-uniform-convergence}
Fix $C_0>0$. Let $u_j,u\in\PSH(X,\theta)$ satisfy $-C_0\leq u_j,u\leq0$ and
\[
 \mathrm{MA}_\theta(u_j)=f_j\,\mathrm{d}V,\qquad
 \mathrm{MA}_\theta(u)=f\,\mathrm{d}V,
 \qquad \sup_j\|f_j\|_p+\|f\|_p<+\infty.
\]
If $u_j\to u$ in $L^1$, then $u_j\to u$ uniformly. More precisely, for every exponent $\eta$ as in the theorem,
\[
 \|u_j-u\|_\infty\leq C\|u_j-u\|_{L^1(\mathrm{d}V)}^\eta.
\]
\end{corollary}
\begin{proof}
Apply the theorem first with solution $u$ and comparison potential $u_j$, and then with solution $u_j$ and comparison potential $u$. The constants are uniform because of the common range and the common density bound. The larger of the two suprema is the uniform norm, and each integral of a positive part is at most the full $L^1$ difference.
\end{proof}

The corollary assumes that the limit $u$ already solves an equation with an $L^p$ density. In an existence proof, this must be established before the corollary can be applied.


\section{H\"older continuity from bounded stability}
\label{sec:holder-from-stability}

Apply bounded stability to the Chern-connection regularizations of a solution. Their $L^1$ error controls the pointwise smoothing error; a gradient estimate for the smooth averages then gives a H\"older modulus. We use the operators $R_t$ and transforms $\Phi_{c,\delta}$ from \cref{lem:chern-kiselman-input,lem:chern-smoothing-integral}. The final interpolation step yields the explicit exponent below.

\Needspace{16\baselineskip}
\begin{theorem}\label{thm:hermitian-holder-elementary}
Let $p>1$ and let $u\in\PSH(X,\omega)\cap L^\infty(X)$ solve
\[
 \mathrm{MA}_\omega(u)=f\,\mathrm{d}V,\qquad 0\leq f\in L^p(\mathrm{d}V).
\]
Then $u$ is H\"older continuous with every exponent
\[
 0<\beta<\frac{2(p-1)}{np+3(p-1)}.
\]
For a fixed bound $-C_0\leq u\leq0$ and a fixed bound for $\|f\|_p$, the H\"older constants are uniform. In particular the continuous solutions constructed in \cref{thm:kn-continuous} have this regularity.
\end{theorem}
\begin{proof}
Normalize by a constant so that $-C_0\leq u\leq0$. Choose
\[
 0<\eta<\frac{p-1}{np+p-1},\qquad \alpha=2\eta<1.
\]
\Needspace{5\baselineskip}
For small $\delta>0$ set
\[
 c=\delta^\alpha,\qquad e=Ac+2K\delta,\qquad
 r=-C_0-1,\qquad
 v_\delta=\frac{\Phi_{c,\delta}+er}{1+e}.
\]
The curvature estimate gives $v_\delta\in\PSH(X,\omega)$. Evaluation of the infimum at $t=\delta$ gives $\Phi_{c,\delta}\leq R_\delta u\leq0$. Its lower bound is
\[
 \Phi_{c,\delta}\geq u-K(\delta^2+\delta),
\]
since $R_tu+Kt^2\geq u$, $Kt\geq0$ and $-c\log(t/\delta)\geq0$. For $\delta$ small this is at least $r$. Thus
\[
 v_\delta\leq\Phi_{c,\delta}\leq R_\delta u\leq0.
\]
Apply \cref{thm:bounded-one-sided-stability} and the integral estimate:
\[
 \sup_X(v_\delta-u)_+
 \leq C\left(\int_X(v_\delta-u)_+\,\mathrm{d}V\right)^\eta
 \leq C\|R_\delta u-u\|_1^\eta\leq C\delta^\alpha.
\]
As $e\leq C\delta^\alpha$ and $u-r\leq C_0+1$, it follows that
\begin{equation}\label{eq:holder-transform-bound}
 \Phi_{c,\delta}-u
   =(1+e)(v_\delta-u)+e(u-r)\leq C\delta^\alpha.
\end{equation}

Fix $z$ and a minimizing radius $t=t(z)\in(0,\delta]$. The lower estimate at this radius and \eqref{eq:holder-transform-bound} give
\[
 -c\log(t/\delta)
 \leq\Phi_{c,\delta}(z)-u(z)+K(\delta^2+\delta)
 \leq C\delta^\alpha=Cc.
\]
Thus $t(z)\geq a\delta$ with a constant $a=\mathrm{e}^{-C}>0$ independent of $z$ and $\delta$. Monotonicity of $R_su+Ks^2$ yields
\[
 \begin{split}
 R_{a\delta}u+K(a\delta)^2
 &\leq R_tu+Kt^2\\
 &=\Phi_{c,\delta}+K\delta^2+K\delta-Kt+c\log(t/\delta)\\
 &\leq\Phi_{c,\delta}+K\delta^2+K\delta.
 \end{split}
\]
Together with \eqref{eq:holder-transform-bound} and the lower bound $R_su-u\geq-Ks^2$, this proves, after setting $s=a\delta$,
\[
 \|R_su-u\|_\infty\leq Cs^\alpha
\]
for every sufficiently small $s>0$.

For points $x,y$ at distance $d$ small enough, the derivative estimate gives
\[
 |u(x)-u(y)|\leq2Cs^\alpha+C d/s.
\]
Choose $s=d^{1/(1+\alpha)}$. Both terms are then bounded by a constant times $d^{\alpha/(1+\alpha)}$. Larger distances are controlled by the oscillation bound. This proves H\"older continuity with exponent $\alpha/(1+\alpha)$. Letting $\eta$ approach its permitted upper endpoint gives every $\beta$ in the theorem; any smaller exponent follows on compact $X$ by changing the constant.
\end{proof}

This is a regularity theorem for bounded solutions in a strictly positive Hermitian background. For a semipositive big background, division by $1+e$ alone does not remove the curvature error. \Cref{sec:local-big-continuity} uses a strict singular barrier to absorb it and proves local continuity wherever that barrier is locally bounded. It retains the general curvature error modulus and does not infer a power rate from it.


\section{Symmetry and the selected projective limit}
\label{sec:symmetric-equation-selection}

Let $\omega$ be a Hermitian form and $0\leq f\in L^p(\mathrm{d}V)$, $p>1$, a nonzero density. When the normalized equation $\mathrm{MA}_\omega(u)=cf\,\mathrm{d}V$ has several solutions, the positive-exponent equations $\mathrm{MA}_\omega(u)=\mathrm{e}^{\lambda u}f\,\mathrm{d}V$ still have unique solutions, and their normalized limits as $\lambda\downarrow0$ single out particular solutions of the normalized equation. We study this selection in the projective example of \cref{thm:degenerate-hermitian-nonuniqueness}, whose density vanishes on an open band. The starting point is a general observation: uniqueness at positive exponent preserves the symmetries of the data, and uniform compactness carries them to every normalized cluster limit. For the projective data this already gives a solution distinct from the two smooth solutions exchanged by the involution. The radial form of the equation then identifies the selected limit (\cref{thm:numa-selected-limit}) and shows that the positive-exponent solutions fail to be $C^2$ along the same hypersurfaces (\cref{prop:numa-positive-exponent-nonc2}).

\begin{proposition}\label{prop:symmetric-equation-clusters}
Let $\omega$ be a Hermitian form on $X$, let $\mathrm{d}V$ be a smooth positive volume form, and let $0\leq f\in L^p(\mathrm{d}V)$, $p>1$, be nonzero. Let $c>0$ be the unique scalar for the bounded equation $\mathrm{MA}_\omega(u)=cf\,\mathrm{d}V$ (\cref{thm:kn-continuous}). For $\lambda>0$ denote by $u_\lambda$ the unique bounded solution of
\[
 \mathrm{MA}_\omega(u_\lambda)=\mathrm{e}^{\lambda u_\lambda}f\,\mathrm{d}V,
 \qquad a_\lambda=u_\lambda-\sup_Xu_\lambda.
\]
The family $\{a_\lambda:0<\lambda\leq1\}$ is relatively compact in $C(X)$. Every cluster limit as $\lambda\downarrow0$ is a normalized bounded solution of $\mathrm{MA}_\omega(a)=cf\,\mathrm{d}V$.

Suppose a group $G$ of biholomorphisms of $X$ preserves $\omega$ and the measure $f\,\mathrm{d}V$. Then every $a_\lambda$ and every such cluster limit is $G$-invariant. In particular, the zero-exponent equation has a $G$-invariant normalized solution.
\end{proposition}
\begin{proof}
The barriers of \cref{prop:vanishing-exponent} give a uniform oscillation bound for $u_\lambda$ and uniform convergence $\mathrm{e}^{\lambda u_\lambda}\to c$. Thus the equation densities are uniformly bounded in $L^p$ for $0<\lambda\leq1$ and tend to $cf$ in $L^p$ as $\lambda\downarrow0$. Apply \cref{thm:lp-uniform-compactness} to obtain relative compactness of $a_\lambda$ in $C(X)$ and the asserted equation for every cluster limit.

For each $g\in G$, pullback of the positive-exponent equation shows that $u_\lambda\circ g$ solves the same equation. Its uniqueness gives $u_\lambda\circ g=u_\lambda$. Surjectivity of $g$ preserves the supremum, so $a_\lambda\circ g=a_\lambda$. Uniform convergence preserves these equalities for every $g$. Compactness gives at least one cluster limit as $\lambda\downarrow0$, yielding the invariant solution.
\end{proof}

\Needspace{10\baselineskip}
\begin{corollary}\label{cor:numa-third-solution}
On $\mathbb P(\mathbb C^2\oplus\mathbb C^2)$, with block homogeneous coordinates $[A:B]$, the data in \cref{thm:degenerate-hermitian-nonuniqueness} can be chosen so that the same zero-exponent equation has at least three distinct normalized bounded solutions. Two are smooth, and a third is H\"older continuous and fixed by the involution $\iota[A:B]=[B:A]$.
\end{corollary}
\begin{proof}
Use the radial coordinate $r$, rank-one forms $P,Q,R$, Fubini--Study form $\omega_X$, and scalar functions $F,x,y,z$ from \cref{subsec:numa-rank-directions,subsec:numa-scalar-coefficients}. Set $w=F(r)$, $\omega_0=xP+yQ+zR$, $H=\ddc w/2$, $A_0=\omega_0-H$ and $B_0=\omega_0+H$ as in \cref{subsec:numa-equal-cubes}. The Hermitian form is $\omega=\omega_0+\ddc\rho$, with $\rho=\varepsilon(r^2-1)\kappa(r)$ from \cref{subsec:numa-positive-background}. Choose the cutoff $\kappa$ even; all its required support and plateau properties are preserved by this choice. Then $\rho$ is even, while $w=F(r)$ is odd. The involution sends $r$ to $-r$, interchanges $P,Q$, and fixes $R$. The coefficient identities $x(-r)=y(r)$ and $z(-r)=z(r)$ show that $\iota^*\omega_0=\omega_0$. Thus $\iota^*\omega=\omega$, and $\iota^*\omega_X=\omega_X$.

Since $\iota^*H=-H$, the involution interchanges $A_0$ and $B_0$. Their cubes are equal, so $f\omega_X^3$ is invariant. It also interchanges $\widetilde u=-w/2-\rho$ and $\widetilde v=w/2-\rho$. These functions therefore have the same supremum, and the normalized solutions $u,v$ are interchanged. They are distinct, so neither is invariant.

Apply \cref{prop:symmetric-equation-clusters} to the group generated by $\iota$. Here $c=1$, as the two constructed solutions already have that scalar. An invariant cluster solution $a$ cannot equal either $u$ or $v$. Its bounded equation with smooth density gives H\"older regularity by \cref{thm:hermitian-holder-elementary}. This gives a third solution. Its regularity is determined in \cref{sec:numa-selected-limit,sec:numa-positive-exponent-regularity}: it is $C^{1,1}$ but not $C^2$.
\end{proof}

To identify the whole normalized limit for these symmetric data, we now use the radial form of the equation. It determines the invariant solution even among bounded potentials whose differentiability is not known in advance.


\Needspace{12\baselineskip}
\section{The selected limit in the symmetric projective example}
\label{sec:numa-selected-limit}

The equation of the projective example has many solutions. Nevertheless, its normalized positive-exponent approximations have a unique limit. We prove this by identifying every bounded solution with all the symmetries of the data. The argument is specific to this example.

Choose the cutoff defining $\rho$ even, and let $K$ be the group generated by the block unitary transformations $U(2)\times U(2)$ and the block exchange $\iota$. These transformations preserve $\omega$ and $f\omega_X^3$. Denote by $u_*$ the normalized even solution constructed using the symmetric transition pair in \cref{prop:numa-radial-continuum}.

\begin{theorem}\label{thm:numa-selected-limit}
For the symmetric projective data, $u_*$ is the unique normalized bounded $K$-invariant solution of
\[
 \mathrm{MA}_\omega(u)=f\omega_X^3.
\]
If $u_\lambda$ is the unique solution of
\[
 \mathrm{MA}_\omega(u_\lambda)
      =\mathrm{e}^{\lambda u_\lambda}f\omega_X^3,\qquad\lambda>0,
\]
then
\[
 u_\lambda-\sup_Xu_\lambda\longrightarrow u_*
       \quad\hbox{uniformly as }\lambda\downarrow0.
\]
The selected limit is $C^{1,1}$. This is a selection theorem for these data; it does not assert uniqueness of all normalized solutions of the zero-exponent equation.
\end{theorem}

We give the radial measure calculation needed to treat bounded solutions, rather than assume differentiability of a cluster limit. Use the notation $x,y,z,P,Q,R$ of \cref{sec:degenerate-hermitian-nonuniqueness}. For a scalar $s$ write
\[
 A(r,s)=(x(r)+s)(y(r)-s),\qquad
 \Phi(r,s)=\int_0^s A(r,t)\,\mathrm{d}t.
\]
Partial derivatives of $\Phi$ with respect to $r$ below keep $s$ fixed.

\begin{lemma}\label{lem:numa-radial-flux}
Suppose $u$ is a bounded block-unitary-invariant solution of the zero-exponent equation. On $U$ write $u+\rho=G(r)$. Then $G$ is locally Lipschitz, its derivative $h=G'$ has locally bounded variation, and its one-sided representatives satisfy
\[
 -x(r)\leq h(r\pm)\leq y(r),\qquad \mathrm{d}h+z\,\mathrm{d}r\geq0.
\]
In fact $h$ is continuous. With
\[
 k(r)=\begin{cases}0,&|r|\leq1,\\x(r)y(r)z(r),&|r|>1,\end{cases}
\]
it satisfies the distributional identity on $\mathbb R$
\begin{equation}\label{eq:numa-flux-identity}
 \mathrm{d}\bigl(\Phi(r,h(r))\bigr)
 =\bigl(k(r)-A(r,h(r))z(r)
                    +\partial_r\Phi(r,h(r))\bigr)\,\mathrm{d}r.
\end{equation}
On every open interval where $A(r,h(r))>0$, $h$ is smooth and
\[
 h'=\frac{k(r)}{A(r,h)}-z.
\]
\end{lemma}
\begin{proof}
Block-unitary orbits in $U$ are exactly the level sets of $r$: normalize the total squared norm of $(A,B)$, and map the two vectors of the same respective lengths by their unitary groups. Thus invariance gives a scalar $G$.

On a compact $r$-interval, the radial direction annihilates $P,Q$, while $R$ is strictly positive in that direction. Testing positivity of $\omega_0+\ddc G(r)$ in this direction gives $G''+z\geq0$ as a one-dimensional distribution. One may obtain this first by smooth approximation of the scalar distribution and then by testing with nonnegative functions. Adding a twice integrated primitive of $z$ makes $G$ convex. A finite bounded convex function on an interval is locally Lipschitz, and its derivative is locally of bounded variation. The quasi-plurisubharmonic representative is the corresponding continuous representative. The horizontal directions give $x+h\geq0$, $y-h\geq0$ almost everywhere. Continuity of $x,y$ and the one-sided limits of a function of bounded variation give these bounds for both one-sided representatives. The radial inequality is precisely $\mathrm{d}h+z\,\mathrm{d}r\geq0$; in particular jumps of $h$ can only be upward.

We next justify the formula for the measure, including possible jumps. Identify $U$ smoothly with $\mathbb R\times Y$, where $Y=\{r=0\}$, by real scaling of one block. The forms $P,Q$ are independent of $r$ in this identification. So is $\dc r$: in local coordinates it is the angular differential of the relative complex scale plus fixed base terms. Therefore
\[
 P\wedge Q\wedge R=\mathrm{d}r\wedge\nu,
\]
where $\nu$ is a fixed smooth positive volume form on the compact orbit $Y$, with the orientation chosen by this equality.

For smooth $G$, the rank-one wedge formula gives the radial coefficient of its Monge--Amp\`ere measure, after dividing by $6\nu$, as
\[
 A(r,h)(z+h')\,\mathrm{d}r
 =A(r,h)z\,\mathrm{d}r+\mathrm{d}\bigl(\Phi(r,h)\bigr)
                         -\partial_r\Phi(r,h)\,\mathrm{d}r.
\]
This identity passes to a locally Lipschitz $G$ with the preceding lower Hessian bound. Indeed, convolve $G$ on a slightly larger compact interval. The derivatives $h_\epsilon$ are uniformly bounded and converge to $h$ almost everywhere and in every finite $L^q$ space. Their second derivatives have a uniform lower bound. Consequently the functions $G_\epsilon(r)$ have a common quasi-plurisubharmonic lower Hessian bound on compact subsets of $U$, and converge uniformly to $G(r)$. A fixed smooth local correction and bounded Bedford--Taylor continuity give convergence of their signed products. On the displayed right side, the two ordinary terms converge in $L^1$, and $\Phi(r,h_\epsilon)\to\Phi(r,h)$ in $L^1$, giving distributional convergence of its derivative. This proves the same formula for the bounded product of $G$.

The specified smooth data measure has radial coefficient $k(r)\,\mathrm{d}r$, giving \eqref{eq:numa-flux-identity} initially for the almost-everywhere representative of $h$. If $h$ jumped at $r_0$, the atom of its left side would be
\[
 \Phi(r_0,h(r_0+))-\Phi(r_0,h(r_0-))
 =\int_{h(r_0-)}^{h(r_0+)}(x(r_0)+s)(y(r_0)-s)\,\mathrm{d}s.
\]
The right side has no atom. An upward jump of positive size lies in $[-x(r_0),y(r_0)]$, where the integrand is strictly positive in the interior. Since $x+y>0$ everywhere, such an integral is strictly positive. Thus there are no jumps. A locally bounded-variation function whose jumps all vanish has a continuous representative here: its distributional derivative consists of an absolutely continuous measure and a nonnegative singular measure, so its one-sided representatives coincide at every point. We henceforth use this representative. In particular $G$ is $C^1$ and $G'=h$ everywhere.

All coefficients on the right of \eqref{eq:numa-flux-identity} are now continuous, so $\Phi(r,h(r))$ is $C^1$. Where $\partial_s\Phi=A>0$, the implicit function theorem gives $h\in C^1$ and the stated ordinary differential equation. Smoothness then follows by the smooth coefficients of that equation. The function $k$ is smooth, since the original density and its flat joins are smooth.
\end{proof}

\begin{proof}[of \cref{thm:numa-selected-limit}]
Let $u$ be a normalized bounded $K$-invariant solution. Apply the radial lemma. Because $\rho$ and $u$ are invariant under block exchange, $G$ is even, so $h$ is odd and $h(0)=0$.

On $0<r<1$ we have $x=p>0$ and $y=p+b\geq p>0$. Starting from zero, as long as $-p<h<y$, the equation is $h'=-z_0$. Thus $h(r)=-\int_0^r z_0(s)\,\mathrm{d}s$. This decreasing function cannot meet the positive upper boundary $y$. Its first meeting with the lower boundary $-p$ occurs at the unique $t_*\in(0,1)$ satisfying
\[
 \int_0^{t_*}z_0(s)\,\mathrm{d}s=p(t_*),
\]
whose existence and uniqueness were proved in \cref{prop:numa-radial-continuum}.

After that meeting, $h$ must stay on $-p$ until $r=1$. To prove it, suppose a connected component $(a,b)$ of $\{h+p>0\}$ lies to its right. Its left endpoint has $a\geq t_*$ and $h(a)=-p(a)<0$. Near that endpoint, $h<0<y$; hence on the component the interior equation applies for as long as $h<0$. It forces $h'=-z_0<0$, so $h$ cannot first reach zero. The equation therefore holds throughout the component. But then
\[
 (h+p)'=p'-z_0<0,
\]
contradicting positivity of $h+p$ immediately to the right of its zero at $a$. Thus $h=-p$ on $[t_*,1]$. Oddness determines the entire negative half of the band. The slope on $[-1,1]$ is exactly that of the symmetric explicit solution.

We also identify the slope outside the band; boundedness at the projective ends matters here. On $r>1$, $k=xyz>0$. Neither horizontal coefficient can vanish at a finite point of this interval. For the lower one define
\[
 \Psi(r)=\int_{-x(r)}^{h(r)}A(r,s)\,\mathrm{d}s\geq0.
\]
Equation \eqref{eq:numa-flux-identity} and the smooth lower endpoint show that $\Psi$ is $C^1$. If $h(r_0)=-x(r_0)$, its derivative at that point is $k(r_0)>0$: the $Az$ term vanishes, and the two partial-$r$ terms agree at the common endpoint. This contradicts that the nonnegative $C^1$ function $\Psi$ has an interior minimum zero there. The upper-endpoint integral $\int_{h(r)}^{y(r)}A(r,s)\,\mathrm{d}s$ gives the same contradiction, with derivative $-k(r_0)$, when $h=y$. Thus $-x<h<y$ everywhere on $r>1$, and the smooth ordinary differential equation of the lemma applies.

For $r\geq2$ the background is Fubini--Study:
\[
 x=\sigma,\qquad y=1-\sigma,\qquad z=\sigma',
 \qquad \sigma=\frac{\mathrm{e}^r}{1+\mathrm{e}^r}.
\]
Set $a=\sigma+h\in(0,1)$. Its equation is
\[
 a(1-a)a'=\sigma(1-\sigma)\sigma'>0.
\]
Hence $a$ is increasing and has a limit in $(0,1]$ at $+\infty$. If that limit were less than one, $h=a-\sigma$ would approach a negative constant and its primitive $G$ would be unbounded below, contrary to boundedness of $u$ and $\rho$. The limit is therefore one. Integrating the equation gives
\[
 \frac{a^2}{2}-\frac{a^3}{3}
  =\frac{\sigma^2}{2}-\frac{\sigma^3}{3}+C.
\]
Taking the limit at infinity gives $C=0$. The scalar function $s^2/2-s^3/3$ is strictly increasing on $(0,1)$, so $a=\sigma$ and $h=0$ for $r\geq2$. Ordinary uniqueness propagates this zero solution throughout $r>1$, since the coefficients are smooth and the two horizontal coefficients stay positive there. Oddness gives $h=0$ on $r<-1$ as well.

We have proved that $G'$ agrees everywhere with the symmetric explicit slope. Thus $u$ differs from $u_*$ by a constant on $U$. The two quasi-plurisubharmonic functions then agree up to that constant on $X$ by their canonical representatives. Their supremum normalization removes the constant. This proves uniqueness within the bounded $K$-invariant class.

Finally \cref{prop:symmetric-equation-clusters} gives uniform precompactness of the normalized positive-exponent family. Every cluster limit solves the zero-exponent equation and is $K$-invariant, because the data are $K$-invariant. The scalar is one for these data. The uniqueness just proved identifies every cluster limit with $u_*$. If convergence of the whole family failed, a sequence staying a fixed positive uniform distance from $u_*$ would have a convergent subsequence, contradicting that identification. This proves uniform convergence and the asserted $C^{1,1}$ regularity of the limit.
\end{proof}


\section{Loss of second derivatives at positive exponent}
\label{sec:numa-positive-exponent-regularity}

Let $X=\mathbb P(\mathbb C^2\oplus\mathbb C^2)$, with block coordinates $[A:B]$, and let $\omega_X$ be the Fubini--Study form. On $U=\{A\ne0,\ B\ne0\}$ use the radial coordinate $r$ and rank-one forms $P,Q,R$ of \cref{subsec:numa-rank-directions}. Take the scalar coefficients $p,q,b,z_0,x,y,z$ from \cref{subsec:numa-scalar-coefficients}. Let $\omega_0=xP+yQ+zR$, and let $\rho$ and $\omega=\omega_0+\ddc\rho$ be the correction and Hermitian form from \cref{subsec:numa-positive-background}. The density $f$ is the common cube there, so $f\omega_X^3=0$ on $|r|\leq1$ and $f\omega_X^3=\omega_0^3$ outside that band. Write $\mathrm{MA}_\omega(u)=(\omega+\ddc u)^3$ for bounded $\omega$-plurisubharmonic $u$. Choose the correction cutoff even and let $K$ be the group generated by block unitary transformations and block exchange. Let $t_*\in(0,1)$ be the unique solution of $\int_0^{t_*}z_0(s)\,\mathrm{d}s=p(t_*)$, and let $u_*$ be the normalized symmetric solution of \cref{prop:numa-radial-continuum}. In the flux formula below use $A(r,s)=(x(r)+s)(y(r)-s)$ and $k(r)=0$ on $|r|\leq1$, $k(r)=x(r)y(r)z(r)$ on $|r|>1$, as in \cref{lem:numa-radial-flux}.

The selected solution is not $C^2$. In fact the same local loss of regularity occurs at every positive exponent, despite smoothness of the background and the nonnegative density. Thus the exponential term restores uniqueness without removing this degeneracy.

\Needspace{12\baselineskip}
\begin{proposition}\label{prop:numa-positive-exponent-nonc2}
For the symmetric projective data of \cref{thm:numa-selected-limit}, let $u_\lambda$ solve
\[
 \mathrm{MA}_\omega(u_\lambda)
       =\mathrm{e}^{\lambda u_\lambda}f\omega_X^3,
 \qquad\lambda>0.
\]
On the open band $\{-1<r<1\}$, $u_\lambda-u_*$ is constant. Both potentials are $C^{1,1}$ there and fail to be $C^2$ at every point of the two hypersurfaces $r=\pm t_*$. In particular, none of these positive-exponent solutions is a globally $C^2$ potential.
\end{proposition}
\begin{proof}
Uniqueness and invariance of the data imply that $u_\lambda$ is $K$-invariant. Write $u_\lambda+\rho=G_\lambda(r)$ on $U$. The positivity argument in \cref{lem:numa-radial-flux} uses only boundedness and invariance, so $G_\lambda$ is locally Lipschitz and $h_\lambda=G_\lambda'$ is locally of bounded variation, with
\[
 -x\leq h_\lambda(r\pm)\leq y,
 \qquad \mathrm{d}h_\lambda+z\,\mathrm{d}r\geq0.
\]
Its weak flux equation is the same one, with the prescribed coefficient replaced by
\[
 k_\lambda(r)=\mathrm{e}^{\lambda(G_\lambda(r)-\rho(r))}k(r).
\]
Indeed, the scalar convolution and signed-product argument of that lemma is unchanged, and this is exactly the radial coefficient of the right-hand measure. At the initial bounded stage $k_\lambda$ is locally bounded. Hence the right side of the flux equation has no atoms. The same strictly positive jump integral eliminates every jump of $h_\lambda$. Thus $h_\lambda$ is continuous and $G_\lambda$ is $C^1$. The coefficient $k_\lambda$ is consequently continuous. Where $A(r,h_\lambda)>0$, inversion of the $C^1$ flux gives
\[
 h_\lambda'=\frac{k_\lambda(r)}{A(r,h_\lambda)}-z.
\]

On $-1<r<1$ the coefficient $k_\lambda$ is identically zero. Exchange invariance gives $h_\lambda(0)=0$ and oddness. On the positive half-band, the first-contact and no-departure arguments in the proof of \cref{thm:numa-selected-limit} apply verbatim: the first branch solves $h'=-z_0$, reaches $-p$ at $t_*$, and cannot leave that lower boundary to its right because $(h+p)'=p'-z_0<0$ on every proposed departure interval. Therefore
\[
 h_\lambda(r)=
 \begin{cases}
  p(r),&-1<r\leq-t_*,\\
  -\displaystyle\int_0^r z_0(s)\,\mathrm{d}s,&-t_*\leq r\leq t_*,\\
  -p(r),&t_*\leq r<1.
 \end{cases}
\]
This is the derivative of $u_*+\rho$. Its primitive differs from $G_\lambda$ by one constant on the connected band. The claimed local $C^{1,1}$ regularity follows from the explicit slope.

At $t_*\in(0,1)$ the two one-sided derivatives of this slope are
\[
 h_\lambda'(t_*-)=-z_0(t_*)<0,
 \qquad
 h_\lambda'(t_*+)=-p'(t_*)>0.
\]
These unequal derivatives prevent the scalar primitive from having a second derivative at $t_*$. Through every point of that level set, the submersion $r$ supplies a smooth curve parametrized by $r$, transferring the failure to $u_\lambda$. Smooth $\rho$ cannot cancel the jump. Oddness gives the conclusion at $-t_*$, and the identical slope proves it for $u_*$.
\end{proof}

The zero set of $f$ contains the whole band. This example therefore does not concern the smooth strictly positive densities in the Hermitian Calabi--Yau theorem. It does show why smoothness of a nonnegative density alone is insufficient for second-order regularity, even when the equation includes a strictly increasing exponential factor.


\section{Local continuity in a semipositive big background}
\label{sec:local-big-continuity}

A strict singular barrier repairs the curvature error of regularization even when the background itself is degenerate. The one-sided stability theorem allows the resulting comparison function to be unbounded. This gives local continuity of bounded solutions wherever the barrier is locally bounded. We retain the general error modulus in the classical regularization theorem; a power rate for that modulus is not needed.

Let $X$ be a connected compact complex manifold of dimension $n\geq1$. Fix a Hermitian form $\omega$ for regularization and a smooth positive probability measure $\mathrm{d}V$. Let $\theta$ be a smooth semipositive big real $(1,1)$-form on $X$. Choose $\rho\in\PSH(X,\theta)$ and $\kappa>0$ such that
\[
 \theta+\ddc\rho\geq\kappa\omega.
\]
Let $U$ be any open set on which $\rho$ is locally bounded. Such a set need not be chosen canonically for the following argument.

Let $X$ be a connected compact complex manifold of dimension $n\geq1$. Fix a Hermitian form $\omega$ for regularization and a smooth positive probability measure $\mathrm{d}V$. Let $\theta$ be a smooth semipositive big real $(1,1)$-form on $X$. Choose $\rho\in\PSH(X,\theta)$ and $\kappa>0$ with $\theta+\ddc\rho\geq\kappa\omega$, and let $U\subset X$ be an open set on which $\rho$ is locally bounded. Write $\mathrm{MA}_\theta(u)=(\theta+\ddc u)^n$ for bounded potentials and $s_+=\max(s,0)$ for a real number $s$. For the potential $u$ in the theorem below, use the averages $R_tu$, truncated transforms $\Phi_{c,\delta}$, constants $A,K$ and curvature modulus $\epsilon$ from \cref{subsec:general-curvature-modulus,lem:general-kiselman-modulus}.

\begin{theorem}\label{thm:local-big-bounded-continuity}
Let $\mathrm{d}V$ be a smooth positive probability measure, $p>1$, and let $u\in\PSH(X,\theta)\cap L^\infty(X)$ satisfy
\[
 \mathrm{MA}_\theta(u)=f\,\mathrm{d}V,\qquad 0\leq f\in L^p(\mathrm{d}V).
\]
Then $u$ is continuous on every open set where the fixed strict barrier $\rho$ is locally bounded. No bounded mass or positive lower-volume property is assumed.
\end{theorem}
\begin{proof}
Normalize $u$ so that $-C_0\leq u\leq0$, and subtract a constant from $\rho$ so that $\rho\leq-C_0-1$. Choose
\[
 0<\eta<\frac{p-1}{np+p-1},\qquad \alpha=2\eta<1,
\]
and set, for small $\delta>0$,
\[
 c=\delta^\alpha+\epsilon(\delta),\qquad
 e=Ac+\epsilon(\delta),\qquad s=e/\kappa,\qquad
 v_\delta=\frac{\Phi_{c,\delta}+s\rho}{1+s}.
\]
The curvature bounds give $v_\delta\in\PSH(X,\theta)$, since $s\kappa=e$. Moreover
\[
 \rho\leq-C_0-1\leq\Phi_{c,\delta}
                  \leq R_\delta u\leq0
\]
for small $\delta$. The lower bound follows from
$\Phi_{c,\delta}\geq u-K(\delta^2+\delta)$ as in the strictly positive case. Hence $v_\delta\leq\Phi_{c,\delta}\leq R_\delta u\leq0$.

Apply \cref{thm:bounded-one-sided-stability}, whose comparison potential may be unbounded. The preceding smoothing estimate gives
\[
 \sup_X(v_\delta-u)_+
 \leq C\left(\int_X(v_\delta-u)_+\,\mathrm{d}V\right)^\eta
 \leq C\delta^\alpha.
\]
Fix $M$ and restrict to points where $\rho\geq-M$. As $s\leq Cc$ and $c\to0$, the exact identity
\[
 \Phi_{c,\delta}-u=(1+s)(v_\delta-u)+s(u-\rho)
\]
gives $\Phi_{c,\delta}-u\leq C_Mc$ there.

At a minimizing radius $t$, the same lower estimate in its defining expression yields
\[
 -c\log(t/\delta)
       \leq\Phi_{c,\delta}-u+K(\delta^2+\delta)
       \leq C_Mc.
\]
Thus $t\geq a_M\delta$ with $a_M>0$ independent of the point and $\delta$. Monotonicity of smoothing now gives
\[
 R_{a_M\delta}u+K(a_M\delta)^2
       \leq R_tu+Kt^2
       \leq\Phi_{c,\delta}+K\delta^2+K\delta.
\]
\Needspace{5\baselineskip}
Its upper bound and the monotone lower bound imply
\[
 \sup_{\{\rho\geq-M\}}
       |R_{a_M\delta}u-u|
       \leq C_M\bigl(\delta^\alpha+\epsilon(\delta)\bigr)
                       \longrightarrow0.
\]
Every compact subset of $U$ has a lower bound for $\rho$. The smooth functions on the left therefore converge uniformly to $u$ on that compact subset. This proves continuity on $U$.
\end{proof}

The proof also gives a modulus on each fixed compact subset: combine its uniform smoothing error with the derivative bound $C/t$. A H\"older exponent would require a proved power bound on the curvature error modulus. For a smooth big background, \cref{lem:smooth-background-regularization-rate} supplies a linear curvature error by refining the classical expansion. Together with stability above a strict barrier, it gives local H\"older continuity even for minimal-singularity solutions in \cref{thm:big-minimal-holder}. Neither argument asserts continuity across the pole set.


\section{Homogeneous replacement on a coordinate ball}
\label{sec:local-homogeneous-dirichlet}

To show that the Monge--Amp\`ere measure of an envelope vanishes off the contact set, we shall replace the envelope on small balls by a solution of the homogeneous equation with the same boundary values. We first recall the smooth Dirichlet problem with a smooth positive density; the replacement is obtained by letting a constant density tend to zero, and continuous boundary values are then handled by uniform approximation.

In this section we work on the ball $B=B_R=\{|z|<R\}\subset\mathbb C^n$, put $\beta=\ddc|z|^2$, and let $\theta$ be any smooth real $(1,1)$-form defined near $\overline B$. Distances and the Laplacian $\Delta$ are Euclidean.

We use smooth comparison on subdomains $D$ of the ball. Let $a,b$ be smooth functions on $\overline D$ with $\theta+\ddc a>0$ and $\theta+\ddc b>0$. If $(\theta+\ddc a)^n\geq(\theta+\ddc b)^n$ in $D$ and $a\leq b$ on $\partial D$, then $a\leq b$ in $D$. Indeed, integrating the derivative of the logarithm of the determinant along the segment between the two positive matrices gives a linear elliptic inequality for $a-b$, and the maximum principle proves the assertion. The boundary of $D$ need not be smooth for this argument.

\begin{lemma}\label{lem:ball-smooth-dirichlet}
Let $\varphi\in C^\infty(\partial B)$, and let $F$ be a smooth strictly positive function on $\overline B$. There is a unique smooth function $v$ on $\overline B$ such that
\[
 \theta+\ddc v>0,\qquad
 (\theta+\ddc v)^n=F\beta^n,
 \qquad v|_{\partial B}=\varphi.
\]
\end{lemma}
\begin{proof}
Extend $\varphi$ to a smooth function $\widetilde\varphi$ on $\overline B$. For sufficiently large $A$, the function $\underline v=\widetilde\varphi+A(|z|^2-R^2)$ is a smooth strict subsolution with the prescribed boundary values: $\theta+\ddc\underline v$ is positive and its determinant is at least $(\max_{\overline B}F)\beta^n$. For $n\geq2$, the classical smooth Dirichlet theorem of Guan--Li therefore applies on the ball, with the arbitrary smooth background $\chi=\theta$ and the smooth positive right-hand side $F$. Their operator is $(\mathrm{i}/2)\partial\bar\partial$, so our potential is $\pi$ times theirs. No closedness condition is needed here; such conditions enter only in their separate theorems on compact manifolds. For $n=1$ the lemma is the Dirichlet problem for the Poisson equation. Uniqueness follows from smooth comparison, applied in both directions.
\end{proof}

\Needspace{15\baselineskip}
\begin{theorem}\label{thm:ball-homogeneous-replacement}\index{Dirichlet problem!homogeneous equation}
For every $\varphi\in C(\partial B)$ there exists a continuous $\theta$-plurisubharmonic function $H_\varphi$ on $\overline B$ such that
\[
 (\theta+\ddc H_\varphi)^n=0\quad\hbox{on }B,
 \qquad H_\varphi|_{\partial B}=\varphi.
\]
It dominates every continuous $\theta$-plurisubharmonic function $p$ on $\overline B$ with $p|_{\partial B}\leq\varphi$. For $\varphi,\psi\in C(\partial B)$, the construction satisfies
\[
 \|H_\varphi-H_\psi\|_{C(\overline B)}
       \leq\|\varphi-\psi\|_{C(\partial B)}.
\]
For smooth boundary data the constructed function is Lipschitz on $\overline B$.
\end{theorem}
\begin{proof}
\emph{Step 1: boundary barriers.} We first take $\varphi$ smooth and $0<\varepsilon\leq1$, and let $v_\varepsilon$ be the solution of \cref{lem:ball-smooth-dirichlet} with the constant density $F\equiv\varepsilon$. The strict subsolution $\underline v$ constructed there for $F\equiv1$ lies below every $v_\varepsilon$, by smooth comparison. Choose a smooth function $q$ with $q=\varphi$ on the boundary and $\Delta q=-C$, where $C$ is so large that $\theta+\ddc v_\varepsilon\geq0$ implies $\Delta v_\varepsilon\geq-C$. The maximum principle for the Laplacian gives $v_\varepsilon\leq q$. Both barriers are smooth and equal $\widetilde\varphi$ on the boundary. Consequently
\begin{equation}\label{eq:ball-boundary-distance}
 |v_\varepsilon(z)-\widetilde\varphi(z)|
       \leq C_1\operatorname{dist}(z,\partial B),
\end{equation}
with $C_1$ independent of $\varepsilon$.

\emph{Step 2: a Lipschitz bound independent of $\varepsilon$.} Let $h\in\mathbb C^n$ be small, and set $D_h=\{z\in B:z+h\in B\}$. Since $\theta$ is smooth, there is $C_2>0$, independent of the translation $h$, such that
\[
 \theta(z)-\theta(z+h)+C_2|h|\beta\geq0.
\]
For a constant $C_3>0$ to be fixed, define on $D_h$
\[
 w(z)=v_\varepsilon(z+h)
          +C_2|h|(|z|^2-R^2)-C_3|h|.
\]
Then $\theta+\ddc w\geq\theta(z+h)+\ddc v_\varepsilon(z+h)>0$, and the determinant of $\theta+\ddc w$ is at least $\varepsilon\beta^n$, because the Euclidean volume form is translation invariant. On $\partial D_h$, one of the points $z$ and $z+h$ lies on $\partial B$, and the other has distance at most $|h|$ from $\partial B$. Estimate \eqref{eq:ball-boundary-distance} and the Lipschitz bound for $\widetilde\varphi$ therefore give $w\leq v_\varepsilon$ on $\partial D_h$ if $C_3$ is chosen sufficiently large, and smooth comparison extends this inequality to all of $D_h$. Together with the same argument for $-h$, this gives
\[
 |v_\varepsilon(z+h)-v_\varepsilon(z)|\leq C_4|h|,
\]
with constants independent of $\varepsilon$. The same bound for arbitrary pairs of points follows by subdividing the segment joining them inside the convex ball, and it extends to $\overline B$ because each $v_\varepsilon$ is continuous there.

\emph{Step 3: the limit $\varepsilon\to0$.} Smooth comparison also shows that $v_\varepsilon$ increases as $\varepsilon$ decreases. By the barriers and the common Lipschitz bound, the limit $H_\varphi$ is Lipschitz and the convergence is uniform on $\overline B$. The inequality $\theta+\ddc v_\varepsilon>0$ of currents passes to the limit, and continuity of local bounded products under uniform convergence gives $(\theta+\ddc H_\varphi)^n=0$. For smooth boundary values $\varphi,\psi$, denote the solutions of \cref{lem:ball-smooth-dirichlet} with $F\equiv\varepsilon$ by $v_{\varepsilon,\varphi}$ and $v_{\varepsilon,\psi}$. Smooth comparison gives $\|v_{\varepsilon,\varphi}-v_{\varepsilon,\psi}\|_\infty\leq\|\varphi-\psi\|_\infty$, and passing to the limit proves the contraction for smooth data.

\emph{Step 4: continuous boundary values.} For continuous $\varphi$, choose smooth $\varphi_j\to\varphi$ uniformly on $\partial B$. By the contraction, $H_{\varphi_j}$ is uniformly Cauchy. Its limit $H_\varphi$ is continuous and $\theta$-plurisubharmonic, has the prescribed boundary values, and satisfies the homogeneous equation, again by continuity of bounded products under uniform convergence. Passing to the limit also proves the contraction for all continuous boundary data.

\emph{Step 5: domination.} It remains to prove the domination property, which does not follow from a comparison of total masses. First take $\varphi$ smooth and $p$ as in the statement. For $\delta>0$ put $p_\delta=p+\delta(|z|^2-R^2)$, so that $\theta+\ddc p_\delta\geq\delta\beta$. By \cref{lem:qpsh-upper-test}, applied with $\theta-\delta\beta$ in place of $\theta$, every $C^2$ function $a$ touching $p_\delta$ from above at an interior point $z_0$ satisfies
\[
 \theta(z_0)+\ddc a(z_0)\geq\delta\beta.
\]
Suppose that $p_\delta-v_\varepsilon$ had a positive maximum. Since this difference is nonpositive on the boundary, the maximum would be attained at an interior point, where $v_\varepsilon$ plus the maximum touches $p_\delta$ from above. The determinant of $\theta+\ddc v_\varepsilon$ at that point would then be at least $\delta^n\beta^n$, contradicting the equation for $v_\varepsilon$ when $\varepsilon<\delta^n$. Hence $p_\delta\leq v_\varepsilon$ for small $\varepsilon$. Letting first $\varepsilon$ and then $\delta$ decrease to zero, we obtain $p\leq H_\varphi$.

Finally, for continuous $\varphi$, take the smooth approximants $\varphi_j$ above and set $a_j=\|\varphi_j-\varphi\|_\infty$. The boundary values of $p-a_j$ are at most $\varphi_j$, so the result for smooth data gives $p-a_j\leq H_{\varphi_j}$. Uniform convergence proves $p\leq H_\varphi$.
\end{proof}


\Needspace{10\baselineskip}
\section{Contact support in a big background}
\label{sec:big-continuous-contact}

Let $\theta$ now be an arbitrary smooth big background. Regularization makes the envelope of a continuous obstacle continuous off the poles of a strict analytic barrier. A homogeneous replacement below the obstacle then forces its Monge--Amp\`ere measure to vanish away from contact.

Choose $\rho$ with analytic singularities and smooth remainder so that
\[
 \theta+\ddc\rho\geq\kappa\omega,\qquad\kappa>0,
\]
and set $Z=\{\rho=-\infty\}$. Such a barrier exists by \cref{thm:resolvable-regularization-input}.

\Needspace{12\baselineskip}
\begin{theorem}\label{thm:big-continuous-envelope-contact}
For $h\in C(X)$, the envelope $p=P_\theta(h)$ is genuine and continuous on $X\setminus Z$. Moreover its extended Monge--Amp\`ere measure satisfies
\[
 \mathbf1_{\{p<h\}}\mathrm{MA}_\theta(p)=0.
\]
The measure in this formula is allowed to have infinite total mass.
\end{theorem}

\begin{proof}
Subtract a constant from $\rho$ so that $\rho+2\leq h$. The elementary envelope construction and upper semicontinuity of the obstacle give
\[
 \rho+2\leq p\leq h.
\]
In particular $p$ is locally bounded off $Z$. We first prove continuity there, without using the contact assertion.

Apply \cref{lem:smooth-background-regularization-rate} to $p-\sup_Xp$ and restore the constant. The transforms $\Phi_{c,\delta}$ from that lemma are defined on $X$ and, with fixed $A,K,C>0$ and $\epsilon(\delta)=C\delta$, satisfy
\begin{equation}\label{eq:big-envelope-regularization}
 \theta+\ddc\Phi_{c,\delta}\geq-(Ac+\epsilon(\delta))\omega.
\end{equation}
Their infimum is attained at a positive radius off $Z$, and
\[
 p-K(\delta^2+\delta)\leq\Phi_{c,\delta}
       \leq R_\delta p\leq R_\delta h.
\]
For small $\delta$ the first bound gives $\Phi_{c,\delta}\geq\rho+1$, hence $\rho\leq\Phi_{c,\delta}$ on $X$.

\Needspace{5\baselineskip}
Set
\[
 m(\delta)=\sup_{0<t\leq\delta}\|R_th-h\|_\infty,
 \qquad c=m(\delta)+\epsilon(\delta)+\delta,
 \qquad s=\frac{Ac+\epsilon(\delta)}{\kappa}.
\]
Positivity, normalization, and localization of the averaging kernels imply $m(\delta)\to0$ by uniform continuity of $h$. Define
\[
 v_\delta=\frac{\Phi_{c,\delta}+s\rho}{1+s}.
\]
The strict barrier and \eqref{eq:big-envelope-regularization} show that $v_\delta$ is a genuine $\theta$-plurisubharmonic function. Since $\rho\leq\Phi_{c,\delta}$,
\[
 v_\delta\leq\Phi_{c,\delta}\leq R_\delta h\leq h+m(\delta).
\]
Thus $v_\delta-m(\delta)$ is an envelope competitor, and $v_\delta\leq p+m(\delta)$.

On a fixed compact subset $L\Subset X\setminus Z$, the difference $p-\rho$ is bounded above. As $s\leq Cc$, the exact identity
\[
 \Phi_{c,\delta}-p=(1+s)(v_\delta-p)+s(p-\rho)
\]
therefore gives $\Phi_{c,\delta}-p\leq C_Lc$. At a minimizing radius $t$,
\[
 -c\log(t/\delta)
        \leq\Phi_{c,\delta}-p+K(\delta^2+\delta)\leq C_Lc.
\]
Since $c\geq\delta$, this yields $t\geq a_L\delta$ with $a_L>0$ fixed. Monotonicity of smoothing now implies
\[
 p-K(a_L\delta)^2\leq R_{a_L\delta}p
       \leq\Phi_{c,\delta}+K\delta^2+K\delta.
\]
Consequently $R_{a_L\delta}p\to p$ uniformly on $L$. These averages are smooth for positive radius even though $p$ is globally unbounded. Applying this to compact neighborhoods proves continuity on $X\setminus Z$.

Now take $x\in\{p<h\}\setminus Z$. By the continuity just proved this is an ordinary open set near $x$. Choose a sufficiently small coordinate ball $B_r$ centered at $x$, with closure in that set. Let $H$ be the homogeneous replacement of \cref{thm:ball-homogeneous-replacement} with boundary data $p|_{\partial B_r}$. Its domination property gives $H\geq p$ on the ball.

We can make the ball small enough that $H\leq h$. Indeed, with $z$ denoting the chosen Euclidean coordinates and $\Delta$ their scalar Laplacian, $\theta+\ddc H\geq0$ implies a fixed bound $\Delta H\geq-C$. Comparison with the function
$\sup_{\partial B_r}p+C'(r^2-|z-x|^2)$, with $C'$ fixed large enough, gives
\[
 H\leq\sup_{\partial B_r}p+C'r^2.
\]
As $r\to0$, the right-hand side tends to $p(x)$, whereas $\inf_{B_r}h\to h(x)>p(x)$. This proves the claim for sufficiently small $r$.

Paste $H$ into $p$ on the ball. The result is $\theta$-plurisubharmonic: for each $a>0$, paste $\max\{p,H-a\}$ inside the ball and $p$ outside. Continuity and equality of the boundary values ensure that $H-a<p$ on a collar, so this is an ordinary local maximum construction. Letting $a\downarrow0$ gives the desired pasted function, with convergence locally uniform near the boundary. It is bounded above by $h$ on all of $X$. The defining property of the envelope therefore forces it to equal $p$. Thus $p=H$ on $B_r$ and $\mathrm{MA}_\theta(p)$ vanishes there.

These balls cover $\{p<h\}\setminus Z$. The local bounded product there agrees with the extended measure, and the latter does not charge the analytic, hence pluripolar, set $Z$. This proves the contact assertion on $X$, without taking or bounding its total mass.
\end{proof}


\section{Quasi-continuous obstacles in a big background}
\label{sec:big-general-contact}

We now remove continuity and boundedness of the obstacle in \cref{thm:big-continuous-envelope-contact}. Two points require care. An envelope competitor lies below its obstacle only quasi-everywhere, so an increasing approximation of the obstacle needs a correction at the exceptional set. Also, truncation by a constant need not preserve positivity in the original background. We use signed bounded products before restricting to a truncation set.

Fix a smooth big real $(1,1)$-form $\theta$ and a strict potential $\rho$ with analytic singularities, as in the preceding section:
\[
 \theta+\ddc\rho\geq\kappa\omega,\qquad \kappa>0.
\]
Its pole set is denoted by $Z$. In particular $\rho$ is smooth on $X\setminus Z$. Capacity below may be computed using a fixed finite cover by coordinate balls. Its null sets are precisely the pluripolar sets, and its quasi-continuity notion agrees with the one already used for envelopes.

\begin{lemma}
\label{lem:big-pointwise-competitor}
Let $f$ be bounded, and let $u\in\PSH(X,\theta)$ satisfy $u\leq f$ quasi-everywhere. There are genuine $u_t\in\PSH(X,\theta)$, $t>0$, such that $u_t\leq f$ everywhere and $u_t\to u$ almost everywhere as $t\downarrow0$.
\end{lemma}
\begin{proof}
The exceptional set $E=\{u>f\}$ is pluripolar. The global form of Josefson's theorem on compact complex manifolds supplies a genuine quasi-plurisubharmonic $\chi$ with $\chi=-\infty$ on $E$; see \cite[Theorem 1.1.5]{SGPT}. Rescale and normalize it so that $\chi\leq0$ and $\ddc\chi\geq-\omega$. Shift $\rho$ so that $\rho\leq\inf_X f$. With $s=t/\kappa$, put
\[
 u_t=\frac{u+t\chi+s\rho}{1+s}.
\]
All three potentials are finite almost everywhere, so this is genuine. Moreover
\[
 \theta+\ddc u_t
 =\frac{\theta+\ddc u+t\ddc\chi+s(\theta+\ddc\rho)}{1+s}
 \geq0.
\]
Outside $E$, each of $u$ and $\rho$ is at most $f$, while $\chi\leq0$, so $u_t\leq f$. On $E$ the numerator is $-\infty$. Finally $u_t\to u$ wherever all three potentials are finite, which is a set of full measure.
\end{proof}

\Needspace{10\baselineskip}
\begin{proposition}\label{prop:big-bounded-lsc-contact}
If $f$ is bounded and lower semicontinuous, then $p=P_\theta(f)$ is locally bounded on $X\setminus Z$, and
\[
 \mathbf1_{\{p<f\}}\mathrm{MA}_\theta(p)=0.
\]
\end{proposition}
\begin{proof}
Take continuous $f_j\uparrow f$ with common bounds $m\leq f_j\leq M$, using the infimum construction in \cref{prop:lsc-obstacle-support}. Shift $\rho$ below $m$. The envelopes $p_j=P_\theta(f_j)$ increase and satisfy
\[
 \rho\leq p_j\leq M.
\]
Their regularized limit $q=(\lim_jp_j)^*$ is genuine. Outside a pluripolar set it equals the pointwise limit and is at most $f$, so $q\leq p=P_\theta(f)$.

We prove the reverse inequality without a domination theorem. Let $u$ be any envelope competitor. Apply \cref{lem:big-pointwise-competitor} to obtain $u_t\leq f$ everywhere. For fixed $t,\eta>0$, the sets
\[
 \{u_t<f_j+\eta\},\qquad j\geq1,
\]
are open, increase with $j$, and cover $X$. Compactness gives an index after which $u_t\leq f_j+\eta$ everywhere. Thus $u_t-\eta$ is a competitor for $p_j$, and $q\geq u_t-\eta$. Let $\eta\downarrow0$ and then $t\downarrow0$. We obtain $q\geq u$ almost everywhere, hence everywhere by the canonical representatives of quasi-plurisubharmonic functions. Taking the envelope over $u$ gives $q\geq p$.

For fixed $k$, the set $U_k=\{p<f_k\}$ is open. If $j\geq k$, then $p_j\leq p<f_k\leq f_j$ on $U_k$. The preceding continuous-obstacle theorem gives $\mathrm{MA}_\theta(p_j)=0$ there. On every relatively compact open subset of $X\setminus Z$, the functions $p_j$ have common bounds. Increasing local Bedford--Taylor continuity, applied after a common smooth correction and finite expansion, gives convergence of their measures to $\mathrm{MA}_\theta(p)$. The weak-limit inequality for open sets therefore gives zero measure on $U_k\setminus Z$. The union of these sets is $\{p<f\}\setminus Z$. The extended measure charges no pluripolar set, so it also charges no part of $Z$.
\end{proof}

Here and below the bounded local increasing continuity is the usual counterpart of decreasing Bedford--Taylor continuity from the source book. Upper regularization of an increasing limit changes it only on a pluripolar set. No assertion about the total mass of an unbounded increasing limit is being used.

\begin{proposition}\label{prop:big-bounded-quasicont-contact}
The conclusion of \cref{prop:big-bounded-lsc-contact} holds for every bounded quasi-continuous obstacle $f$.
\end{proposition}
\begin{proof}
By \cref{lem:quasicont-lsc-approximation}, there are bounded lower semicontinuous $f_j$ with common bounds and $f_j\downarrow f$ quasi-everywhere. The envelopes $p_j=P_\theta(f_j)$ decrease to $p=P_\theta(f)$. Indeed, the common lower barrier makes the limit genuine; the limit lies below $f$ quasi-everywhere and is thus an envelope competitor, while monotonicity gives the other inequality. All these functions have common local bounds off $Z$.

For fixed $k$, put $H_k=\min\{1,(f-p_k)_+\}$ on $X\setminus Z$. This is a bounded quasi-continuous function there. If $j\geq k$, then on $\{H_k>0\}$ we have $p_j\leq p_k<f\leq f_j$, outside a pluripolar set. The preceding proposition gives $H_k\mathrm{MA}_\theta(p_j)=0$ off $Z$.

Let $g$ be continuous with compact support in $X\setminus Z$. Bounded local product convergence holds against the fixed quasi-continuous coefficient $gH_k$. To verify this, remove an open set of small capacity where the coefficient may be discontinuous, extend it continuously from the compact complement, and use weak convergence there. Local uniform bounds for the potentials give uniform capacity control for the products, so the error on the removed set tends to zero. Consequently $H_k\mathrm{MA}_\theta(p)=0$ off $Z$. Letting $k\to\infty$ exhausts $\{p<f\}\setminus Z$. Non-pluripolarity removes $Z$.
\end{proof}

The remaining approximation needs this capacity argument for signed measures. We record the precise bounded fact to avoid imposing positivity where it is unavailable.

\begin{lemma}\label{lem:signed-truncation-quasicont}
Let $\Omega$ be Hermitian and $\theta$ any smooth real $(1,1)$-form. Suppose $v_j$ are $\Omega$-plurisubharmonic, have common finite upper and lower bounds, and decrease to $v$. Then the signed measures
\[
 \nu_j=(\theta+\ddc v_j)^n
\]
converge to $(\theta+\ddc v)^n$ against every fixed bounded quasi-continuous function. Their total variations are uniformly controlled by local capacity on a fixed finite cover.
\end{lemma}
\begin{proof}
On a relatively compact coordinate ball, choose a fixed quadratic $q$ such that $v_j+q$ are plurisubharmonic. Their common bounds remain uniform. Expand
\[
 (\theta+\ddc v_j)^n
 =\sum_{\ell=0}^n\binom n\ell
   \bigl(\ddc(v_j+q)\bigr)^\ell
       \wedge(\theta-\ddc q)^{n-\ell}.
\]
The absolute value of the $\ell$th summand is bounded by a fixed constant times
\[
 \bigl(\ddc(v_j+q)\bigr)^\ell\wedge\kappa_0^{n-\ell},
\]
where $\kappa_0$ is a fixed positive Euclidean form. This mixed positive measure is uniformly bounded by a constant times local capacity: add a fixed strictly plurisubharmonic quadratic to $v_j+q$, expand its top power, and rescale its uniformly bounded range into $[0,1]$ for the capacity definition. This also handles $\ell=0$. Summing on a finite cover gives the asserted variation control. Ordinary decreasing Bedford--Taylor continuity gives weak convergence term by term. For a bounded quasi-continuous coefficient, remove an open set of small capacity and extend the coefficient continuously from its compact complement. The variation bound controls the error uniformly for the sequence and its limit. Letting that capacity tend to zero proves the assertion.
\end{proof}

\Needspace{12\baselineskip}
\begin{theorem}
\label{thm:big-general-envelope-contact}
Let $\theta$ be smooth and big. Let $f$ be a quasi-continuous obstacle, finite almost everywhere, and suppose $p=P_\theta(f)$ is genuine. Then $p\leq f$ quasi-everywhere and
\[
 \mathbf1_{\{p<f\}}\mathrm{MA}_\theta(p)=0.
\]
The background need not be semipositive. No bounded mass or positive lower-volume property is assumed, and the measure may have infinite total mass.
\end{theorem}
\begin{proof}
Set $f_j=\max(f,-j)$ and $p_j=P_\theta(f_j)$. These envelopes are genuine because $p$ is a competitor. They have a common finite upper bound. Here is the elementary reason: $f_1$ is finite almost everywhere, so some set $A=\{f_1\leq N\}$ has positive smooth volume. For any competitor $a\leq f_1$ quasi-everywhere, the uniform $L^1$ bound for $a-\sup_Xa$ gives
\[
 \operatorname{Vol}(A)(\sup_Xa-N)
       \leq\int_A(\sup_Xa-a)\,\mathrm{d}V\leq C.
\]
The same bound applies to its envelope, and hence to all $p_j$. Choose $M$ strictly larger than this common upper bound.

The obstacle $g_j=\min(f_j,M)$ is bounded and quasi-continuous. Since $p_j<M$ and $p_j\leq f_j$ quasi-everywhere, $p_j$ is a competitor for $P_\theta(g_j)$; the reverse inequality follows from $g_j\leq f_j$. Thus $p_j=P_\theta(g_j)$. The preceding bounded result gives contact support for $p_j$ at $g_j$. The strict inequality $p_j<M$ excludes contact at the artificial upper cutoff, so
\begin{equation}\label{eq:big-truncated-obstacle-contact}
 \mathbf1_{\{p_j<f_j\}}\mathrm{MA}_\theta(p_j)=0.
\end{equation}
Each $p_j$ is locally bounded off $Z$, with a lower barrier of the form $\rho-C_j$. The sequence decreases to $p$: its limit is genuine since it dominates $p$, lies below $f$ quasi-everywhere, and is therefore itself an envelope competitor.

Fix a Hermitian form $\Omega\geq\theta$. For integers $C,k\geq1$, define
\[
 v_{j,C}=\max(p_j,-C),\qquad v_C=\max(p,-C),
\]
which are bounded $\Omega$-plurisubharmonic functions, and set
\[
 H_{k,C}
  =\min\{1,(f-p_k)_+\}\min\{1,(p+C)_+\}
       \quad\hbox{on }X\setminus Z.
\]
Extend $H_{k,C}$ by zero on $Z$. It is bounded and quasi-continuous: off $Z$ the first difference has a finite quasi-plurisubharmonic second term; bounded continuous truncations preserve quasi-continuity, and the pluripolar set $Z$ can be removed with arbitrarily small capacity. At a pole of $p$ the second factor is zero.

For $j\geq k$, on $\{H_{k,C}>0\}$ we have, outside irrelevant pluripolar sets,
\[
 p_j\geq p>-C,\qquad p_j\leq p_k<f\leq f_j.
\]
The signed measure $\nu_{j,C}=(\theta+\ddc v_{j,C})^n$ therefore agrees there with $\mathrm{MA}_\theta(p_j)$ by plurifine locality. Equation \eqref{eq:big-truncated-obstacle-contact} gives
\[
 H_{k,C}\nu_{j,C}=0.
\]
For fixed $C$, the potentials $v_{j,C}$ have common global bounds and decrease to $v_C$. The signed convergence lemma yields
\[
 \int_X H_{k,C}(\theta+\ddc v_C)^n=0.
\]
On the set where $H_{k,C}$ is positive, $p>-C$. The restricted product in \cref{thm:arbitrary-extended-ma} identifies the last integrand with $H_{k,C}\mathrm{MA}_\theta(p)$, a positive measure; on the complement the coefficient is zero. It follows that $\mathrm{MA}_\theta(p)$ vanishes on $\{H_{k,C}>0\}$.

If $p(x)$ is finite and $p(x)<f(x)$, then $p_k(x)<f(x)$ for large $k$, and $p(x)>-C$ for large $C$. The countable union of these positive sets therefore covers $\{p<f\}$ outside $Z\cup\{p=-\infty\}$. Both excluded sets are pluripolar. This proves the theorem.
\end{proof}

Every passage to a limit in the last argument takes place at a fixed truncation level. Its capacity bound may depend on that level. We never sum such bounds or assume that the resulting extended measure is finite. This is why the argument applies when bounded mass fails.


\section{Compactness with a common lower barrier}

A common lower barrier makes all finite and infinite rooftops genuine. An error integral against each approximating measure bounds the rooftop error by a summable tail. Weighted lower semicontinuity passes this bound to the limit, and singular domination identifies the lower envelope with the $L^1$ limit. Sandwiching by upper envelopes then gives capacity convergence; suitable domination also controls the measure tails.

\Needspace{16\baselineskip}
\begin{theorem}\label{thm:common-barrier-criterion}\index{Barrier!common lower bound}
Let $\theta\geq0$ be big. Suppose $\varphi\in\PSH(X,\theta)$ admits a genuine $P_{\delta\theta}(\varphi)$ for some $0<\delta<1$. Let $u_j\in\PSH(X,\theta)$ satisfy, for a fixed $C$,
\[
 \varphi-C\leq u_j\leq\varphi+C,\qquad u_j\longrightarrow u\text{ in }L^1(X).
\]
If
\begin{equation}\label{eq:common-barrier-self-error}
 I_j:=\int_X|\mathrm{e}^{u_j}-\mathrm{e}^u|\,
               \mathrm{MA}^{\mathrm{ext}}_\theta(u_j)\longrightarrow0,
\end{equation}
then $u_j\to u$ in capacity.
\end{theorem}
\begin{proof}
The canonical limit lies between $\varphi-C$ and $\varphi+C$. It suffices to prove capacity convergence after further extraction from an arbitrary subsequence. Choose an extraction with $I_j\leq2^{-j}$ and almost-everywhere convergence for smooth volume. For $N\geq j$ define
\[
 p_{j,N}=P_\theta(u_j,\ldots,u_N),\qquad
 p_j=P_\theta\left(\inf_{\ell\geq j}u_\ell\right).
\]
\Needspace{8\baselineskip}
All these envelopes lie between $\varphi-C$ and $\varphi+C$. The contact theorem and ordered contact give a finite-rooftop inequality: partition the contact set according to the first index $\ell$ for which $p_{j,N}=u_\ell$. On that piece the measure of $p_{j,N}$ is bounded by that of $u_\ell$, and $\mathrm{e}^{p_{j,N}}=\mathrm{e}^{u_\ell}$. Therefore
\[
 \int_X|\mathrm{e}^{p_{j,N}}-\mathrm{e}^u|\,
            \mathrm{MA}^{\mathrm{ext}}_\theta(p_{j,N})
 \leq\sum_{\ell=j}^N I_\ell\leq2^{1-j}.
\]
This uses only positive measures, including when their total masses are infinite.

As $N\to\infty$, $p_{j,N}\downarrow p_j$; the limit is the envelope of the infimum because every competitor is below every finite rooftop. The sequence $p_j$ increases to a genuine canonical limit $p$. Both passages hold in capacity by \cref{lem:monotone-psh-capacity}. Apply \cref{lem:weighted-extended-lsc} to the bounded nonnegative weights $|\mathrm{e}^{p_{j,N}}-\mathrm{e}^u|$, first in $N$ and then in $j$. We obtain
\[
 \int_X|\mathrm{e}^p-\mathrm{e}^u|\,\mathrm{MA}^{\mathrm{ext}}_\theta(p)=0.
\]
The inequalities $p_j\leq u_\ell$ for all $\ell\geq j$ imply $p\leq u$, first almost everywhere and then for canonical representatives. Also $p\sim u\sim\varphi$. Off their poles, $\mathrm{e}^p<\mathrm{e}^u$ exactly on $\{p<u\}$, so
$\mathrm{MA}^{\mathrm{ext}}_\theta(p)(\{p<u\})=0$.
The competitor $P_{\delta\theta}(\varphi)-C$ makes $P_{\delta\theta}(p)$ genuine. Apply \cref{thm:singular-strict-domination} with coefficient zero to obtain $p\geq u$, and hence $p=u$.

Finally $U_j=(\sup_{\ell\geq j}u_\ell)^*$ decreases to $u$ by Hartogs' lemma and the almost-everywhere extraction. The inequalities
$p_j\leq u_j\leq U_j$, together with the monotone capacity convergence of both outer sequences, prove the claim.
\end{proof}

\Needspace{10\baselineskip}
\begin{proposition}\label{prop:lower-barrier-criterion}
In \cref{thm:common-barrier-criterion}, one may replace the two bounds on $u_j$ by
\[
 \varphi-C\leq u_j\leq C
\]
provided the $L^1$ limit satisfies $u\preceq\varphi$. In particular the approximants may be bounded while $\varphi$ is unbounded.
\end{proposition}
\begin{proof}
The $L^1$ limit satisfies $u\geq\varphi-C$, so the additional hypothesis implies $u\sim\varphi$. Repeat the finite rooftop construction in the theorem. Every $p_{j,N}$ remains genuine and lies between $\varphi-C$ and $C$, which is enough for contact support and for the bounded exponential weights. The resulting inequality is still
\[
 \int_X|\mathrm{e}^{p_{j,N}}-\mathrm{e}^u|\,
             \mathrm{MA}^{\mathrm{ext}}_\theta(p_{j,N})\leq2^{1-j}.
\]
The limits $p_j$ exist by the common lower bound, and $p_j\leq u$ follows from $p_j\leq u_\ell$ for every $\ell\geq j$ and almost-everywhere convergence. Thus
$\varphi-C\leq p_j\leq u$, and all $p_j$ and their increasing limit $p$ have the same ordinary type as $\varphi$. Monotone capacity convergence and weighted lower semicontinuity give
$\mathrm{MA}^{\mathrm{ext}}_\theta(p)(\{p<u\})=0$. Singular domination applies to $(p,u)$ exactly as before, so $p=u$. The upper tail envelopes decrease to $u$, and the same sandwich gives capacity convergence. No singularity-type assertion was needed for the finite rooftops or for the individual approximants.
\end{proof}

\Needspace{12\baselineskip}
\begin{corollary}\label{cor:common-barrier-continuity}
Under the geometric and common-barrier hypotheses of \cref{thm:common-barrier-criterion}, suppose either
\begin{enumerate}
\item $\mathrm{MA}^{\mathrm{ext}}_\theta(u_j)\leq\nu$ for a fixed finite non-pluripolar measure $\nu$; or
\item $\mathrm{MA}^{\mathrm{ext}}_\theta(u_j)\leq A\omega_{w_j}^n$, where $A$ is fixed and $w_j\in\PSH(X,\omega)$ are uniformly bounded and converge in capacity to a bounded $w$.
\end{enumerate}
Then $u_j\to u$ in capacity, all the measures have finite mass, and
\[
 \mathrm{MA}^{\mathrm{ext}}_\theta(u_j)
 \rightharpoonup\mathrm{MA}^{\mathrm{ext}}_\theta(u).
\]
\end{corollary}
\begin{proof}
In the first case \cref{lem:fixed-npp-integrals} applied to exponentials proves \eqref{eq:common-barrier-self-error}. In the second case use \cref{lem:moving-bounded-integrals}, since the exponentials are bounded and have a common lower curvature bound. The criterion gives capacity convergence in both cases.

The masses are uniformly bounded by the corresponding dominating measures. For the deep tails use $u_j\geq\varphi-C$. In the first case their masses are at most $\nu(\{\varphi<C-T\})$, which tends to zero because $\nu$ ignores the pole set of $\varphi$. In the second case \cref{lem:bounded-mixed-capacity} bounds them by
$C'\operatorname{Cap}_\omega(\{\varphi<C-T\})$, which also tends to zero. Apply the raw-tail continuity theorem \cref{thm:undamped-tail-continuity} in its general-background form. It gives finiteness and weak convergence for the limiting measure as well.
\end{proof}

Its reference need not be an envelope fixed point; the smaller-background condition is enough. A model reference becomes useful when an a priori estimate supplies the two uniform barriers. If such a lower barrier is unavailable, constructing one in a relative envelope class is an additional problem, treated in \cref{sec:general-measure-relative}.


\section{A bounded subsolution theorem}

A bounded density relative to the Monge--Amp\`ere measure of one bounded potential also admits a bounded exponential solution. This permits right-hand sides singular with respect to smooth volume. Regularizing the reference gives a moving dominating measure; the mixed density inequality and fixed-reference trace estimate of \cref{lem:nonclosed-mixed-density,lem:bounded-reference-trace} bound the normalizations when the density is allowed to vanish.

\Needspace{15\baselineskip}
\begin{theorem}\label{thm:bounded-subsolution}
Let $\theta$ be semipositive and big. Suppose
\[
 \mu=f\,\mathrm{MA}_\theta(\psi),\qquad
 \psi\in\PSH(X,\theta)\cap L^\infty(X),\qquad
 0\leq f\in L^\infty(\mathrm{MA}_\theta(\psi)),
\]
and $\mu(X)>0$. For every $\lambda>0$ there is a unique bounded $u\in\PSH(X,\theta)$ satisfying
\[
 \mathrm{MA}_\theta(u)=\mathrm{e}^{\lambda u}\mu.
\]

\end{theorem}
\begin{proof}
Write $\nu=\mathrm{MA}_\theta(\psi)$, a finite non-pluripolar measure.

\emph{Smooth strictly positive density.} First suppose $f$ is smooth and strictly positive. Choose a Hermitian $\Omega\geq\theta$ and smooth $\Omega$-plurisubharmonic $\psi_j\downarrow\psi$, using the bounded regularization theorem. The rooftops
\[
 q_j=P_\theta(\psi_j)
\]
are uniformly bounded and decrease to $\psi$: they lie between $\psi$ and $\psi_j$. Contact support and ordered contact in background $\Omega$ give
\[
 \mathrm{MA}_\theta(q_j)
 \leq\mathbf1_{\{q_j=\psi_j\}}\Omega_{\psi_j}^n.
\]
Each right side has a bounded density with respect to smooth volume. Moreover $\mathrm{MA}_\theta(q_j)(X)>0$ by individual noncollapse. The bounded big-background theorem therefore supplies $u_j$ with
\[
 \mathrm{MA}_\theta(u_j)=\mathrm{e}^{\lambda u_j}f\,
                                     \mathrm{MA}_\theta(q_j).
\]
Since $f$ is bounded above and away from zero, $q_j-C$ and $q_j+C$ are lower and upper barriers for a fixed $C$. More explicitly choose $C$ so that
$\mathrm{e}^{\lambda(q_j-C)}f\leq1\leq \mathrm{e}^{\lambda(q_j+C)}f$ for all $j$. Exponential order then gives
$q_j-C\leq u_j\leq q_j+C$.

Extract an $L^1$ limit $u$, still uniformly bounded. The moving domination
\[
 \mathrm{MA}_\theta(u_j)\leq C'\Omega_{\psi_j}^n
\]
and \cref{cor:common-barrier-continuity}, with reference zero, give convergence in capacity. Bounded capacity continuity with the quasi-continuous weights $\mathrm{e}^{\lambda u_j}f$ gives
\[
 \mathrm{MA}_\theta(u)=\mathrm{e}^{\lambda u}f\nu.
\]

\emph{Bounded measurable density.} Choose positive smooth $f_j$ with a common upper bound $L>0$, converging to $f$ in $L^1(\nu)$. Such an approximation follows from regularity of the finite measure, continuous approximation on large compact sets, and smooth uniform approximation of the continuous functions; truncation to a fixed bounded interval and addition of a positive constant tending to zero preserve the bounds and convergence. The first step gives bounded $u_j$ with
\[
 \mathrm{MA}_\theta(u_j)=\mathrm{e}^{\lambda u_j}f_j\nu.
\]
A fixed $\psi-C$ is a lower barrier, since
$\mathrm{e}^{\lambda(\psi-C)}f_j\leq1$ for all $j$. Hence $u_j\geq\psi-C$. The upper bound requires a separate argument because $\inf_Xf_j$ can tend to zero.

The construction in the first step also gives the mixed inequality
\begin{equation}\label{eq:bounded-subsolution-mixed}
 \theta_{u_j}\wedge\theta_\psi^{n-1}
 \geq \mathrm{e}^{\lambda u_j/n}f_j^{1/n}\nu.
\end{equation}
To see this without assuming a mixed inequality for arbitrary singular measures, use that step's approximants $u_{j,k}$ and $q_k$. Their top measures have smooth-volume densities and
\[
 \mathrm{MA}_\theta(u_{j,k})
     =\mathrm{e}^{\lambda u_{j,k}}f_j\mathrm{MA}_\theta(q_k).
\]
Apply \cref{lem:nonclosed-mixed-density} to $(u_{j,k},q_k)$. For this fixed $j$ both sequences are uniformly bounded and converge in capacity to $(u_j,\psi)$. Bounded weighted continuity passes the mixed inequality to the limit and yields \eqref{eq:bounded-subsolution-mixed}.

Put $M_j=\sup_Xu_j$, $v_j=u_j-M_j$ and $c_j=\mathrm{e}^{\lambda M_j}$. If the $M_j$ were unbounded above, extract $M_j\to+\infty$ and then $v_j\to v$ in $L^1(X)$, with $v$ a genuine normalized potential. By \cref{lem:fixed-npp-integrals},
$\mathrm{e}^{\lambda v_j/n}\to \mathrm{e}^{\lambda v/n}$ in $L^1(\nu)$. The inequalities
\[
 |f_j^{1/n}-f^{1/n}|^n\leq|f_j-f|,
 \qquad0\leq f_j\leq L,
\]
then give
\[
 \int_X \mathrm{e}^{\lambda v_j/n}f_j^{1/n}\,\mathrm{d}\nu
 \longrightarrow\int_X\mathrm{e}^{\lambda v/n}f^{1/n}\,\mathrm{d}\nu>0.
\]
The last integral is positive because $\mu(X)>0$ and $\nu$ gives no mass to the pole set of $v$. Integrate \eqref{eq:bounded-subsolution-mixed} and apply \cref{lem:bounded-reference-trace}. For some $a>0$ along the extraction, this gives
\[
 a c_j^{1/n}\leq
 \int_X\theta_{v_j}\wedge\theta_\psi^{n-1}\leq C_\psi,
\]
contradicting $M_j\to+\infty$. Thus $u_j$ is uniformly bounded above as well as below.

Extract an $L^1$ limit. Now $\mathrm{MA}_\theta(u_j)\leq C''\nu$, so \cref{cor:common-barrier-continuity} gives capacity and weak Monge--Amp\`ere convergence. The right sides converge in total variation: use the fixed-measure exponential convergence and
\[
 \int_X|\mathrm{e}^{\lambda u_j}f_j-\mathrm{e}^{\lambda u}f|\,\mathrm{d}\nu
 \leq C\|f_j-f\|_{L^1(\nu)}
       +\|f\|_\infty\int_X|\mathrm{e}^{\lambda u_j}-\mathrm{e}^{\lambda u}|\,\mathrm{d}\nu.
\]
Consequently $\mathrm{MA}_\theta(u)=\mathrm{e}^{\lambda u}\mu$. Uniqueness is bounded exponential order.
\end{proof}


\section{Adding weighted Monge--Amp\`ere measures}

Weighted measures of two solutions in the same ordinary type can be added while retaining that type. Truncation gives bounded equations whose solutions have a common lower barrier. Their exponential error tends to zero even if the approximating masses grow, so the lower-barrier criterion identifies the limit. This addition step will remove a measure deficit in the prescribed-type existence proof.

\Needspace{15\baselineskip}
\begin{theorem}\label{thm:weighted-measure-addition}\index{Monge--Ampere measure@Monge--Amp\`ere measure!weighted addition}
Let $\theta\geq0$ be big, and suppose $\varphi\in\PSH(X,\theta)$ admits $P_{\delta\theta}(\varphi)\not\equiv-\infty$ for some $0<\delta<1$. Let $u,v\in\PSH(X,\theta)$ satisfy $u,v\sim\varphi$, and let $\lambda>0$. If
\[
 \mu=\mathrm{e}^{-\lambda u}\mathrm{MA}^{\mathrm{ext}}_\theta(u)
       +\mathrm{e}^{-\lambda v}\mathrm{MA}^{\mathrm{ext}}_\theta(v)
\]
is a finite measure, then there exists $a\in\PSH(X,\theta)$, with $a\sim\varphi$, such that
\[
 \mathrm{MA}^{\mathrm{ext}}_\theta(a)=\mathrm{e}^{\lambda a}\mu.
\]
Weights are interpreted off the pluripolar pole sets. The solution is unique in this ordinary singularity type.
\end{theorem}
\begin{proof}
Choose $S\geq\max(\sup_Xu,\sup_Xv)$. A common translation reduces to $u,v\leq0$: translating by $-S$ multiplies $\mu$ by $\mathrm{e}^{\lambda S}$ and translates the resulting solution back by $S$.

For integers $j\geq1$, set $u_j=\max(u,-j)$, $v_j=\max(v,-j)$ and
\[
 \mu_j=\mathrm{e}^{-\lambda u_j}\mathrm{MA}_\theta(u_j)
           +\mathrm{e}^{-\lambda v_j}\mathrm{MA}_\theta(v_j).
\]
For $s_j=(u_j+v_j)/2$, positivity gives
\[
 \mu_j\leq2^n \mathrm{e}^{\lambda j}\mathrm{MA}_\theta(s_j).
\]
Each measure is finite and nonzero by individual bounded noncollapse. The bounded subsolution theorem gives $a_j$ with
$\mathrm{MA}_\theta(a_j)=\mathrm{e}^{\lambda a_j}\mu_j$.
Bounded exponential order yields
\[
 a_j\leq\min(u_j,v_j)\leq0.
\]
Choose $C$ with $|u-v|\leq C$ off the common poles; then $|u_j-v_j|\leq C$. Put
\[
 D=C/2+n\lambda^{-1}\log2,\qquad b_j=s_j-D.
\]
The inequality
\[
 \mathrm{e}^{-\lambda b_j}\mathrm{MA}_\theta(b_j)
 \geq2^{-n}\mathrm{e}^{\lambda D-\lambda s_j}
            \bigl(\mathrm{MA}_\theta(u_j)+\mathrm{MA}_\theta(v_j)\bigr)
 \geq\mu_j
\]
follows from $|u_j-v_j|\leq C$. Hence $a_j\geq b_j\geq w$, where
\[
 w=(u+v)/2-D\sim\varphi
\]
is a fixed genuine potential. Extract $a_j\to a$ in $L^1(X)$. The upper barriers decrease to $\min(u,v)$, so $w\leq a\leq\min(u,v)$ and $a\sim\varphi$.

To prove convergence in capacity, write $E_j=\{\min(u,v)>-j\}$. Locality gives $\mu_j=\mu$ there. Thus
\[
 \int_{E_j}|\mathrm{e}^{a_j}-\mathrm{e}^a|\,\mathrm{MA}_\theta(a_j)
 \leq\int_X|\mathrm{e}^{a_j}-\mathrm{e}^a|\,\mathrm{d}\mu\longrightarrow0
\]
by \cref{lem:fixed-npp-integrals}. On $X\setminus E_j$ we have $a_j\leq-j$ and $a\leq-j$, whence $|\mathrm{e}^{a_j}-\mathrm{e}^a|\leq \mathrm{e}^{-j}$. Also
\[
 \mathrm{MA}_\theta(a_j)
 =\mathrm{e}^{\lambda(a_j-u_j)}\mathrm{MA}_\theta(u_j)
    +\mathrm{e}^{\lambda(a_j-v_j)}\mathrm{MA}_\theta(v_j)
 \leq\mathrm{MA}_\theta(u_j)+\mathrm{MA}_\theta(v_j).
\]
Choose a Hermitian $\Omega\geq\theta$. Write $r_n=\max(n-2,0)$. The bounded-range estimate gives
\[
 \int_{X\setminus E_j}|\mathrm{e}^{a_j}-\mathrm{e}^a|\,\mathrm{MA}_\theta(a_j)
 \leq C'\mathrm{e}^{-j}(1+j)^{r_n}\longrightarrow0.
\]
The lower-barrier criterion \cref{prop:lower-barrier-criterion} now gives $a_j\to a$ in capacity. Only polynomial control of the approximating masses was needed.

It remains to identify the equation. This can be done on bounded levels without claiming raw weak convergence prematurely. For $T>0$ choose the continuous cutoff $g_T$ from \cref{lem:weighted-extended-lsc} and put $\zeta_T=g_T(w)$. Where $\zeta_T>0$, we have $w>-2T$, hence every $a_j$ and $a$ is above $-2T$. If $j>2T$ and $x\notin E_j$, the inequalities $w(x)\leq a_j(x)\leq-j$ force $\zeta_T(x)=0$. Consequently, for all $j>2T$,
\[
 \zeta_T\mathrm{MA}_\theta(a_j)
       =\zeta_T \mathrm{e}^{\lambda a_j}\mu.
\]
Test this identity against any continuous function. On the left, replace $a_j$ by $\max(a_j,-2T)$ and use bounded capacity continuity with the fixed bounded quasi-continuous weight $\zeta_T$. On the right, fixed-measure exponential convergence gives convergence in total variation. Therefore
\[
 \zeta_T\mathrm{MA}^{\mathrm{ext}}_\theta(a)
       =\zeta_T\mathrm{e}^{\lambda a}\mu.
\]
Let $T\to\infty$. The cutoffs increase to one off the pole set of the genuine $w$; all measures ignore that set. This proves the equation and, in particular, finiteness of the limiting raw measure. Uniqueness follows from singular exponential order.
\end{proof}


\section{Comparison above a strict singular barrier}
\label{sec:big-minimal-comparison}

Let $\theta$ now be any smooth big real $(1,1)$-form. A strict analytic barrier replaces the smaller-background construction used above. It confines a prospective negative minimum away from its pole set; a local comparison then rules out that minimum.

Fix $\rho\leq0$ with analytic singularities and smooth remainder, and decrease the Hermitian metric if necessary so that $\theta_\rho\geq\omega$. Put $Z=\{\rho=-\infty\}$ and $\Omega=X\setminus Z$. These data remain fixed through \cref{sec:nonclosed-interior-smoothness}. A potential bounded below by $\rho-C$ is locally bounded on $\Omega$, where its extended measure agrees with its bounded product.

\begin{lemma}\label{lem:big-local-comparison}
Let $B\Subset\Omega$ be a coordinate ball, and let $u,v$ be bounded $\theta$-plurisubharmonic functions on $B$. Assume
\[
 \liminf_{x\to\partial B}(u-v)(x)\geq0,
 \qquad \mathbf1_{\{u<v\}}\mathrm{MA}_\theta(u)\leq
             \mathbf1_{\{u<v\}}\mathrm{MA}_\theta(v).
\]
Then $u\geq v$ on $B$.
\end{lemma}
\begin{proof}
Fix $\varepsilon>0$ and set $v_\varepsilon=\max(u,v-\varepsilon)$. It equals $u$ near the boundary. Locality on $\{u>v-\varepsilon\}$ and ordered contact on $\{u\leq v-\varepsilon\}$ show that
\[
 \mathrm{MA}_\theta(u)\leq\mathrm{MA}_\theta(v_\varepsilon)
\]
throughout $B$: the latter set is contained in $\{u<v\}$, where the assumed inequality applies.

Let $q$ be the ordinary plurisubharmonic envelope on $B$ of $u-v_\varepsilon$. This bounded quasi-continuous obstacle is zero near the boundary and lies between $Ah$ and zero for a sufficiently large $A$, where $h<0$ is a smooth strictly plurisubharmonic defining function for the ball. Thus $q$ tends to zero at the boundary, and \eqref{eq:local-envelope-contact} shows that $(\ddc q)^n$ is carried by $D=\{q=u-v_\varepsilon\}$.

Canonical representatives give $q+v_\varepsilon\leq u$. Ordered contact, followed by the preceding measure inequality, yields
\[
 \mathbf1_D\mathrm{MA}_\theta(q+v_\varepsilon)
 \leq\mathbf1_D\mathrm{MA}_\theta(u)
 \leq\mathbf1_D\mathrm{MA}_\theta(v_\varepsilon).
\]
The positive expansion of $(\theta_{v_\varepsilon}+\ddc q)^n$ contains both $\mathrm{MA}_\theta(v_\varepsilon)$ and $(\ddc q)^n$. Thus $\mathbf1_D(\ddc q)^n=0$, and contact support gives $(\ddc q)^n=0$. The classical bounded Bedford--Taylor comparison principle with zero boundary values implies $q=0$. Consequently $u\geq v_\varepsilon$, first quasi-everywhere and then everywhere. Let $\varepsilon\downarrow0$.
\end{proof}

\begin{theorem}\label{thm:big-barrier-domination}
Suppose $u,v\in\PSH(X,\theta)$ satisfy, for constants $C_0,C_1\geq0$,
\[
 \min(u,v)\geq\rho-C_0,\qquad v\leq u+C_1.
\]
If, for some $0\leq c<1$,
\[
 \mathbf1_{\{u<v\}}\mathrm{MA}_\theta(u)\leq
 c\,\mathbf1_{\{u<v\}}\mathrm{MA}_\theta(v),
\]
then $u\geq v$. The measures need not have finite total mass.
\end{theorem}
\begin{proof}
Replace $\rho$ by $\widehat\rho=a\rho-A$, where $a>1$ is sufficiently close to one and $A$ is sufficiently large. Since
\[
 \theta+\ddc(a\rho)=a\theta_\rho+(1-a)\theta,
\]
we can arrange $\theta_{\widehat\rho}\geq\omega/2$, $\widehat\rho-v\leq-1$, and
$u-\widehat\rho\to+\infty$ when approaching $Z$. When $Z$ is empty, the last boundary condition is unnecessary. In the proof below write $r=\widehat\rho$.

Choose $0<\delta<1$ with $(1-\delta)^n>c$, and put $v_\delta=(1-\delta)v+\delta r$. Then
\[
 u-v_\delta\geq\delta(u-r)-(1-\delta)C_1.
\]
If its infimum $m$ on $\Omega$ were negative, a sequence $x_j$ approaching that infimum would have a subsequence converging to a point $x_0\in\Omega$. Choose a coordinate ball $B\Subset\Omega$ centered at $x_0$, identify it with $\{|z|<1\}$, and take
\[
 \eta(z)=b(1-|z|^2),\qquad 0<b<1,
 \qquad \omega/2+\ddc\eta\geq0.
\]
The function $w=v_\delta+\delta\eta+m$ is $\theta$-plurisubharmonic on $B$, and $\liminf_{\partial B}(u-w)\geq0$. Also
\[
 w-v=\delta(r-v+\eta)+m<0.
\]
Thus $\{u<w\}\subset\{u<v\}$, and positivity gives on this set
\[
 \mathrm{MA}_\theta(u)\leq c\mathrm{MA}_\theta(v)\leq(1-\delta)^n\mathrm{MA}_\theta(v)\leq\mathrm{MA}_\theta(w).
\]
The local comparison lemma implies $u\geq w$ on $B$. Evaluating at $x_j$ gives $u(x_j)-v_\delta(x_j)\geq m+\delta\eta(x_j)$, contrary to $\eta(x_0)>0$. We have proved $u\geq v_\delta$. Let $\delta\downarrow0$ along values with $(1-\delta)^n>c$, and extend the resulting inequality from $\Omega$ by canonical representatives.
\end{proof}

\Needspace{20\baselineskip}
\begin{corollary}\label{cor:big-barrier-orders}
Let $u,v\in\PSH(X,\theta)$ satisfy $u,v\geq\rho-C$ for some $C\geq0$ and $u\sim v$.
\begin{enumerate}
\item If $\mathrm{MA}_\theta(u)\leq c\mathrm{MA}_\theta(v)$ for $c\geq0$, then $c\geq1$.
\item If $\mathrm{MA}_\theta(u)\leq \mathrm{e}^{\lambda(u-v)}\mathrm{MA}_\theta(v)$ for $\lambda>0$, then $u\geq v$.
\end{enumerate}
In particular the positive-exponent equation has at most one solution in each such ordinary singularity type, and the normalized equation has at most one positive scalar in that type.
\end{corollary}
\begin{proof}
If $c<1$, apply the theorem to $u,v+A$ for every constant $A>0$. Its lower-barrier constant may depend on $A$, which causes no difficulty. The conclusion $u\geq v+A$ is impossible for all $A$. For the second assertion, fix $a>0$ and restrict the measure inequality to $\{u<v-a\}$. There the factor is at most $\mathrm{e}^{-\lambda a}<1$, so the theorem gives $u\geq v-a$. Let $a\downarrow0$. Apply this assertion in both orders for two exponential solutions. Apply the first assertion in both orders for two normalized solutions.
\end{proof}

\section{Subsolutions and the Gauduchon trace}
\label{sec:big-minimal-estimates}

The strict barrier also lets us construct a subsolution within a bounded distance of any reference between $\rho$ and zero. A convex combination of positive currents gives the curvature estimate. Pairing the equation with a Gauduchon metric then bounds its normalization through a linear trace.

\begin{lemma}\label{lem:big-reference-subsolution}
Fix $p>1$ and $L>0$. There is $m>0$ such that, whenever
\[
 \rho\leq\eta\leq0,\quad\eta\in\PSH(X,\theta),
 \qquad g\geq0,\quad\|g\|_{L^p(\mathrm{d}V)}\leq L,
\]
there exists $v\in\PSH(X,\theta)$ with
\[
 \eta\leq v\leq\eta+1,\qquad
 \mathrm{MA}_\theta(v)\geq m g\,\mathrm{d}V.
\]
The constant depends only on the fixed geometry, $\rho,p,L$.
\end{lemma}
\begin{proof}
Choose $1<q<p$. The functions $h=g(2-\rho)^{2n}$ have uniformly bounded $L^q$ norms, by H\"older's inequality and the finite polynomial moments of a fixed quasi-plurisubharmonic function. Apply \cref{lem:big-lp-subsolution} with the Hermitian background $\omega$. It gives $-1\leq s\leq0$, $s\in\PSH(X,\omega)$, and
\[
 \omega_s^n\geq m h\,\mathrm{d}V
\]
for a fixed $m>0$. Put $u_1=\rho+s$. Then $\theta_{u_1}\geq\omega_s$. On $\Omega$ define
\[
 t=1+\eta-u_1,\qquad v=\eta+t^{-1}.
\]
Here $1\leq t\leq2-\rho$. Apply \cref{lem:bounded-convex-pair} on relatively compact subsets of $\Omega$, with $a=\eta$, $h=u_1$ and $\chi(s)=(1-s)^{-1}$. It gives
\[
 \theta_v=(1-t^{-2})\theta_\eta+t^{-2}\theta_{u_1}
       +2t^{-3}\,\mathrm{d}(\eta-u_1)\wedge\dc(\eta-u_1)
 \geq t^{-2}\theta_{u_1}.
\]
Since $\eta\leq v\leq\eta+1$, \cref{lem:qpsh-removable-analytic} extends $v$ across $Z$, and the bounds extend by \cref{lem:qpsh-almost-everywhere-order}. On $\Omega$, positivity of the local products gives
\[
 \mathrm{MA}_\theta(v)\geq t^{-2n}\theta_{u_1}^n
       \geq(2-\rho)^{-2n}\omega_s^n
       \geq m g\,\mathrm{d}V.
\]
Both sides give zero mass to $Z$, proving the global inequality.
\end{proof}

\begin{proposition}\label{prop:big-singular-gauduchon-trace}
Let $\gamma$ be Gauduchon, write $\gamma^n=G\,\mathrm{d}V$, and put
\[
 T_\gamma=\int_X\theta\wedge\gamma^{n-1}>0.
\]
If $u\geq\rho-C$ is $\theta$-plurisubharmonic and $\mathrm{MA}_\theta(u)\geq F\,\mathrm{d}V$, $F\geq0$, then
\begin{equation}\label{eq:big-singular-gauduchon-trace}
 \int_X F^{1/n}G^{1-1/n}\,\mathrm{d}V\leq T_\gamma.
\end{equation}
Consequently, for a fixed nonzero $0\leq f\in L^p(\mathrm{d}V)$, $p>1$, solutions above this barrier of
$\mathrm{MA}_\theta(u)=\mathrm{e}^{\lambda u}f\,\mathrm{d}V$ have uniformly bounded suprema when $\lambda>0$ is fixed. If instead $v_j\in\PSH(X,\theta)$, $j\geq1$, and constants $C_j\geq0$ satisfy
\[
 \sup_Xv_j=0,\qquad
 \mathrm{MA}_\theta(v_j)=c_j \mathrm{e}^{v_j/j}f\,\mathrm{d}V,
 \qquad v_j\geq\rho-C_j,
\]
with $c_j>0$, then the constants $c_j$ are uniformly bounded above.
\end{proposition}
\begin{proof}
The linear current $\theta_u\wedge\gamma^{n-1}$ is positive and globally defined, even at $Z$. Distributional Stokes gives
\[
 \int_X\theta_u\wedge\gamma^{n-1}
 =\int_X\theta\wedge\gamma^{n-1},
\]
since $\ddc\gamma^{n-1}=0$. Applying the identity to $\rho$ proves $T_\gamma>0$.

On $\Omega$, apply \cref{lem:interior-mixed-background} to $\theta_u$ and $\gamma$ to obtain
\[
 \theta_u\wedge\gamma^{n-1}\geq
 F^{1/n}G^{1-1/n}\,\mathrm{d}V.
\]
Exhaust $\Omega$ and use positivity of the global trace measure. The complement has zero smooth volume, giving \eqref{eq:big-singular-gauduchon-trace}.

\Needspace{5\baselineskip}
For the exponential assertion write $a=\sup_Xu$ and $w=u-a$. The estimate becomes
\[
 \mathrm{e}^{\lambda a/n}\int_X\mathrm{e}^{\lambda w/n}f^{1/n}G^{1-1/n}\,\mathrm{d}V
 \leq T_\gamma.
\]
The integral has a positive lower bound over all normalized $\theta$-plurisubharmonic $w$. Otherwise normalized compactness supplies a subsequence converging almost everywhere to a genuine potential, and dominated convergence gives a strictly positive limit integral, a contradiction. We used $w\leq0$ and the integrability of $f^{1/n}$. This bounds $a$ above.

For the last assertion the same estimate contains $c_j^{1/n}$ and $\mathrm{e}^{v_j/(jn)}$. If those integrals approached zero along a subsequence with $j\to\infty$, normalized compactness and almost-everywhere convergence would give $v_j/j\to0$ almost everywhere. Dominated convergence then gives the positive limit $\int_Xf^{1/n}G^{1-1/n}\,\mathrm{d}V$. This contradiction bounds $c_j$; finitely many remaining indices do not matter. No total top-degree mass was used in either argument.
\end{proof}


\section{Stability above a strict barrier}
\label{sec:big-reference-stability}

The subsolutions of \cref{lem:big-reference-subsolution} stay within a bounded distance of a given reference. With them, the proof of \cref{thm:bounded-one-sided-stability} extends to unbounded potentials lying above $\rho$, and the resulting estimate is uniform over all references between $\rho$ and zero. For a fixed reference, it upgrades $L^1$ compactness of solutions to uniform convergence of their differences on $\Omega$.

\begin{theorem}\label{thm:big-reference-one-sided-stability}\index{Stability!varying references}
Fix $C_0,L>0$, $p>1$, and
\[
 0<\gamma<\frac{p-1}{np+p-1}.
\]
Let $\eta\in\PSH(X,\theta)$ satisfy $\rho\leq\eta\leq0$, and suppose that $u\in\PSH(X,\theta)$ satisfies
\[
 \eta-C_0\leq u\leq\eta,\qquad
 \mathrm{MA}_\theta(u)=f\,\mathrm{d}V,\qquad 0\leq f\in L^p(\mathrm{d}V),
 \quad\|f\|_p\leq L.
\]
There is a constant $C$, depending only on the fixed geometry and on $\rho,C_0,L,p,\gamma$, such that every $v\in\PSH(X,\theta)$ with $v\leq\eta$ satisfies
\begin{equation}\label{eq:big-reference-stability}
 \sup_\Omega(v-u)_+
 \leq C\left(\int_\Omega(v-u)_+\,\mathrm{d}V\right)^\gamma.
\end{equation}
\end{theorem}

The comparison potential $v$ need not have the ordinary singularity type of $\eta$. Since $Z$ is pluripolar and has zero volume, the integral may equally be taken over $X$, and the resulting inequality $v\leq u+C\bigl(\int_\Omega(v-u)_+\,\mathrm{d}V\bigr)^\gamma$ on $\Omega$ extends to $X$.

\begin{proof}
Replacing $v$ by $\max(v,u)$, we may assume $\eta-C_0\leq u\leq v\leq\eta$, so that $0\leq v-u\leq C_0$ on $\Omega$. We repeat the proof of \cref{thm:bounded-one-sided-stability}, with $1<q<p$, $E_\delta=\{u<v-\delta\}$, $g_\delta=\mathbf1_{E_\delta}f$ and $m_\delta=\|g_\delta\|_q$, and with suprema taken over $\Omega$. Three ingredients change. First, when $m_\delta>0$, \cref{lem:big-reference-subsolution}, applied to $g_\delta/m_\delta$ and followed by subtracting one from its potential, gives $h\in\PSH(X,\theta)$ and a fixed $b>0$ such that
\[
 \eta-1\leq h\leq\eta,\qquad
 \mathrm{MA}_\theta(h)\geq b\,g_\delta/m_\delta\,\mathrm{d}V;
\]
the constant $b$ is uniform in $\eta$, because the auxiliary density has $L^q$ norm one and the strict barrier is fixed. These bounds give $h-v\leq C_0$ and $v-h\leq1$, the two properties of $h$ used there. Second, all potentials involved are bounded below by $\rho$ minus a constant, so positivity of mixed products holds on $\Omega$, where they are locally bounded, and $Z$ is not charged. Third, bounded domination is replaced by \cref{thm:big-barrier-domination}: the potentials $v-\delta$ and $w=(1-t)v+th-\delta-C_0t$ lie below $u+C_0$, so this theorem applies with the coefficients zero and $1/2$. The rest of the argument is unchanged. It gives \eqref{eq:big-reference-stability} with the exponent $\gamma_q=(p-q)/(npq+p-q)$, and we choose $q>1$ so close to one that $\gamma_q>\gamma$.
\end{proof}

\Needspace{15\baselineskip}
\begin{corollary}\label{cor:big-reference-uniform-convergence}
Fix $\eta\in\PSH(X,\theta)$ with $\rho\leq\eta\leq0$, constants $C_0,L>0$, and $p>1$. Consider the family of potentials $u\in\PSH(X,\theta)$ satisfying
\[
 \eta-C_0\leq u\leq\eta,\qquad
 \mathrm{MA}_\theta(u)=f_u\,\mathrm{d}V,\qquad0\leq f_u\in L^p(\mathrm{d}V),
 \quad\|f_u\|_p\leq L.
\]
This family is compact for uniform convergence of differences on $\Omega$. If $u_j$ belong to the family and $u_j\to u$ in $L^1$, then $u$ belongs to the family and
\[
 \sup_\Omega|u_j-u|\leq C\|u_j-u\|_{L^1(\mathrm{d}V)}^\gamma
 \longrightarrow0
\]
for every exponent $\gamma$ as in \cref{thm:big-reference-one-sided-stability}. Moreover, $f_{u_j}$ converges weakly to $f_u$ in $L^p$.
\end{corollary}
\begin{proof}
Applying the theorem in both orders to $u_j$ and $u_k$, with the common reference $\eta$, we see that an $L^1$-Cauchy sequence in the family is uniformly Cauchy in differences on $\Omega$. Its uniform limit on $\Omega$ is locally $\theta$-plurisubharmonic and agrees almost everywhere, hence everywhere for the canonical representatives, with the $L^1$ limit $u$. The bounds $\eta-C_0\leq u\leq\eta$ hold on all of $X$. Letting $k\to\infty$ in the pairwise estimate proves the quantitative bound, before any equation for $u$ is used.

The densities form a bounded set in the reflexive space $L^p$. Let $f$ be the weak limit of a subsequence $f_{u_{j_k}}$. Since the potentials are locally bounded and converge locally uniformly on $\Omega$, Bedford--Taylor continuity, which holds also for the smooth nonclosed background, gives $\mathrm{MA}_\theta(u)=f\,\mathrm{d}V$ on $\Omega$. Neither measure charges the analytic set $Z$, so this equality holds on $X$. Positivity and weak lower semicontinuity of the norm give $f\geq0$ and $\|f\|_p\leq L$. Thus every weak cluster point equals the density $f_u$ of $\mathrm{MA}_\theta(u)$, which proves weak convergence of the whole sequence. Finally, the common bounds by $\eta$ and $\eta-C_0$ give $L^1$ compactness of the family, by the compactness of quasi-plurisubharmonic functions, and the preceding argument upgrades each convergent subsequence to uniform convergence of differences on $\Omega$.
\end{proof}

For solutions with minimal singularities, normalized by $\sup_Xu=0$, we take $\eta=V_\theta$. Then $\rho\leq V_\theta\leq0$, and every normalized comparison potential $v\in\PSH(X,\theta)$ satisfies $v\leq V_\theta$, as the theorem requires.


\section{Minimal singularities in an arbitrary big background}
\label{sec:big-minimal-existence}

We retain the assumptions and the strict potential of \cref{sec:big-minimal-comparison}. The reference is now
\[
 V_\theta=\left(\sup\{w\in\PSH(X,\theta):w\leq0\}\right)^*.
\]
In particular $\rho\leq V_\theta\leq0$ and $\sup_XV_\theta=0$. We will solve the equations in the ordinary singularity type of $V_\theta$, with no bounded mass or uniform positive lower volume assumption. The resolution step uses the following precise version of classical regularization.

\begin{theorem}[Demailly regularization]\label{thm:resolvable-regularization-input}
Let $\alpha$ be any smooth real $(1,1)$-form and let $u\in\PSH(X,\alpha)$. There are $\delta_j\downarrow0$ and $\widetilde u_j\in\PSH(X,\alpha+\delta_j\omega)$ decreasing to $u$ such that each $\widetilde u_j$ is smooth off an analytic pole set and, after a finite sequence of blowups with smooth centers, has local expressions
\[
 \pi_j^*\widetilde u_j=\sum_\ell a_\ell\log|s_\ell|+b,
 \qquad a_\ell>0,\quad b\in C^\infty.
\]
The divisors may be taken to have simple normal crossings. The coefficients include the fixed normalization of $\ddc$.
\end{theorem}

This is the standard input \cite[Theorem 3.2]{DP04}, applied to the closed current $\ddc u$ with smooth lower bound $-\alpha$. Its lower bound is not required to be closed, and its ambient manifold is only Hermitian. Item (iv) supplies the smooth remainder after resolution; ordinary analytic type with merely a bounded remainder would not suffice here. This version need not be the particular multiplier-ideal sequence in \cref{thm:demailly-hermitian-input}. We do not assert that the additional conclusions of these two regularization statements hold for the same sequence. Applying this theorem to any strictly $\theta$-plurisubharmonic potential with a sufficiently small curvature error also justifies the choice of $\rho$ above.

\begin{lemma}\label{lem:big-resolved-equation}
Let $\psi\in\PSH(X,\theta)$ satisfy $\theta_\psi\geq a\omega$ for some $a>0$ and have the resolved smooth-remainder form in the preceding theorem. For every nonzero $0\leq f\in L^p(\mathrm{d}V)$, $p>1$, and $\lambda>0$, there exists $u\in\PSH(X,\theta)$ with $u\simeq\psi$ such that
\[
 \mu_u=\mathrm{e}^{\lambda u}f\,\mathrm{d}V.
\]
\end{lemma}
\begin{proof}
Resolve $\psi$ by $\pi:Y\to X$. The pulled-back current has the form
\[
 \pi^*\theta+\ddc\pi^*\psi=\beta+[D],
\]
where $D$ is effective and $\beta$ is smooth. Off $D$ one has $\beta\geq a\pi^*\omega$, so this inequality holds everywhere by smoothness. Thus $\beta$ is semipositive.

We check its bigness. For a blowup with smooth center and exceptional divisor $E$, choose a metric on $\mathcal O(-E)$ with curvature positive in the directions tangent to the projective fibers. The pullback of a Hermitian metric is positive in the complementary directions along $E$. Consequently its sum with a sufficiently small positive multiple of this curvature is a Hermitian metric: first check positivity along $E$, then in a neighborhood, and finally on the compact complement, where the pullback is already positive. Iterating over the blowups gives an effective real exceptional divisor $F$, with logarithmic potential $\zeta$, and a smooth Hermitian form $\kappa$ such that
\[
 \pi^*\omega+\ddc\zeta=\kappa+[F].
\]
For the iteration, pull back the previous smooth positive form, apply the one-blowup construction, and add the pulled-back previous divisor to the new positive multiple of $E$. Hence
$\beta+a\ddc\zeta\geq a\kappa$, proving bigness of $\beta$. This argument concerns Hermitian metrics and never requires a closed positive form on $Y$.

Fix a smooth positive volume form $\mathrm{d}V_Y$, and write $\pi^*\mathrm{d}V=J\,\mathrm{d}V_Y$. The function $J$ is smooth, nonnegative and bounded, and is positive off the exceptional locus. Define
\[
 H=\mathrm{e}^{\lambda\pi^*\psi}(f\circ\pi)J.
\]
It is nonzero and belongs to $L^p(\mathrm{d}V_Y)$. Indeed $\psi$ is bounded above, and change of variables gives
\[
 \int_Y|(f\circ\pi)J|^p\,\mathrm{d}V_Y
 \leq\|J\|_\infty^{p-1}\int_X|f|^p\,\mathrm{d}V.
\]
Values on the analytic exceptional sets do not affect these integrals. The semipositive big-background theorem \cref{thm:big-bounded-exponential} produces a bounded $s\in\PSH(Y,\beta)$ with
\[
 \mathrm{MA}_\beta(s)=\mathrm{e}^{\lambda s}H\,\mathrm{d}V_Y.
\]
The function $\pi^*\psi+s$ is $\pi^*\theta$-plurisubharmonic, because its current is $\beta_s+[D]\geq0$. By \cref{thm:hermitian-space-descent} it descends to $u\in\PSH(X,\theta)$. Boundedness of $s$ gives $u\simeq\psi$. Off the exceptional and polar sets, change of variables proves the required equation. Both its extended left-hand side and its smooth-volume right-hand side ignore those analytic sets. The equation therefore holds on all of $X$.
\end{proof}

\begin{theorem}\label{thm:big-minimal-exponential}
Let $\theta$ be any smooth big real $(1,1)$-form and let $0\leq f\in L^p(\mathrm{d}V)$, $p>1$, be nonzero. For every $\lambda>0$ there is a unique $u\in\PSH(X,\theta)$ such that
\[
 u\simeq V_\theta,\qquad
 \mathrm{MA}_\theta(u)=\mathrm{e}^{\lambda u}f\,\mathrm{d}V.
\]
No bounded mass or uniform positive lower volume hypothesis is needed.
\end{theorem}
\begin{proof}
Apply \cref{thm:resolvable-regularization-input} to $V_\theta$. Set $\delta_j=2^{-j}$ and take a subsequence whose curvature errors are at most $\delta_j$. Denote its potentials by $\widetilde\psi_j\downarrow V_\theta$; then
$\theta_{\widetilde\psi_j}\geq-\delta_j\omega$. Set
\[
 a_j=(1-\delta_j)\sup_X\widetilde\psi_j\geq0,
 \qquad
 \psi_j=(1-\delta_j)\widetilde\psi_j+\delta_j\rho-a_j.
\]
Normalized compactness, or monotonicity and Hartogs' lemma, gives $a_j\to0$. Resolving the singularities of $\widetilde\psi_j$ and $\rho$ together shows that $\psi_j$ has a smooth remainder after resolution. Moreover,
\begin{equation}\label{eq:big-minimal-approximation}
 \theta_{\psi_j}\geq\delta_j^2\omega,
 \qquad
 (1-\delta_j)V_\theta+\delta_j\rho-a_j
 \leq\psi_j\leq V_\theta.
\end{equation}
The upper inequality follows because $\psi_j$ is $\theta$-plurisubharmonic and at most zero. The explicit subtraction of $a_j$ keeps the normalization separate from the fixed pole barrier.

By the resolved equation lemma there are $u_j\simeq\psi_j$ satisfying
\[
 \mu_{u_j}=\mathrm{e}^{\lambda u_j}f\,\mathrm{d}V.
\]
Each $u_j$ is bounded below by $\rho$ up to a constant depending on $j$. The Gauduchon trace proposition gives a uniform upper bound for $\sup_Xu_j$.

There is also a uniform lower estimate. Translate the fixed $\rho$ down by $1+\sup_j a_j$ when applying \cref{lem:big-reference-subsolution}, and take its reference to be $\psi_j-1$. The lemma gives $\psi_j-1\leq v_j\leq\psi_j\leq0$ and $\mu_{v_j}\geq mf\,\mathrm{d}V$, with $m>0$ independent of $j$. Put $w_j=v_j+\lambda^{-1}\log m$. Then
\[
 \mu_{w_j}\geq \mathrm{e}^{\lambda w_j}f\,\mathrm{d}V,
 \qquad u_j\simeq w_j\simeq\psi_j.
\]
The exponential order in \cref{cor:big-barrier-orders} gives $u_j\geq w_j$. Thus for a uniform $C$,
\begin{equation}\label{eq:big-minimal-solution-bounds}
 \psi_j-C\leq u_j\leq V_\theta+C.
\end{equation}
After extraction, $u_j$ converges in $L^1$ and almost everywhere to a genuine $u$. The preceding bounds give $V_\theta-C\leq u\leq V_\theta+C$.

\Needspace{14\baselineskip}
We verify its equation using envelopes. Put $A_j=\sup_{k\geq j}a_k$ and, for $N\geq j$, define
\[
 p_{j,N}=P_\theta(u_j,\ldots,u_N),\qquad
 t_{j,N}=\max_{j\leq k\leq N}u_k.
\]
Both functions lie above the genuine competitor
\[
 b_j=(1-\delta_j)V_\theta+\delta_j\rho-C-A_j.
\]
In particular every envelope is genuine. All these functions are uniformly locally bounded on $\Omega$, also as $j$ varies. Contact support from \cref{thm:big-general-envelope-contact} and a disjoint partition by the first contact index give
\[
 \mu_{p_{j,N}}\leq \mathrm{e}^{\lambda p_{j,N}}f\,\mathrm{d}V.
\]
There is no factor $N$ here: on each piece the envelope equals one $u_k$, and ordered contact bounds its measure by that single measure. Locality and ordered contact for a finite maximum similarly give
\[
 \mu_{t_{j,N}}\geq \mathrm{e}^{\lambda t_{j,N}}f\,\mathrm{d}V.
\]
Let $N\to\infty$, setting $p_j=\lim_Np_{j,N}$ and $t_j=(\lim_Nt_{j,N})^*$. Local bounded monotone continuity passes both inequalities to these limits. Next $p_j$ increases to a regularized limit $p$, while $t_j$ decreases to $u$ by Hartogs' lemma. The bound by $b_j$ gives $p\simeq V_\theta$, and $p\leq u$. A second use of local bounded monotone continuity gives
\[
 \mu_p\leq \mathrm{e}^{\lambda p}f\,\mathrm{d}V,
 \qquad \mu_u\geq \mathrm{e}^{\lambda u}f\,\mathrm{d}V.
\]
These inequalities hold on $\Omega$ and extend across $Z$ because all measures involved ignore it. Thus
$\mu_p\leq \mathrm{e}^{\lambda(p-u)}\mu_u$. The exponential order gives $p\geq u$, so $p=u$ and the equation follows. Applying that order in both directions also proves uniqueness. Every passage of the Monge--Amp\`ere operator above was along a locally bounded monotone sequence on $\Omega$.
\end{proof}

\begin{theorem}\label{thm:big-minimal-normalized}
Under the assumptions of \cref{thm:big-minimal-exponential}, there are $v\in\PSH(X,\theta)$ and a uniquely determined $c>0$ such that
\[
 \sup_Xv=0,
 \qquad V_\theta-C\leq v\leq V_\theta,
 \qquad \mathrm{MA}_\theta(v)=cf\,\mathrm{d}V.
\]
Uniqueness of the normalized potential is not asserted.
\end{theorem}
\begin{proof}
Solve the positive-exponent equations with $\lambda=1/j$, and write
\[
 v_j=u_j-\sup_Xu_j,\qquad
 c_j=\mathrm{e}^{\sup_Xu_j/j},\qquad
 \mu_{v_j}=c_j \mathrm{e}^{v_j/j}f\,\mathrm{d}V.
\]
Then $v_j\simeq V_\theta$ and $v_j\leq V_\theta\leq0$. The Gauduchon trace gives $c_j\leq C_0$ uniformly.

\Needspace{12\baselineskip}
Choose $\varepsilon>0$ small and $1<q<p$ such that $\mathrm{e}^{-\varepsilon v_j}f$ is bounded in $L^q$. This follows from H\"older's inequality and uniform Skoda integrability of normalized potentials in a fixed enlarged Hermitian background. The subsolution lemma, with reference $V_\theta-1$ and the fixed barrier $\rho-1$, produces
\[
 V_\theta-1\leq\eta_j\leq V_\theta,
 \qquad
 \mu_{\eta_j}\geq m \mathrm{e}^{-\varepsilon v_j}f\,\mathrm{d}V
                 \geq mf\,\mathrm{d}V
\]
for a uniform $m>0$. Since $\mu_{v_j}\leq c_jf\,\mathrm{d}V$, scalar order gives $c_j\geq m$.

\Needspace{8\baselineskip}
Take $D=\varepsilon^{-1}\max\{0,\log(C_0/m)\}$ and put $w_j=\eta_j-D$. We have
\[
 \mathrm{e}^{\varepsilon(w_j-v_j)}\mu_{v_j}
 =c_j \mathrm{e}^{v_j/j+\varepsilon\eta_j-\varepsilon D-\varepsilon v_j}
           f\,\mathrm{d}V
 \leq m \mathrm{e}^{-\varepsilon v_j}f\,\mathrm{d}V
 \leq\mu_{w_j}.
\]
Exponential order therefore gives $v_j\geq w_j$. In particular
\begin{equation}\label{eq:big-normalized-uniform-type}
 V_\theta-1-D\leq v_j\leq V_\theta,
 \qquad m\leq c_j\leq C_0.
\end{equation}
Extract a subsequence with $v_j\to v$ in $L^1$ and almost everywhere and $c_j\to c\in[m,C_0]$. The limit is normalized, genuine and of minimal ordinary type. After relabeling, denote the retained exponents by $\lambda_j\downarrow0$; thus $\mu_{v_j}=c_j\mathrm{e}^{\lambda_jv_j}f\,\mathrm{d}V$.

One more comparison is needed to identify the measure. Define finite lower envelopes $p_{j,N}=P_\theta(v_j,\ldots,v_N)$ and finite maxima $t_{j,N}=\max_{j\leq k\leq N}v_k$. They all lie between $V_\theta-1-D$ and $V_\theta$. On $\Omega$ put
\[
 h_k=c_k \mathrm{e}^{(\lambda_k-\varepsilon)v_k},\qquad
 H_j=\sup_{k\geq j}h_k,\qquad L_j=\inf_{k\geq j}h_k.
\]
These functions are bounded on each relatively compact subset of $\Omega$, uniformly in $j$. Almost everywhere they converge, as $j\to\infty$, to $c \mathrm{e}^{-\varepsilon v}$. Contact partition and the maximum principle give
\[
 \mu_{p_{j,N}}\leq \mathrm{e}^{\varepsilon p_{j,N}}H_jf\,\mathrm{d}V,
 \qquad
 \mu_{t_{j,N}}\geq \mathrm{e}^{\varepsilon t_{j,N}}L_jf\,\mathrm{d}V.
\]
First let $N\to\infty$, obtaining $p_j$ and $t_j$, and then let $j\to\infty$. The lower envelopes increase to $p\simeq V_\theta$ with $p\leq v$, while the upper envelopes decrease to $v$. Local bounded monotone continuity applies to the left sides, and dominated convergence on relatively compact subsets of $\Omega$ applies to the densities. We obtain
\[
 \mu_p\leq c \mathrm{e}^{\varepsilon(p-v)}f\,\mathrm{d}V,
 \qquad \mu_v\geq cf\,\mathrm{d}V.
\]
Extend across $Z$ and apply exponential order to get $p\geq v$. Hence $p=v$ and $\mu_v=cf\,\mathrm{d}V$. The artificial fixed exponent $\varepsilon$ in this last argument is what allows domination to identify the two envelope limits; the unweighted inequalities alone would not do so. Scalar uniqueness is \cref{cor:big-barrier-orders}.
\end{proof}

These theorems reconstruct \cite[Theorems 4.6--4.7]{BGL25}. Their equations prove finite mass for the solutions, not a mass bound for their whole ordinary singularity type. At exponent zero only the scalar is unique; \cref{thm:degenerate-hermitian-nonuniqueness} gives counterexamples to uniqueness of the potential.


\Needspace{11\baselineskip}
\section{Local H\"older continuity in a big background}
\label{sec:big-minimal-holder}

We now combine the stability estimate above $\rho$ with the regularization of \cref{lem:smooth-background-regularization-rate}, whose curvature error is linear in the regularization parameter. As in the proof of \cref{thm:hermitian-holder-elementary}, a small multiple of the strict barrier absorbs the curvature error of the transforms, and a lower bound for the minimizing radius converts the estimate for the transforms into an estimate for the smoothing on compact subsets of $\Omega$. This gives a local H\"older modulus for solutions with minimal singularities.

\begin{theorem}\label{thm:big-minimal-holder}\index{Holder continuity@H\"older continuity!minimal singularities}
Fix $C_0>0$. Let $u\in\PSH(X,\theta)$ satisfy
\[
 \sup_Xu=0,\qquad V_\theta-C_0\leq u\leq V_\theta,
 \qquad\mathrm{MA}_\theta(u)=f\,\mathrm{d}V,
\]
where $0\leq f\in L^p(\mathrm{d}V)$ and $p>1$.
Then $u$ is locally H\"older continuous on $\Omega$, with every exponent
\[
 0<\beta<\frac{2(p-1)}{np+3(p-1)}.
\]
For fixed bounds on $C_0$ and $\|f\|_p$, the H\"older constants are uniform on
each fixed compact subset of $\Omega$.
\end{theorem}

\begin{proof}
Since $u\geq V_\theta-C_0\geq\rho-C_0$, we may subtract a constant from the strict barrier so that $\rho\leq u-1$; it still satisfies $\theta+\ddc\rho\geq\kappa\omega$ for a fixed
$\kappa>0$. Choose
$0<\gamma<(p-1)/(np+p-1)$ and set $\alpha=2\gamma<1$.
Apply \cref{lem:smooth-background-regularization-rate} with $a=u$, and let $A,K,C$ and $\Phi_{c,\delta}$ be the constants and transforms given there. For small $\delta>0$ set
\[
 c=\delta^\alpha,\qquad e=Ac+C\delta,\qquad s=e/\kappa,
 \qquad v_\delta=\frac{\Phi_{c,\delta}+s\rho}{1+s}.
\]
By the curvature bound of \cref{lem:smooth-background-regularization-rate}, $v_\delta\in\PSH(X,\theta)$. The lower bound for $\Phi_{c,\delta}$ in that lemma implies $\rho\leq\Phi_{c,\delta}$ when $\delta$
is small. Thus
\[
 v_\delta\leq\Phi_{c,\delta}\leq R_\delta u\leq0.
\]
In particular $v_\delta\leq V_\theta$, so \cref{thm:big-reference-one-sided-stability} applies with the reference
$V_\theta$ and the comparison potential $v_\delta$. Together with the $L^1$ smoothing estimate, it gives
\[
 \sup_\Omega(v_\delta-u)_+
 \leq C\|(R_\delta u-u)_+\|_1^\gamma\leq C\delta^\alpha.
\]
Fix a compact subset of $\Omega$ on which $\rho\geq-M$. Since $s\leq C\delta^\alpha$
and $u\leq0$, the identity
\[
 \Phi_{c,\delta}-u=(1+s)(v_\delta-u)+s(u-\rho)
\]
gives $\Phi_{c,\delta}-u\leq C_M\delta^\alpha$ there.

At a radius $t$ where the infimum defining $\Phi_{c,\delta}$ is attained, the monotone lower bound yields
\[
 -c\log(t/\delta)\leq\Phi_{c,\delta}-u+K(\delta^2+\delta)
 \leq C_Mc.
\]
Hence $t\geq a_M\delta$ for a fixed $a_M>0$. As in the proof of
\cref{thm:hermitian-holder-elementary}, monotonicity of the smoothing now gives
\[
 R_{a_M\delta}u-u\leq C_M\delta^\alpha,
 \qquad R_{a_M\delta}u-u\geq-K(a_M\delta)^2
\]
on the chosen compact set. Since every small radius has the form $a_M\delta$, we obtain $|R_tu-u|\leq C_Mt^\alpha$ there for all small $t>0$, with constants uniform on that compact set.
\Needspace{8\baselineskip}
We apply this to a slightly larger compact neighborhood, on which also
$\|\nabla R_tu\|\leq C_M/t$. For nearby points $x,y$ at distance $d>0$, the smooth
function $R_tu$ and a short path in that neighborhood give
\[
 |u(x)-u(y)|\leq C_M(t^\alpha+d/t).
\]
The choice $t=d^{1/(1+\alpha)}$ gives the exponent
$\alpha/(1+\alpha)$, and letting $\gamma$ approach its upper endpoint gives
the asserted range. Pairs of points farther apart are handled by the boundedness of $u$ on the compact neighborhood.
\end{proof}

It suffices that the equation hold on $\Omega$. Suppose that
$\mathrm{MA}_\theta(u)=cf\,\mathrm{d}V$ on $\Omega$, where $c>0$ and
$0\leq f\in L^p(X,\mathrm{d}V)$, and that $u$ satisfies the same bounds as in the theorem.
The lower bound $u\geq V_\theta-C_0\geq\rho-C_0$ shows that $u$ is locally
bounded on $\Omega$, so there $\mathrm{MA}_\theta(u)$ is the local bounded product.
Since $Z$ has zero volume and is not charged by $\mathrm{MA}_\theta(u)$, the equation holds on $X$, and the theorem applies with
$cf$ in place of $f$. The uniform estimates then depend on a bound for
$\|cf\|_p$, which includes the constant $c$.

\subsection{Limits of normalized exponential solutions}

For the normalized exponential family $(v_\lambda,c_\lambda)$ of \cref{lem:big-exponential-uniform-type}, the reference interval $[V_\theta-B,V_\theta]$ is uniform in $0<\lambda\leq1$. Stability therefore controls differences on all of $\Omega$, and the local H\"older estimates strengthen the convergence on compact subsets. We fix the nonzero density $0\leq f\in L^p(\mathrm{d}V)$, $p>1$, and use the constants $B,C,m$ of that lemma.

\Needspace{16\baselineskip}
\begin{proposition}\label{prop:big-exponential-clusters}
For the family $(v_\lambda,c_\lambda)$ of
\cref{lem:big-exponential-uniform-type}, the normalized potentials
are relatively compact for uniform convergence of
differences on $\Omega$. As $\lambda\downarrow0$, every cluster limit
$v$ is normalized, has minimal singularities, and satisfies
\[
 \mathrm{MA}_\theta(v)=cf\,\mathrm{d}V,
\]
where $c>0$ is the unique constant of
\cref{thm:big-minimal-normalized}. Moreover,
\[
 c_\lambda\longrightarrow c,
 \qquad c_\lambda\mathrm{e}^{\lambda v_\lambda}f
          \longrightarrow cf\quad\text{in }L^p(\mathrm{d}V).
\]
Every sequence of parameters in $(0,1]$ has a subsequence for which the
potentials converge in $C^\beta$ on each compact subset of $\Omega$, for
all
\[
 0<\beta<\frac{2(p-1)}{np+3(p-1)}.
\]
\end{proposition}

\begin{proof}
By \cref{lem:big-exponential-uniform-type}, all $v_\lambda$ lie in a common reference interval, and their densities are bounded in $L^p$ by $C\|f\|_p$. We may therefore apply
\cref{cor:big-reference-uniform-convergence} with $\eta=V_\theta$.
Given any sequence $\lambda_j\downarrow0$, we pass to a subsequence with
$v_{\lambda_j}\to v$ uniformly in differences on $\Omega$ and in
$L^1(X)$, and with $c_{\lambda_j}\to c_*\in[m,C]$.
The limit satisfies the same reference bounds and is normalized. Off the pole set,
\[
 0\leq1-\mathrm{e}^{\lambda v_\lambda}
 \leq1-\mathrm{e}^{\lambda(V_\theta-B)}\longrightarrow0
\]
as $\lambda\downarrow0$. Multiplying by $f$ and using dominated convergence, we obtain
$\mathrm{e}^{\lambda v_\lambda}f\to f$ in $L^p$ for the whole family.
The compactness corollary therefore identifies the limit equation as
$\mathrm{MA}_\theta(v)=c_*f\,\mathrm{d}V$, and the uniqueness of the constant gives $c_*=c$.
Since every subsequential limit of $c_\lambda$ as $\lambda\downarrow0$ equals $c$, we conclude that
$c_\lambda\to c$, and the stated convergence of the densities follows.

Finally, \cref{thm:big-minimal-holder} gives uniform local $C^\alpha$
bounds for every $\alpha$ below its endpoint, and the uniform limit
inherits these bounds by pointwise convergence. For
$0<\beta<\alpha$, uniform convergence and the elementary interpolation
inequality on smaller coordinate balls give
\[
 [v_{\lambda_j}-v]_{C^\beta}
 \leq C\|v_{\lambda_j}-v\|_\infty^{1-\beta/\alpha}
          [v_{\lambda_j}-v]_{C^\alpha}^{\beta/\alpha}
 \longrightarrow0,
\]
where the H\"older bounds are taken on slightly larger compact neighborhoods. The
uniformly convergent subsequence given by compactness thus converges in $C^\beta$
on every compact subset of $\Omega$ and for every exponent below the endpoint; no further
extraction depending on the exponent is needed.
\end{proof}

\begin{corollary}\label{cor:big-exponential-symmetries}
Suppose a group $G$ of biholomorphisms of $X$ preserves $\theta$ and the measure
$f\,\mathrm{d}V$. Then every $v_\lambda$ and every cluster limit as
$\lambda\downarrow0$ is $G$-invariant. If the normalized zero-exponent
equation has only one $G$-invariant solution with minimal singularities,
then the whole family $v_\lambda$ converges to it as $\lambda\downarrow0$, uniformly in
differences on $\Omega$ and in each local $C^\beta$ topology of the
proposition.
\end{corollary}
\begin{proof}
For $g\in G$, pullback by $g$ preserves $\theta$-plurisubharmonicity, and
$V_\theta\circ g=V_\theta$ by the definition of the envelope. Pullback also
preserves the positive-exponent equation, so its uniqueness gives
$u_\lambda\circ g=u_\lambda$, and the normalization preserves this equality.
Pullback by a fixed smooth biholomorphism is continuous on $L^1(X)$,
since its Jacobian relative to the fixed smooth volume is bounded.
Hence every cluster limit is invariant almost everywhere, and therefore everywhere for the
canonical representatives. This argument does not require $g$ to
preserve the barrier $\rho$ or the set $\Omega$.

If the invariant solution is unique, every cluster
limit is the same potential. Relative compactness then implies
convergence of the whole family for uniform convergence of differences on $\Omega$:
a sequence staying at a fixed positive distance from that potential would
have a convergent subsequence, a contradiction. The interpolation
argument above upgrades this to the asserted local $C^\beta$ convergence.
\end{proof}

The proposition gives compactness and identifies the equation of every cluster limit, and the symmetry criterion is one way to identify the limit itself. In general, \cref{thm:vanishing-exponential-selection} shows that the whole family converges, and it describes the limit through the greatest subsolution of the zero-exponent equation that satisfies a family of integral normalizations determined by how the constant $c$ depends on the density. This holds even when the zero-exponent equation has several normalized solutions.


\Needspace{12\baselineskip}
% Book adaptation of the supplied nonclosed regularity manuscript.
% Review scope: research/editorial/nonclosed-interior-regularity-review.md.
\section{Interior smoothness on the regular locus}
\label{sec:nonclosed-interior-smoothness}

Throughout this section, $\theta$ is a smooth big real $(1,1)$-form, and $\rho$, $Z=\{\rho=-\infty\}$ and $\Omega=X\setminus Z$ are the strict analytic barrier and its regular locus fixed in \cref{sec:big-minimal-comparison}. Thus $\rho\leq0$ has analytic singularities with smooth remainder, and $\theta+\ddc\rho\geq\omega$. If $\varphi\in\PSH(X,\theta)$ has minimal singularities, then $\varphi\geq V_\theta-C\geq\rho-C$ for a constant $C$. Hence $\varphi$ is locally bounded on $\Omega$, and its Monge--Amp\`ere measure there is the bounded-potential product.

For such a potential, \cref{thm:big-minimal-holder} gives local H\"older continuity on $\Omega$ when the density belongs to $L^p$ for some $p>1$. We show that a smooth strictly positive density on the whole of $\Omega$ then makes the potential smooth on $\Omega$. Local information alone cannot give this: \cref{ex:local-lipschitz-constant-measure} below is a locally Lipschitz plurisubharmonic function with constant Monge--Amp\`ere density that is not $C^1$. The argument therefore uses the barrier globally. By the proper gap of \cref{lem:interior-proper-gap}, the difference between the solution and the barrier attains a minimum on the closed set where smoothness fails, and strict positivity of $\theta+\ddc\rho$ turns this minimum into a quadratic separation. Near the minimum, smooth Dirichlet approximants have traces of polynomial growth; the separation improves this growth until a uniform trace bound yields elliptic regularity.

\Needspace{12\baselineskip}
\begin{theorem}\label{thm:nonclosed-interior-smoothness}\index{Interior regularity}
Suppose $\varphi\in\PSH(X,\theta)$ has minimal singularities and is locally
H\"older continuous on $\Omega$. Let $c>0$ and let $f$ be smooth and strictly positive on $\Omega$. Assume
\[
 (\theta+\ddc\varphi)^n=cf\,\omega^n\quad\text{on }\Omega.
\]
Then
$\varphi\in C^\infty(\Omega)$ and $\theta+\ddc\varphi>0$ there.
\end{theorem}

The H\"older hypothesis holds when, in addition, $f\in L^p(X,\omega^n)$ for some $p>1$, by \cref{thm:big-minimal-holder}; see \cref{cor:big-lp-interior-smoothness}. The theorem also applies when local H\"older continuity is obtained by another argument.

The proof is by contradiction and occupies the rest of the section. Let $\mathcal S\subset\Omega$ be the closed set of points near which $\varphi$ is not smooth, suppose that $\mathcal S\neq\varnothing$, and let $p\in\mathcal S$ minimize $\varphi-\rho$ on $\mathcal S$. On a small ball centered at $p$ the argument has four steps.
\begin{enumerate}
\item A comparison principle for continuous solutions (\cref{sec:interior-calculus}) and the smooth Dirichlet problem give smooth solutions $u_j$ of the equation that converge uniformly to $\varphi$ at a polynomial rate in $j$, while the traces of $\theta+\ddc u_j$, measured in Euclidean coordinates, grow at most polynomially in $j$ (\cref{sec:interior-approximants}).
\item Near $p$, the quadratic separation along $\mathcal S$, together with uniform trace bounds on compact sets where $\varphi$ is already smooth, improves the trace growth to $O(j^\gamma)$ for every $\gamma>0$. Interpolation with the convergence rate then gives uniform $C^{1,\alpha}$ bounds for every $\alpha<1$ (\cref{sec:interior-localization}).
\item A third-order estimate whose constants depend polynomially on the ellipticity turns uniform $C^{1,\alpha}$ bounds, for $\alpha$ close to one, into a uniform trace bound (\cref{sec:interior-calabi}).
\item A continuous solution with bounded trace is smooth (\cref{sec:interior-criterion}). Hence $\varphi$ is smooth near $p$, contradicting $p\in\mathcal S$ (\cref{sec:interior-main-proof}).
\end{enumerate}

\subsection{Local comparison}\label{sec:interior-calculus}

The comparison principle below, for continuous sub- and supersolutions on a bounded domain, identifies the given continuous potential with limits of smooth Dirichlet solutions. Since the background need not be closed, the classical Bedford--Taylor comparison does not apply directly. We freeze the background at a point, apply the closed-current comparison formula \eqref{eq:interior-power-comparison} to quadratic modifications, and control the error by the mixed inequalities of \cref{lem:interior-mixed-background,lem:interior-convolution}.

The local arguments of this section take place in coordinate balls, always with compact closure in a larger coordinate chart. There we put $\beta=\ddc|z|^2$, write $\partial_i=\partial/\partial z_i$ and $\partial_{\bar i}=\partial/\partial\bar z_i$, and take coefficient matrices of real $(1,1)$-forms relative to $\beta$. For such a form $A$, $\operatorname{tr}_\beta A$ and $\det_\beta A$ denote the trace and the determinant of its coefficient matrix, so that $A^n=(\det_\beta A)\beta^n$, and $\Delta_\beta=\operatorname{tr}_\beta\ddc$. For a positive form $T$, $\operatorname{tr}_T$ is the trace computed with the inverse of the coefficient matrix of $T$. Indices run from $1$ to $n$, repeated indices in tensor expressions are summed, and $I$ denotes the identity matrix. In local statements, $\vartheta$ denotes a smooth real $(1,1)$-form, not necessarily closed, which plays the role of $\theta$.

\begin{proposition}\label{prop:interior-comparison}
Let $D\Subset\mathbb C^n$ be a bounded domain and let $\vartheta$ be a smooth real $(1,1)$-form near $\overline D$. Suppose $u,v\in C(\overline D)$ are $\vartheta$-plurisubharmonic,
\[
  (\vartheta+\ddc u)^n\geq F\beta^n,
  \qquad
  (\vartheta+\ddc v)^n\leq F\beta^n,
\]
where $F>0$ is continuous on $D$. If $u\leq v$ on $\partial D$, then $u\leq v$ in $D$.
\end{proposition}

\begin{proof}
\emph{Step 1: $v$ admits no strict smooth lower test.} Suppose, to the contrary, that a $C^2$ function $q$ touches $v$ from below at a point $p\in D$ and that, at $p$,
\[
  \vartheta+\ddc q>0,
  \qquad
  \det_\beta(\vartheta+\ddc q)>F.
\]
\Needspace{6\baselineskip}
Choose a sufficiently small $a>0$. After replacing $q$ by $q-a|z-p|^2$, which makes the contact strict, and restricting to a small ball $B_r(p)$, there are fixed $\lambda,\eta>0$ such that
\[
  A:=\vartheta+\ddc q\geq\lambda\beta,
  \qquad A^n\geq(F+\eta)\beta^n.
\]
To compare with closed currents, freeze the background at $p$: put $\kappa=\sup_{B_r}\lVert\vartheta-\vartheta(p)\rVert_\beta$ and $B_0=\vartheta(p)+\kappa\beta$. Then $S=B_0-\vartheta$ satisfies $0\leq S\leq2\kappa\beta$. For small $s>0$ the set $E=\{v<q+s\}$ is nonempty and relatively compact in $B_r$. If $Q_0$ is a quadratic potential for the constant form $B_0$, then
\[
  P=\ddc(v+Q_0)=\vartheta+\ddc v+S,
  \qquad
  Q=\ddc(q+s+Q_0)=A+S\geq\lambda\beta
\]
are closed positive Hessians. Applying \eqref{eq:interior-power-comparison}, expanding $P$, and using $S\leq2\kappa\beta\leq2\kappa\lambda^{-1}Q$, we obtain
\[
  \int_E P^n\leq\int_E F\beta^n+e_\kappa\int_E P^n,
  \qquad
  e_\kappa=\left(1+\frac{2\kappa}{\lambda}\right)^n-1.
\]
On the other hand, \eqref{eq:interior-power-comparison} with $j=n$ gives
\[
  \int_E P^n\geq\int_E Q^n\geq\int_E(F+\eta)\beta^n.
\]
If $M$ bounds $F$ on the ball, the two estimates give $(1-e_\kappa)\eta\leq e_\kappa M$. Since $\kappa\to0$, and hence $e_\kappa\to0$, as $r\downarrow0$, this is impossible for small $r$. Thus no such lower test exists.

\emph{Step 2: comparison.} Suppose that $\max_D(u-v)>0$. We regularize $u$ to produce a forbidden lower test for $v$. Let $\chi_\delta$ be a nonnegative smooth mollifier of mass one and radius $\delta>0$. For small $t>0$ and $C\geq\sup_D|z|^2$, set
\[
  u_t=u+t(|z|^2-C).
\]
Then $u_t\leq v$ on $\partial D$, the positive maximum persists, and $T_t=\vartheta+\ddc u_t\geq t\beta$. \Cref{lem:interior-mixed-background} gives
\[
  T_t^n\geq(F^{1/n}+t)^n\beta^n.
\]
Choose a relatively compact neighborhood $W$ of the positive maximum set. By \cref{lem:interior-convolution},
\[
  (T_t*\chi_\delta)^n
  \geq\bigl((F^{1/n}+t)*\chi_\delta\bigr)^n\beta^n.
\]
Since $T_t\geq t\beta$, the convolution satisfies $T_t*\chi_\delta\geq t\beta$. Moreover,
\[
  \vartheta+\ddc(u_t*\chi_\delta)
  =T_t*\chi_\delta+(\vartheta-\vartheta*\chi_\delta),
\]
so if
$\epsilon_\delta=\|\vartheta-\vartheta*\chi_\delta\|_{C^0(W)}$, then
\[
 \vartheta+\ddc(u_t*\chi_\delta)
 \geq(1-\epsilon_\delta/t)(T_t*\chi_\delta).
\]
For fixed $t>0$, uniform convergence of the convolved density and
$\epsilon_\delta\to0$ show that the left-hand side is positive and has
determinant strictly larger than $F$ on $W$ for small $\delta$. Since $u_t*\chi_\delta\to u_t$ uniformly on $W$, the difference $u_t*\chi_\delta-v$ still attains its maximum over $\overline W$ at an interior point of $W$ for small $\delta$. Subtracting this maximum from $u_t*\chi_\delta$ produces a smooth function touching $v$ from below with the properties excluded in Step~1, a contradiction.
\end{proof}

\subsection{Bounded traces and smoothness}\label{sec:interior-criterion}

We begin with the criterion used in the last step of the proof. A bounded trace of a positive current bounds all its coefficient measures, so a continuous solution with bounded trace is a strong solution with $W^{2,q}$ regularity for every finite $q$. The equation then bounds the eigenvalues from below, so it becomes uniformly elliptic. Since the solution is not yet known to be $C^2$, the Evans--Krylov theory is applied to it as a viscosity solution of a concave extension of the operator.

\begin{proposition}\label{prop:interior-criterion}
Let $B\subset\mathbb C^n$ be a coordinate ball, let $\vartheta$ be a smooth real $(1,1)$-form there and let $F\in C^\infty(B)$ be strictly positive. Suppose $u$ is continuous, $T=\vartheta+\ddc u\geq0$, and
\begin{equation}\label{eq:interior-local-ma}
  T^n=F\beta^n.
\end{equation}
For a point $p\in B$, the following conditions, each required on some neighborhood of $p$, are equivalent:
\begin{enumerate}
\item $u$ is smooth;
\item $T\geq\delta\beta$ for some $\delta>0$;
\item $\operatorname{tr}_\beta T$ has a locally bounded measure density.
\end{enumerate}
\end{proposition}

\begin{proof}
Choose $C$ so that $v=u+C|z|^2$ is plurisubharmonic, and use nonnegative radial mollifiers $\chi_\epsilon$ for which $v_\epsilon=v*\chi_\epsilon\downarrow v$. The smooth matrix convolutions $T_\epsilon=T*\chi_\epsilon$ can be written as
\[
  T_\epsilon=(\vartheta*\chi_\epsilon-C\beta)+\ddc v_\epsilon.
\]
Bedford--Taylor monotone continuity and smooth convergence of the background imply $T_\epsilon^k\to T^k$ for every $k$.

\emph{(2)$\Rightarrow$(3).} Assume (2). For every Hermitian matrix $M\geq\delta I$,
\[
  \det M\geq\delta^{n-1}\bigl(\operatorname{tr} M-(n-1)\delta\bigr).
\]
Apply this to the regularizations and pass to currents. Equation \eqref{eq:interior-local-ma} gives
\[
  \operatorname{tr}_\beta T\leq\delta^{1-n}F+(n-1)\delta,
\]
which is (3).

\emph{(3)$\Rightarrow$(1).} Assume (3), say $\operatorname{tr}_\beta T\leq L$. Positivity makes every coefficient of the matrix-valued current $T$ a bounded density. Hence the Euclidean Laplacian of $u$ is locally bounded, and interior Poisson estimates give $u\in W^{2,q}_{\mathrm{loc}}$ for every finite $q$. Mollification converges in $W^{2,q}$; taking $q>2n$ identifies the bounded-potential determinant with the ordinary almost-everywhere determinant. Thus $\det_\beta T=F$ almost everywhere. On a smaller ball with $F\geq m>0$, all eigenvalues of $T$ are at most $L$ and their product is at least $m$, so
\[
  \frac{m}{L^{n-1}}I\leq T\leq LI
\]
almost everywhere.

\Needspace{9\baselineskip}
It remains to show that this strong solution is smooth. Let $P$ be the real linear map from real symmetric $2n\times2n$ Hessian matrices to Hermitian $n\times n$ matrices, sending $D^2u$ to the coefficient matrix of $\ddc u$ relative to $\beta$, and let $\Theta(x)$ be the matrix of $\vartheta$. With $\delta=mL^{1-n}$, define the concave extension
\[
  G(x,M)=\inf_{\delta I\leq A\leq LI}
  \left\{\log\det A+\operatorname{tr}\!\left[A^{-1}\bigl(\Theta(x)+P(M)-A\bigr)\right]\right\}-\log F(x).
\]
For fixed $A$, write $\mathcal E_A(x,M)$ for the expression in braces
minus $\log F(x)$. The tangent-plane formula for $\log\det$ gives
\[
  G(x,D^2u)=0\quad\text{almost everywhere}.
\]
The operator $G$ is concave and uniformly elliptic on real symmetric matrices. The affine operators $\mathcal E_A$ have constant principal matrices with eigenvalues bounded above and below in terms of $\delta$ and $L$. Their $x$-dependence is smooth uniformly in $A$, so $G$ has a uniform local Lipschitz modulus in $x$.

We claim that $u$ is a viscosity solution of $G(x,D^2u)=0$. Suppose that a smooth function $\psi$ touches $u$ from above at $x_0$ and that $G(x_0,D^2\psi)<0$. Some fixed matrix $A$ then satisfies $\mathcal E_A(x,D^2\psi)<-\eta$ on a small ball. Since $G(x,D^2u)=0$ almost everywhere, $\mathcal E_A(x,D^2u)\geq0$ there. Thus $u-\psi$ has an interior maximum and satisfies
\[
  \operatorname{tr}\!\left(A^{-1}P(D^2(u-\psi))\right)\geq\eta
  \quad\text{almost everywhere}.
\]
After a small quadratic perturbation makes the contact strict, this contradicts the Aleksandrov maximum principle \cite[Chapter~9, Theorem~9.1]{GT01}. If instead $\psi$ touches $u$ from below and $G(x_0,D^2\psi)>0$, then $\mathcal E_A(x,D^2\psi)>\eta$ for every admissible $A$ on a small ball. At almost every point choose an admissible $A(x)$ with $\mathcal E_{A(x)}(x,D^2u)<\eta/2$. A measurable choice is obtained by taking the first qualifying element of a countable dense family. The function $\psi-u$ then violates the same maximum principle for the measurable uniformly elliptic coefficients $A(x)^{-1}$. This proves the claim.

The classical variable-coefficient Evans--Krylov theorem gives $u\in C^{2,\alpha}_{\mathrm{loc}}$ for some $0<\alpha<1$. The original equation is now uniformly elliptic with smooth data. Difference quotients satisfy linear equations whose principal coefficients are averages of inverse matrices along positive definite segments. Interior Schauder estimates first give $u\in C^{3,\alpha}_{\mathrm{loc}}$ and then, by iteration, $u\in C^\infty$. This proves (3)$\Rightarrow$(1).

\emph{(1)$\Rightarrow$(2).} If $u$ is smooth, then $T$ is a smooth semipositive form with determinant $F>0$, so it is positive definite; on a smaller neighborhood this gives (2).
\end{proof}

\Needspace{12\baselineskip}
\subsection{Smooth exact approximants}\label{sec:interior-approximants}

We now construct the approximants of the first step. The smooth Dirichlet theorem, \cref{lem:ball-smooth-dirichlet}, applies to every smooth positive density $F$. Approximating the H\"older boundary values of $\varphi$ by smooth functions therefore gives smooth solutions, and the point is to control them quantitatively. The estimates below bound the trace of a smooth solution by a power of the $C^4$ norm of its boundary value. The $C^4$ norms of the boundary approximations grow polynomially while their uniform errors decay polynomially, so both the convergence and the trace growth have polynomial rates.

In this subsection balls are centered at the origin. On the boundary sphere, $\nu$ is the outward Euclidean unit normal and $\nabla_{\tan}$ denotes tangential differentiation. For a positive form $T$, inner products $\langle\cdot,\cdot\rangle_T$ and norms $|\cdot|_T$ of $(1,0)$-covectors are computed with the inverse of the coefficient matrix of $T$.

\Needspace{20\baselineskip}
\begin{lemma}\label{lem:interior-trace-peak}
Let $B$ be a coordinate ball, let $\vartheta$ and $F>0$ be smooth near $\overline B$, and let $q$ be smooth with $R=\vartheta+\ddc q\geq\varepsilon\beta$ for a fixed $\varepsilon>0$. Fix a finite $L\geq0$. If $u\in C^\infty(B)$ satisfies
\[
  T=\vartheta+\ddc u>0,
  \qquad T^n=F\beta^n,
  \qquad |u-q|\leq L,
\]
put $S=\operatorname{tr}_\beta T$ and $v=u-q$. For $A>0$, define
\[
  Q_A=\log S-Av+\mathrm{e}^{-v}.
\]
There are $A_0,C$, depending only on the fixed data and $L$, such that at every interior local maximum of $Q_A$, with $A\geq A_0$,
\begin{equation}\label{eq:interior-trace-peak}
  S\leq C\sqrt A.
\end{equation}
The constants do not depend on any bound for the gradient of $u$.
\end{lemma}

\begin{proof}
This is the maximum-principle computation behind the classical second-order estimate. The non-closedness of $\vartheta$ produces a first-order error term, which is absorbed by the convex term $\mathrm{e}^{-v}$.

At the point in question make a constant unitary change of coordinates so that the matrix of $T$ is diagonal, with eigenvalues $\lambda_i$. Write
\[
  \Delta_T=\sum_i\lambda_i^{-1}\partial_i\partial_{\bar i},
  \qquad H=\operatorname{tr}_T\beta.
\]
Differentiating $\log\det T=\log F$ and commuting derivatives gives
\begin{equation}\label{eq:interior-trace-differentiation}
  \Delta_TS=\Delta_\beta\log F+P+E_0,
  \quad
  P=\sum_{k,i,j}\frac{|\partial_kT_{i\bar j}|^2}{\lambda_i\lambda_j},
  \quad |E_0|\leq C_0H.
\end{equation}
The non-closedness of $\vartheta$ enters through the bounded vector
\[
  E_p=\sum_k(\partial_p\vartheta_{k\bar k}-\partial_k\vartheta_{p\bar k}).
\]
\Needspace{6\baselineskip}
Since $\partial_pS-E_p=\sum_k\partial_kT_{p\bar k}$, weighted Cauchy--Schwarz gives $|\partial S-E|_T^2\leq SP$. From \eqref{eq:interior-trace-differentiation},
\begin{equation}\label{eq:interior-log-trace}
  \Delta_T\log S
  \geq-\frac{C_0(1+H)}S-
  \frac2S\operatorname{Re}\langle\partial\log S,E\rangle_T.
\end{equation}

Set $g_A(s)=-As+\mathrm{e}^{-s}$. At a local maximum of $Q_A$, $\partial\log S=-g_A'(v)\partial v$ and $\Delta_Tv=n-\operatorname{tr}_TR$. Combining these identities with \eqref{eq:interior-log-trace}, then completing the square in $\partial v$, yields
\begin{equation}\label{eq:interior-peak-inequality}
  0\geq
  \left[\varepsilon A-\frac{C_2}{S}-\frac{C_2(A+1)^2}{S^2}\right]H
  -C_3(A+1)-\frac{C_2}{S}.
\end{equation}
Here $|E|_T^2\leq CH$ and the bounds for $v$ have been used. If $S\geq K\sqrt A$ and $K$ is large, the bracket is at least $\varepsilon A/2$. Thus $H\leq C$ independently of $A$ and $u$. The determinant upper bound then gives $S\leq C'$. Enlarging $K$ makes this incompatible with $S\geq K\sqrt A$, proving \eqref{eq:interior-trace-peak}.
\end{proof}

We shall also use the following variant, in which a gradient bound replaces the convex term. If $|\nabla(u-q)|_\beta$ is bounded, then at an interior maximum of $\log S-A(u-q)$ the trace is bounded independently of $A$ for all large $A$. Indeed, at the maximum $\partial\log S=A\partial(u-q)$, and \eqref{eq:interior-log-trace} gives
\[
  0\geq\left[A\varepsilon-\frac CS-\frac{CA}{S}\right]H-nA-\frac CS.
\]
For $S$ larger than a fixed constant the bracket is at least $A\varepsilon/2$; the determinant again controls $S$.

\Needspace{19\baselineskip}
\paragraph{A polynomial boundary trace estimate}

\begin{proposition}\label{prop:interior-global-trace}
Let $B$ be the unit ball. Fix smooth $\vartheta$, smooth $F>0$ near $\overline B$, and $M<\infty$. There are $C,N<\infty$, depending only on these fixed data, such that every smooth solution $u$ with boundary value $g\in C^\infty(\partial B)$ of
\[
  T=\vartheta+\ddc u>0,
  \qquad T^n=F\beta^n,
  \qquad \lVert u\rVert_{L^\infty(B)}\leq M,
  \qquad u|_{\partial B}=g,
\]
satisfies
\begin{equation}\label{eq:interior-polynomial-global-trace}
  \sup_{\overline B}\operatorname{tr}_\beta T
  \leq C\bigl(1+\lVert g\rVert_{C^4(\partial B)}\bigr)^N.
\end{equation}
\end{proposition}

\begin{proof}
We follow the classical boundary estimates for the Dirichlet problem on a ball and keep track of the dependence on $L=1+\lVert g\rVert_{C^4}$. Barriers bound the first derivatives, and differentiation along rotations bounds the mixed normal--tangential second derivatives on the boundary. A lower bound for the tangential block of $T$ then bounds the double normal derivative through the equation, and \cref{lem:interior-trace-peak} reduces the interior bound to the boundary bound.

\emph{Step 1: first derivatives.} Fix a sufficiently large constant $K>0$, depending only on the fixed data. Extend $g$ to $G$ on $\overline B$ with $\lVert G\rVert_{C^4}\leq CL$. For $A=CL$ large enough,
\[
  \underline u=G+A(|z|^2-1),
  \qquad
  \underline T=\vartheta+\ddc\underline u\geq K\beta,
  \qquad
  \underline T^n\geq F\beta^n.
\]
\Cref{prop:interior-comparison} gives $\underline u\leq u$. If $h$ solves $\Delta_\beta h=-C_\vartheta$ with $h=g$ on $\partial B$ and $C_\vartheta\geq\sup\operatorname{tr}_\beta\vartheta$, then $u\leq h$. Boundary differentiation of these barriers gives
\begin{equation}\label{eq:interior-first-boundary}
  |\partial_\nu u|+|\nabla_{\tan}u|\leq CL
  \quad\text{on }\partial B.
\end{equation}

Let $H=\operatorname{tr}_T\beta$, let $\mathcal L=\operatorname{tr}_T\ddc$ be the linearized operator, and let $w=u-\underline u\geq0$ be the barrier function used below. Since $H\geq nF^{-1/n}$ is bounded below, taking $K$ fixed and sufficiently large gives
\begin{equation}\label{eq:interior-w-barrier}
  \mathcal Lw=n-\operatorname{tr}_T\underline T\leq-(1+H).
\end{equation}
Let $E=\sum_k(z_k\partial_{z_k}+\bar z_k\partial_{\bar z_k})$ be the real radial field. Write $\Theta$ for the coefficient matrix of $\vartheta$. Differentiating $\log\det T=\log F$ and commuting $E$ with complex second derivatives yields
\[
  \mathcal L(Eu)=E(\log F)+2n-\operatorname{tr}_T(E\Theta+2\Theta).
\]
Hence $|\mathcal L(Eu)|\leq C(1+H)$. The same estimate holds for $Zu$ when $Z$ is an infinitesimal unitary rotation. Applying the minimum principle to $CL+C'w\pm Yu$, with $Y=E$ or $Z$, gives
\begin{equation}\label{eq:interior-first-fields}
  |Eu|+|Zu|\leq CL\quad\text{on }B.
\end{equation}

\emph{Step 2: mixed second derivatives on the boundary.} For a rotation $Z$, the function $R_Z=Z(u-G)$ vanishes on $\partial B$ and satisfies $|\mathcal LR_Z|\leq CL(1+H)$. The barriers $CLw\pm R_Z$ and \eqref{eq:interior-w-barrier} imply $|\partial_\nu R_Z|\leq CL^2$ on the boundary. Rotations span the real tangent space of $\partial B$, so all real mixed normal--tangential second derivatives are $O(L^2)$. Pure tangential second derivatives are $O(L)$ by differentiating $u=g$ twice and using \eqref{eq:interior-first-boundary}. Thus all complex Hessian entries except the complex normal--normal entry are $O(L^2)$.

\emph{Step 3: the double normal derivative.} It remains to bound the tangential block of $T$ from below. For $n\geq2$, let
\[
  \lambda_0=\min\{T_p(\xi,\bar\xi):p\in\partial B,
  \ \xi\in T_p^{1,0}(\partial B),\ |\xi|_\beta=1\}.
\]
Choose a minimizing pair $(p,\xi)$ and extend $\xi$ to a smooth unit tangential field near $p$, with uniformly bounded derivatives. On the sphere, any smooth function vanishing there satisfies
\[
  \ddc(u-G)(\xi,\bar\xi)=\frac12E(u-G).
\]
Define
\[
  a(z)=(\vartheta+\ddc G)_z(\xi(z),\overline{\xi(z)})-\frac12EG(z),
  \qquad
  P=a+\frac12Eu-\lambda_0.
\]
Then $P\geq0$ on the nearby spherical boundary and $P(p)=0$.
The boundary first-derivative bound also gives $\lambda_0\leq CL$ by
the tangential identity above. Consequently
\[
  |P|\leq CL,
  \qquad
  |\mathcal LP|\leq CL(1+H).
\]
On a fixed cap $D=B\cap B_r(p)$ consider
\[
  \Psi=P+C_1Lw+C_2L|z-p|^2.
\]
Choose $C_2$ so that $\Psi\geq0$ on the interior spherical part of $\partial D$, and then $C_1$ so that $\mathcal L\Psi\leq0$ by \eqref{eq:interior-w-barrier}. The minimum principle yields $\Psi\geq0$. At $p$ the normal derivative of $|z-p|^2$ vanishes, and therefore
\[
  \partial_\nu P+C_1L\partial_\nu w\leq0.
\]
Equations \eqref{eq:interior-first-boundary}--\eqref{eq:interior-first-fields} now give $D^2u(\nu,\nu)(p)\leq CL^2$, hence $\operatorname{tr}_\beta T(p)\leq CL^2$. Positivity and $\det_\beta T=F$ imply
\begin{equation}\label{eq:interior-tangential-lower}
  \lambda_0\geq cL^{-2(n-1)}.
\end{equation}

\Needspace{5\baselineskip}
At any boundary point write
\[
  T=\begin{pmatrix}B_0&b\\ b^*&d\end{pmatrix}
\]
in tangential--normal blocks, with $B_0$ an $(n-1)\times(n-1)$ Hermitian matrix, $b\in\mathbb C^{n-1}$ and $d\in\mathbb R$. We have $|b|\leq CL^2$ and $B_0\geq\lambda_0I$. The determinant equation and the Schur complement give
\[
  d=\frac{F}{\det B_0}+b^*B_0^{-1}b
  \leq(\sup F)\lambda_0^{-(n-1)}+CL^4\lambda_0^{-1},
\]
which is polynomial in $L$. For $n=1$, the trace is $F$ and this argument is unnecessary.

\emph{Step 4: the interior.} Choose a fixed quadratic $q$ with $\vartheta+\ddc q\geq\varepsilon\beta$. Since $|u-q|\leq M+\sup_B|q|$, \cref{lem:interior-trace-peak} applies with this bound in place of its constant $L$; we use it with one fixed $A$. At an interior maximum of its auxiliary function the trace is uniformly bounded, while at a boundary maximum the polynomial boundary bound of Steps~2 and~3 applies. Since the remaining terms are uniformly bounded, \eqref{eq:interior-polynomial-global-trace} follows.
\end{proof}

\Needspace{23\baselineskip}
\paragraph{Exact approximation}

We apply these estimates to smooth approximations of the boundary values of $\varphi$ on a ball in $\Omega$. Recall that $\theta$, $\omega$, $c$ and $f$ are the data of \cref{thm:nonclosed-interior-smoothness}.

\Needspace{18\baselineskip}
\begin{proposition}\label{prop:interior-tame-approximants}
Let $B\Subset\Omega$ be a coordinate ball, identified with $\{|z|<r_0\}$, with closure in a larger coordinate chart, and put
\[
  F=\frac{cf\,\omega^n}{\beta^n}.
\]
There are $a>0$, $b,C<\infty$, and $u_j\in C^\infty(\overline B)$ such that
\begin{equation}\label{eq:interior-exact-approximants}
  \theta+\ddc u_j>0,
  \qquad
  (\theta+\ddc u_j)^n=F\beta^n,
\end{equation}
and
\begin{equation}\label{eq:interior-tame-rates}
  \lVert u_j-\varphi\rVert_{L^\infty(B)}\leq Cj^{-a},
  \qquad
  \sup_{\overline B}\operatorname{tr}_\beta(\theta+\ddc u_j)\leq Cj^b.
\end{equation}
\end{proposition}

\begin{proof}
Since $\varphi$ is H\"older continuous near $\overline B$, mollification at scale $1/j$ gives smooth boundary approximations $g_j$ with
\[
  \lVert g_j-\varphi\rVert_{L^\infty(\partial B)}\leq Cj^{-a},
  \qquad
  \lVert g_j\rVert_{C^4(\partial B)}\leq Cj^4.
\]
\Needspace{7\baselineskip}
Choose a fixed quadratic $q$ with $\eta=\theta+\ddc q>0$ near $\overline B$. Extend $g_j$ smoothly to $G_j$. For a sufficiently large, polynomially bounded $A_j$, the function
\[
  \underline w_j=G_j-q+A_j(|z|^2-r_0^2)
\]
has boundary value $g_j-q$ and satisfies
\[
  \eta+\ddc\underline w_j
  =\theta+\ddc G_j+A_j\beta>0,
  \qquad
  (\eta+\ddc\underline w_j)^n\geq F\beta^n.
\]
Apply the smooth Dirichlet theorem recalled at the beginning of this subsection with the background $\eta$, the density $F$ and the boundary value $g_j-q$. It gives a smooth solution $w_j$ of $(\eta+\ddc w_j)^n=F\beta^n$ with $\eta+\ddc w_j>0$; put $u_j=w_j+q$. Comparing $u_j$ with $\varphi\pm\lVert g_j-\varphi\rVert_\infty$ by \cref{prop:interior-comparison} proves the first bound in \eqref{eq:interior-tame-rates}. The functions $u_j$ are therefore uniformly bounded. \Cref{prop:interior-global-trace}, applied after rescaling $B$ to the unit ball, and the polynomial $C^4$ boundary norm give the second bound.
\end{proof}

\subsection{Localization near the nonregular set}\label{sec:interior-localization}

In the auxiliary function of \cref{lem:interior-trace-peak}, the term $-A(u_j-q)$ penalizes points where $u_j-q$ is large. A quadratic gap of $u-q$ along the nonregular set therefore suppresses the contribution of the boundary of a neighborhood of that set. Taking $A$ proportional to $\log j$ makes this suppression beat the polynomial growth of the boundary trace, and near the point of minimal $u-q$ the trace then grows more slowly than every positive power of $j$. Interpolation with the polynomial approximation rate gives uniform $C^{1,\alpha}$ bounds.

Distances and norms below are Euclidean. By a smooth exact solution for $(\vartheta,F)$ we mean a smooth function $u_j$ with $\vartheta+\ddc u_j>0$ and $(\vartheta+\ddc u_j)^n=F\beta^n$.

\Needspace{28\baselineskip}
\begin{proposition}\label{prop:interior-subpower}\label{prop:interior-uniform-c1a}
For $R>0$, let $B=B_R(0)$, let $\vartheta$ and $F>0$ be smooth near $\overline B$, and let $u\in C(\overline B)$. Suppose $q$ is smooth and $\vartheta+\ddc q\geq\varepsilon\beta$ for some $\varepsilon>0$. Put $v=u-q$ and $m=v(0)$, and fix $\kappa>0$. Let $\mathcal S\subset\overline B$ be closed, contain $0$, and satisfy
\begin{equation}\label{eq:interior-quadratic-separation}
  v(0)=m,
  \qquad
  v(z)\geq m+\kappa|z|^2\quad(z\in\mathcal S).
\end{equation}
Assume that smooth exact solutions $u_j$, $j\geq2$, satisfy, for fixed constants $a>0$, $b\geq0$ and $C>0$,
\begin{align}
  \lVert u_j-u\rVert_\infty&\leq Cj^{-a}, \label{eq:interior-localization-a}\\
  \sup_{\overline B}\operatorname{tr}_\beta(\vartheta+\ddc u_j)&\leq Cj^b, \label{eq:interior-localization-b}\\
  \sup_j\sup_K\operatorname{tr}_\beta(\vartheta+\ddc u_j)&<\infty
  \quad\text{for every }K\Subset\overline B\setminus\mathcal S. \label{eq:interior-localization-c}
\end{align}
Then for every $\gamma>0$ there are a neighborhood $U_\gamma$ of $0$ and $C_\gamma$ such that
\begin{equation}\label{eq:interior-subpower-trace}
  \operatorname{tr}_\beta(\vartheta+\ddc u_j)\leq C_\gamma j^\gamma
  \quad\text{on }U_\gamma.
\end{equation}
For every $0<\alpha<1$, there is also a neighborhood $U$ of $0$ on which $u_j\to u$ in $C^{1,\alpha}$. In particular the approximants are uniformly bounded in that norm and $u\in C^{1,\alpha}(U)$.
\end{proposition}

\begin{proof}
\emph{Step 1: the trace bound \eqref{eq:interior-subpower-trace}.} Write $e_j=Cj^{-a}$, fix $d=\kappa R^2/2$, and choose $D>(b+1)/d$. Choose $\eta>0$ so that $2D\eta<\gamma/2$. For small $\delta>0$, the neighborhood
\[
  W_\delta=\{z\in B:\operatorname{dist}(z,\mathcal S)<\delta\}
\]
contains $0$, has $v\geq m-\eta$ on its closure, and satisfies $v\geq m+d$ on $\partial B\cap\overline W_\delta$ with the fixed $d$ above. These properties follow from \eqref{eq:interior-quadratic-separation}, continuity, and compactness. The remaining part of the boundary of $W_\delta$ is a compact subset of $\overline B\setminus\mathcal S$, where \eqref{eq:interior-localization-c} applies.

Set $A_j=D\log j$. Apply the maximum principle on $\overline W_\delta$ to
\[
  Q_j=\log\operatorname{tr}_\beta(\vartheta+\ddc u_j)-A_j(u_j-q)+\mathrm{e}^{-(u_j-q)}.
\]
At an interior maximum, \cref{lem:interior-trace-peak} gives $\log\operatorname{tr}_\beta(\vartheta+\ddc u_j)\leq C+\frac12\log A_j$. On the part of $\partial W_\delta$ inside $B$, which stays away from $\mathcal S$, the trace has a bound independent of $j$. On the outer part $\partial B\cap\overline W_\delta$, \eqref{eq:interior-localization-b} and the gap $d$ give
\[
  Q_j\leq C+b\log j-A_j(m+d-e_j),
\]
and the choice $Dd>b+1$ makes this smaller than the bound below for large $j$. Hence throughout $W_\delta$,
\[
  Q_j\leq C_\delta+\tfrac12\log A_j-A_j(m-\eta-e_j).
\]
On a small ball about $0$ where $v\leq m+\eta$, this implies
\[
  \log\operatorname{tr}_\beta(\vartheta+\ddc u_j)
  \leq C_\delta+\tfrac12\log A_j+A_j(2\eta+2e_j).
\]
By the choice of $\eta$, the logarithmic factor and $A_je_j$ are absorbed into the remaining power. This proves \eqref{eq:interior-subpower-trace}.

\emph{Step 2: $C^{1,\alpha}$ convergence.} On a smaller ball, the real Laplacian of $u_j$ is bounded by $Cj^\gamma$. Interior Poisson estimates give a $W^{2,q_0}$ bound for every finite $q_0$, and Sobolev embedding with $q_0$ large then gives, for every $\sigma<1$,
\[
  \lVert u_j\rVert_{W^{2,q_0}}+\lVert u_j\rVert_{C^{1,\sigma}}
  \leq Cj^\gamma.
\]
For $d_k=u_{2^{k+1}}-u_{2^k}$,
\[
  \lVert d_k\rVert_\infty\leq C2^{-ka},
  \qquad
  \lVert d_k\rVert_{C^{1,\sigma}}\leq C2^{k\gamma}.
\]
If $\alpha<\sigma$ and $t=(1+\alpha)/(1+\sigma)$, H\"older interpolation gives
\[
  \lVert d_k\rVert_{C^{1,\alpha}}
  \leq C2^{k[-a(1-t)+\gamma t]}.
\]
Since $\gamma$ may be arbitrarily small, the exponent can be made negative. The dyadic series converges in $C^{1,\alpha}$ to its uniform limit $u$. If $2^k\leq j<2^{k+1}$, interpolation applied to $u_j-u_{2^k}$ gives the same decaying bound. Hence the full sequence converges to $u$ in $C^{1,\alpha}$ on a smaller ball.
\end{proof}

\begin{proposition}\label{prop:interior-regular-stability}
Suppose $u_j$ solve the same smooth equation as in \cref{prop:interior-subpower}, satisfy \eqref{eq:interior-localization-a}--\eqref{eq:interior-localization-b}, and $u$ is smooth with $\vartheta+\ddc u>0$ on an open set $O\subset B$. Then, for every $K\Subset O$,
\begin{equation}\label{eq:interior-regular-stability}
  \sup_j\sup_K\operatorname{tr}_\beta(\vartheta+\ddc u_j)<\infty.
\end{equation}
\end{proposition}

\begin{proof}
Where $u$ is smooth and strictly $\vartheta$-plurisubharmonic, it provides its own quadratic separation, so the argument of \cref{prop:interior-subpower} applies around every point of $O$. A first application gives trace growth of order $\sqrt{\log j}$; the resulting gradient convergence then allows the variant of \cref{lem:interior-trace-peak} with a gradient bound, which removes this growth.

For $p$ in a fixed compact subset of $O$, take a uniformly sized ball $B_r(p)\Subset O$ and set
\[
  q_p(z)=u(z)-\kappa|z-p|^2
\]
with $\kappa>0$ small, so $\vartheta+\ddc q_p\geq\varepsilon\beta$. If $e_j=Cj^{-a}$ and $v_{j,p}=u_j-q_p$, then
\[
  v_{j,p}\geq-e_j,
  \quad v_{j,p}(p)\leq e_j,
  \quad v_{j,p}\geq\kappa r^2-e_j\text{ on }\partial B_r(p).
\]
With $A_j=D\log j$ and $D\kappa r^2>b+1$, the maximum argument of Step~1 in the proof of \cref{prop:interior-subpower} gives
\begin{equation}\label{eq:interior-sqrtlog-trace}
  \operatorname{tr}_\beta(\vartheta+\ddc u_j)(p)\leq C\sqrt{\log j}.
\end{equation}
The constants are uniform on compact subsets of $O$.

Put $T_j=\vartheta+\ddc u_j$ and $w_j=u_j-u$. On slightly larger compact sets, \eqref{eq:interior-sqrtlog-trace} bounds $|\Delta w_j|$ by $C\sqrt{\log j}$. Poisson estimates and Sobolev embedding give
\[
  \lVert w_j\rVert_{C^{1,\sigma}}\leq C_\sigma\sqrt{\log j},
  \qquad \sigma<1.
\]
Interpolating with $\lVert w_j\rVert_\infty\leq Cj^{-a}$ shows $\nabla w_j\to0$ uniformly on smaller compact sets. The gradients of $v_{j,p}$ are therefore uniformly bounded. Repeat the maximum argument with $\log\operatorname{tr}_\beta T_j-A_jv_{j,p}$, now using the variant of \cref{lem:interior-trace-peak} with a gradient bound. Its interior trace bound is now independent of $A_j$, while the boundary term is suppressed as before. Evaluation at $p$ gives
\[
  \operatorname{tr}_\beta T_j(p)\leq C\exp(2A_je_j)\leq C,
\]
proving \eqref{eq:interior-regular-stability}.
\end{proof}



\Needspace{24\baselineskip}
\subsection{Derivative estimates and bounded traces}\label{sec:interior-calabi}

Uniform $C^{1,\alpha}$ bounds control averages of the trace $\operatorname{tr}_\beta(\vartheta+\ddc u_j)=\operatorname{tr}_\beta\vartheta+\Delta_\beta u_j$ over small balls. To pass from averages to pointwise values we need a bound for the gradient of the trace, that is, a third-order estimate, whose dependence on the ellipticity constants is polynomial. We prove it by Calabi's maximum-principle argument, computed with the Chern connection of the solution so that the torsion of the nonclosed background appears explicitly. The resulting averaging inequality is nonlinear in the trace, and it closes under iteration when $\alpha$ is sufficiently close to one.

In this subsection $B$ is the unit ball and $B_r(p)$ is the Euclidean ball of radius $r$ centered at $p$. For a Hermitian metric $g$, $\partial g$ denotes the ordinary coordinate derivatives of its coefficient matrix, while tensor derivatives in the proofs are taken with respect to the Chern connection of $g$.

\begin{proposition}\label{prop:interior-calabi}
Fix smooth $\vartheta$ and smooth $F>0$ near the closed unit ball. There are $C,N<\infty$, depending only on finitely many derivatives of the fixed data, such that any smooth solution
\[
  g=\vartheta+\ddc u,
  \qquad g^n=F\beta^n,
  \qquad M^{-1}\beta\leq g\leq M\beta,
  \qquad M\geq1,
\]
satisfies
\begin{equation}\label{eq:interior-calabi-polynomial}
  \sup_{B_{1/2}}|\partial g|_\beta\leq CM^N.
\end{equation}
If only $\sup_B\operatorname{tr}_\beta g\leq L$ is given, the same bound holds with a fixed power of $1+L$. On a ball of radius $r\leq1$,
\begin{equation}\label{eq:interior-gradient-trace-scaled}
  \sup_{B_{r/2}}|\nabla\operatorname{tr}_\beta g|_\beta
  \leq Cr^{-1}(1+L)^N.
\end{equation}
\end{proposition}

\begin{proof}
The quantity $\mathcal R$ below is the squared norm of the difference between the Chern connection of $g$ and the flat connection. We derive a differential inequality for $\mathcal R$ in which the torsion enters through fixed derivatives of $\vartheta$, and combine it with the trace in a maximum principle; every constant is kept a power of $M$.

We compare the flat connection of $\beta$ with the Chern connection $\nabla$ of $g$. Write $g^{k\bar\ell}$ for the coefficients of the inverse matrix and $\bar\nabla$ for covariant differentiation in $(0,1)$ directions; tensor norms and inner products without subscript are computed with $g$. The difference of the two connections and its squared norm are
\[
  \Psi^k_{ij}=g^{k\bar\ell}\partial_i g_{j\bar\ell},
  \qquad \mathcal R=|\Psi|_g^2.
\]
The torsion of $\nabla$ is determined by the fixed background, because the third derivatives of $u$ cancel in it. Its lowered and raised components are
\begin{equation}\label{eq:interior-torsion-background}
  C_{ij\bar\ell}
  =\partial_i\vartheta_{j\bar\ell}-\partial_j\vartheta_{i\bar\ell},
  \qquad
  \mathcal T^k_{ij}=g^{k\bar\ell}C_{ij\bar\ell}.
\end{equation}
\Needspace{8\baselineskip}
Consequently, for a fixed integer $N$,
\begin{align}
  |\mathcal T|&\leq CM^N,\nonumber\\
  |\nabla\mathcal T|+|\bar\nabla\mathcal T|
  &\leq CM^N(1+\sqrt{\mathcal R}),\label{eq:interior-torsion-derivatives}\\
  |\nabla^2\mathcal T|+|\nabla\bar\nabla\mathcal T|+|\bar\nabla^2\mathcal T|
  &\leq CM^N(1+\mathcal R+|\nabla\Psi|+|\bar\nabla\Psi|).\nonumber
\end{align}
Every factor here is either a fixed derivative of $\vartheta$, a connection factor $\Psi$, or a factor of $g^{\pm1}$; this explains the polynomial dependence.

For tensor fields, use $\Delta_g=g^{p\bar q}\nabla_p\nabla_{\bar q}$.
We need the following curvature identities. With
\[
  \operatorname{Rm}^k{}_{i\bar j l}=-\partial_{\bar j}\Psi^k_{il},
  \qquad
  \operatorname{Rm}_{i\bar j l\bar k}=g_{m\bar k}\operatorname{Rm}^m{}_{i\bar j l},
\]
set
\[
  \operatorname{Ric}^{(2)}_{l\bar k}=g^{i\bar j}\operatorname{Rm}_{i\bar j l\bar k},
  \qquad
  \operatorname{Ric}^{(1)}_{l\bar k}=g^{i\bar j}\operatorname{Rm}_{l\bar k i\bar j}
  =-\partial_l\partial_{\bar k}\log F.
\]
The two curvature interchanges are
\begin{align*}
  \operatorname{Rm}_{i\bar j l\bar k}-\operatorname{Rm}_{l\bar j i\bar k}
  &=-\nabla_{\bar j}C_{il\bar k},\\
  \operatorname{Rm}_{i\bar j l\bar k}-\operatorname{Rm}_{i\bar k l\bar j}
  &=-\nabla_i\overline{C}_{\bar j\bar k l}.
\end{align*}
They follow by differentiating the torsion identity and its conjugate. Successive interchange and contraction give
\begin{equation}\label{eq:interior-ricci-contraction}
  \operatorname{Ric}^{(2)}_{l\bar k}
  =\operatorname{Ric}^{(1)}_{l\bar k}
  -g^{i\bar j}\bigl(
    \nabla_l\overline C_{\bar j\bar k i}
    +\nabla_{\bar j}C_{il\bar k}
  \bigr).
\end{equation}
Metric compatibility, \eqref{eq:interior-torsion-derivatives}, and the fixed smooth formula for $\operatorname{Ric}^{(1)}$ imply
\begin{equation}\label{eq:interior-ricci2}
  |\operatorname{Ric}^{(2)}|\leq CM^N(1+\sqrt{\mathcal R}),
  \qquad
  |\nabla\operatorname{Ric}^{(2)}|
  \leq CM^N(1+\mathcal R+|\nabla\Psi|+|\bar\nabla\Psi|).
\end{equation}
Moreover,
\[
  \nabla_{\bar q}\Psi^k_{ij}=-\operatorname{Rm}^k{}_{i\bar q j},
\]
and the Chern--Bianchi identity, including its torsion term, is
\[
  \nabla_p\operatorname{Rm}^k{}_{i\bar q j}-\nabla_i\operatorname{Rm}^k{}_{p\bar q j}
  =-\mathcal T^r_{pi}\operatorname{Rm}^k{}_{r\bar q j}.
\]
Consequently,
\begin{equation}\label{eq:interior-laplacian-connection-difference}
  \Delta_g\Psi^k_{ij}
  =-\nabla_i(\operatorname{Ric}^{(2)})^k{}_j
  +g^{p\bar q}\mathcal T^r_{pi}\operatorname{Rm}^k{}_{r\bar q j}.
\end{equation}
If $A^k{}_l=(\operatorname{Ric}^{(2)})^k{}_l$, define its action on a tensor of type $(1,2)$ by
\[
  (A\mathbin{\cdot}\Psi)^k_{ij}
  =A^k{}_l\Psi^l_{ij}-A^l{}_i\Psi^k_{lj}-A^l{}_j\Psi^k_{il}.
\]
Differentiating the squared norm and using metric compatibility gives
\begin{equation}\label{eq:interior-laplacian-calabi-norm}
  \Delta_g\mathcal R
  =|\nabla\Psi|^2+|\bar\nabla\Psi|^2
   +2\operatorname{Re}\langle\Delta_g\Psi,\Psi\rangle
   -\operatorname{Re}\langle A\mathbin{\cdot}\Psi,\Psi\rangle.
\end{equation}
There is no mixed torsion term because the Chern torsion has type $(2,0)$. Using \eqref{eq:interior-torsion-derivatives}, \eqref{eq:interior-ricci2}, and \eqref{eq:interior-laplacian-connection-difference} in \eqref{eq:interior-laplacian-calabi-norm}, Young's inequality absorbs the products of $\sqrt{\mathcal R}$ with $|\nabla\Psi|+|\bar\nabla\Psi|$. We obtain
\begin{equation}\label{eq:interior-calabi-differential}
  \Delta_g\mathcal R\geq-CM^N(\mathcal R^{3/2}+1).
\end{equation}
Let $s=\operatorname{tr}_\beta g$. Direct differentiation as in \eqref{eq:interior-trace-differentiation} gives
\begin{equation}\label{eq:interior-trace-calabi-coupling}
  \Delta_gs\geq M^{-1}\mathcal R-CM,
  \qquad
  |\partial s|_g^2\leq nM^2\mathcal R.
\end{equation}
Take a cutoff $\zeta$ supported in $B$, equal to one on $B_{1/2}$, set $\eta=\zeta^2$ and $h=2nM-s$, so $nM\leq h\leq2nM$, and maximize
\[
  Q=\frac{\eta\mathcal R}{h}.
\]
At a positive interior maximum of $Q$, the identity
$\partial\mathcal R/\mathcal R=\partial h/h-\partial\eta/\eta$ gives
\[
 \begin{split}
 0\geq\frac{\Delta_gQ}{Q}
 &=\frac{\Delta_g\mathcal R}{\mathcal R}
   +\frac{\Delta_g\eta}{\eta}-\frac{\Delta_gh}{h}
   -\frac{2|\partial\eta|_g^2}{\eta^2}
   +\frac{2\operatorname{Re}\langle\partial\eta,\partial h\rangle_g}{\eta h}.
 \end{split}
\]
We have $|\partial\eta|_g^2\leq CM\eta$,
$|\Delta_g\eta|\leq CM$ and
$|\partial h|_g\leq\sqrt n M\sqrt{\mathcal R}$.
Insert \eqref{eq:interior-calabi-differential} and
\eqref{eq:interior-trace-calabi-coupling} and multiply by $\eta Q$.
The positive term is $M^{-1}Q^2$. The two terms involving
$\sqrt{\mathcal R}$ are bounded in absolute value by a fixed power of
$M$ times $Q^{3/2}$, since $\mathcal R=hQ/\eta$ and $h\leq2nM$.
The remaining cutoff terms are bounded by $CM^N(Q+1)$. Therefore
\[
 0\geq M^{-1}Q^2-CM^N(Q^{3/2}+Q+1).
\]
Thus $Q$, and then $\mathcal R$ on $B_{1/2}$, are bounded by a fixed polynomial in $M$. Since $|\partial g|_\beta^2\leq M^3\mathcal R$, this proves \eqref{eq:interior-calabi-polynomial}.

If $s\leq L$, the determinant lower bound gives $\lambda_{\min}(g)\geq(\inf F)L^{1-n}$, so $M$ is a fixed power of $1+L$. For a ball $B_r(p)$, set $u_r(z)=r^{-2}u(p+rz)$ and use the
coefficient matrices $\vartheta_r(z)=\vartheta(p+rz)$,
$g_r(z)=g(p+rz)$ and $F_r(z)=F(p+rz)$. Then
$g_r=\vartheta_r+\ddc u_r$ in the unit coordinates and
$\det g_r=F_r$. For $r\leq1$ the derivatives of the fixed data remain
uniformly bounded. The unit-ball estimate for the ordinary coefficient derivatives $\partial g_r$ and their conjugates gives the
factor $r^{-1}$ in \eqref{eq:interior-gradient-trace-scaled}.
\end{proof}

\Needspace{12\baselineskip}
\begin{proposition}\label{prop:interior-c1a-to-trace}
There is $\alpha_*<1$, depending only on $n$ and the exponent $N$ in \cref{prop:interior-calabi}, with the following property. Let $B_0$ be a Euclidean ball, fix a smooth real $(1,1)$-form $\vartheta$ and a smooth function $F>0$ near $\overline{B_0}$, and let $u_j$ be smooth exact solutions of $(\vartheta+\ddc u_j)^n=F\beta^n$ with $\vartheta+\ddc u_j>0$. If
\[
  \sup_j\lVert u_j\rVert_{C^{1,\alpha}(B_0)}<\infty
  \qquad\text{for some }\alpha_*<\alpha<1,
\]
then $\operatorname{tr}_\beta(\vartheta+\ddc u_j)$ is uniformly bounded on every smaller ball.
\end{proposition}

\begin{proof}
Fix a center for concentric balls $B_r\Subset B_0$, put $s_j=\operatorname{tr}_\beta(\vartheta+\ddc u_j)$, and set
\[
  M_j(r)=1+\sup_{B_r}s_j.
\]
For concentric balls $B_r\Subset B_R\Subset B_0$, $d=R-r\leq1$, \cref{prop:interior-calabi} gives
\begin{equation}\label{eq:interior-trace-gradient-M}
  |\nabla s_j|\leq Cd^{-1}M_j(R)^N
  \quad\text{on }B_{r+d/2}.
\end{equation}
Fix $p\in B_r$ and let $\chi_h$ be a nonnegative mollifier of radius $h\leq d/4$. Since
\[
  s_j=\operatorname{tr}_\beta\vartheta+\Delta_\beta u_j,
\]
integration by parts, after subtracting the constant vector $\nabla u_j(p)$, and the uniform $C^{1,\alpha}$ bound give
\[
  (s_j*\chi_h)(p)\leq C(1+h^{\alpha-1}).
\]
Equation \eqref{eq:interior-trace-gradient-M} gives
\[
  |s_j(p)-(s_j*\chi_h)(p)|
  \leq Cd^{-1}M_j(R)^Nh.
\]
Choose
\[
  h=\frac d4M_j(R)^{-N/(2-\alpha)}.
\]
Taking the supremum over $p\in B_r$ and absorbing fixed powers of $d^{-1}$, we obtain
\begin{equation}\label{eq:interior-nonlinear-iteration}
  M_j(r)\leq Cd^{-1}M_j(R)^q,
  \qquad
  q=\frac{N(1-\alpha)}{2-\alpha}.
\end{equation}
Choose $\alpha_*$ so close to one that $q<1$ for $\alpha>\alpha_*$. Choose $0<r_0<R_0$ with $B_{R_0}\Subset B_0$. Let $r_k\uparrow R_0$ with $r_{k+1}-r_k=(R_0-r_0)2^{-k-1}$. Iterating \eqref{eq:interior-nonlinear-iteration} gives
\[
  M_j(r_0)
  \leq
  \prod_{k=0}^{m-1}\left[C(r_{k+1}-r_k)^{-1}\right]^{q^k}
  M_j(r_m)^{q^m}.
\]
The infinite product converges because $\sum q^k(1+k)<\infty$. For each fixed $j$, the final factor tends to one as $m\to\infty$. The resulting bound is independent of $j$.
\end{proof}

\Needspace{8\baselineskip}
\subsection{Proof of the main theorem}\label{sec:interior-main-proof}

We now assemble the four steps outlined after \cref{thm:nonclosed-interior-smoothness}. The barrier enters only once, but globally: it locates a point of the nonsmooth set at which the quadratic separation required in \cref{prop:interior-subpower} holds.

\begin{proof}[of \cref{thm:nonclosed-interior-smoothness}]
If $n=1$, the equation is the linear identity
\[
  \ddc\varphi=cf\,\omega-\theta
  \quad\text{on }\Omega.
\]
Its right-hand side is smooth, so interior elliptic regularity gives $\varphi\in C^\infty(\Omega)$, and the original equation gives $\theta+\ddc\varphi=cf\,\omega>0$. We henceforth assume $n\geq2$.

Let
\[
  \mathcal S=\{p\in\Omega:\varphi\text{ is not smooth on any neighborhood of }p\}.
\]
This is closed in $\Omega$. Suppose $\mathcal S\neq\varnothing$. By \cref{lem:interior-proper-gap}, the restriction of $\varphi-\rho$ to $\mathcal S$ attains its minimum $m$ at some $p\in\mathcal S$.

Choose nested coordinate balls
\[
  B\Subset B_0\Subset\Omega
\]
centered at $p$. On $B_0$, \cref{prop:interior-tame-approximants} gives smooth exact Dirichlet solutions $u_j$ converging uniformly to $\varphi$ at a polynomial rate and having polynomial global trace growth.

On the open set $\Omega\setminus\mathcal S$ the function $\varphi$ is smooth, and $\theta+\ddc\varphi$ is semipositive with positive determinant, hence positive definite. \Cref{prop:interior-regular-stability} therefore gives a uniform trace bound for the $u_j$ on every compact subset of $B_0\setminus\mathcal S$.

In coordinates centered at $p$, take $\kappa>0$ small and set
\[
  q(z)=\rho(z)-\kappa|z|^2.
\]
After shrinking $B$, the lower bound $\theta+\ddc\rho\geq\omega$ gives $\theta+\ddc q>0$ on $\overline B$. If $v=\varphi-q$, then the minimizing property of $p$ gives
\[
  v(p)=m,
  \qquad
  v(z)=(\varphi-\rho)(z)+\kappa|z|^2
  \geq m+\kappa|z|^2
  \quad(z\in\mathcal S\cap\overline B).
\]
Thus \cref{prop:interior-subpower} applies on $B$ with the closed set $\mathcal S\cap\overline B$: the approximation rate and global polynomial trace bound come from \cref{prop:interior-tame-approximants}, and the bound away from $\mathcal S$ comes from \cref{prop:interior-regular-stability}. Its second assertion gives, for every $0<\alpha<1$, a neighborhood of $p$ on which the approximants have a uniform $C^{1,\alpha}$ bound.

Choose $\alpha>\alpha_*$ as in \cref{prop:interior-c1a-to-trace}, and apply that proposition on a ball centered at $p$ inside this neighborhood. It yields
\[
  \sup_j\operatorname{tr}_\beta(\theta+\ddc u_j)<\infty
\]
on a still smaller neighborhood of $p$. Since $u_j\to\varphi$ uniformly, $\ddc u_j\to\ddc\varphi$ as currents. The trace bound passes to the limit, so $\operatorname{tr}_\beta(\theta+\ddc\varphi)$ has a bounded density near $p$. \Cref{prop:interior-criterion} now implies that $\varphi$ is smooth near $p$, contradicting $p\in\mathcal S$.

Hence $\mathcal S=\varnothing$ and $\varphi\in C^\infty(\Omega)$. The form $\theta+\ddc\varphi$ is smooth and semipositive there, while
\[
  (\theta+\ddc\varphi)^n=cf\,\omega^n>0.
\]
Its nonnegative eigenvalues have positive product, so every eigenvalue is positive. Therefore $\theta+\ddc\varphi>0$ on $\Omega$.
\end{proof}



\begin{corollary}\label{cor:big-lp-interior-smoothness}
Let $0\leq f\in L^p(X,\omega^n)$, $p>1$, be smooth and strictly positive
on $\Omega$. Every minimal-singularity solution of
$(\theta+\ddc\varphi)^n=cf\,\omega^n$, $c>0$, is smooth and strictly
$\theta$-plurisubharmonic on $\Omega$.
\end{corollary}
\begin{proof}
The analytic complement of $\Omega$ has zero mass for both measures,
so it suffices to know the equation in the bounded-potential sense on
$\Omega$. Apply \cref{thm:big-minimal-holder} with density $cf$ to obtain
local H\"older continuity. All hypotheses of
\cref{thm:nonclosed-interior-smoothness} are now satisfied.
\end{proof}

By \cref{thm:big-minimal-normalized}, such a solution exists, with a uniquely determined $c>0$. For a smooth closed big representative on a compact K\"ahler manifold,
\cref{thm:big-class-ample-smoothness} applies this conclusion to the
classical full-mass solution and identifies the regular open set with
the ample locus.


\section{Local density and boundary control}
\label{sec:local-density-boundary}

Let $X$ be a compact complex manifold of dimension $n\geq1$, let $\omega$ be a Hermitian form, and let $\theta$ be a smooth real $(1,1)$-form. Write $\mathrm{MA}_\theta(u)$ for the extended product of \cref{chap:products}. A function $u\in\PSH(X,\theta)$ has minimal ordinary singularities if every $v\in\PSH(X,\theta)$ satisfies $v\leq u+C_v$ for some constant $C_v$. We examine an arbitrary open set $\Omega\subset X$ on which the equation density is smooth and positive. The boundary criterion and examples below reconstruct \cite[Theorems~1.1--1.2 and Sections~3--5]{SuppliedLocalDensity26}.

The distinction from \cref{thm:nonclosed-interior-smoothness} is the domain: there the density is smooth on the entire regular locus of a fixed global strict barrier with analytic singularities, and the solution is locally H\"older on that locus. Local information on an arbitrary open subset does not supply these hypotheses.

\subsection{Smooth boundary traces}

Retain $X$, $\omega$ and $\theta$, and let $\mathrm{d}V$ be a smooth positive volume form. Fix $\rho\in\PSH(X,\theta)$ with $\theta+\ddc\rho\geq\varepsilon\omega$ for some $\varepsilon>0$. Let $\Omega$ be an open set where $\rho$ is smooth and $\mathrm{d}\theta=0$. For a coordinate ball $B\Subset\Omega$, a boundary trace $g$ means the full interior limit
\[
 \lim_{B\ni z\to\xi}u(z)=g(\xi)
 \qquad(\xi\in\partial B),
\]
not merely a radial or almost-everywhere limit.

\begin{theorem}\label{thm:smooth-boundary-trace-regularity}
Suppose that $u\in\PSH(X,\theta)$ has minimal ordinary singularities and
\[
 \mathrm{MA}_\theta(u)=f\,\mathrm{d}V\quad\text{on }\Omega,
 \qquad f\in C^\infty(\Omega),\quad f>0.
\]
Assume that $\Omega$ is covered by coordinate Euclidean balls $B_j\Subset\Omega$ such that $\theta$ has a smooth local potential near $\overline{B_j}$ and $u$ has a smooth boundary trace on $\partial B_j$. Then $u\in C^\infty(\Omega)$ and $\theta+\ddc u>0$ there.
\end{theorem}

\begin{proof}
Minimality gives $u\geq\rho-C$ for a constant $C$. Thus $u$ is locally bounded on $\Omega$. Fix one of the balls, write it as $B_R(z_0)$, and choose a smooth function $h$ near its closure with $\ddc h=\theta$. The bounded plurisubharmonic function $v=u+h$ has, by \cref{prop:bounded-products}, Bedford--Taylor measure $f\,\mathrm{d}V$ and smooth boundary trace $\varphi=g+h|_{\partial B}$.

Write $\mathrm{d}V=J\,\mathrm{d}\Lambda$ in these coordinates, where $\mathrm{d}\Lambda$ is Euclidean volume, and set $G=\pi^n fJ/n!$. Our normalization gives
\begin{equation}\label{eq:local-density-determinant}
 (\ddc V)^n=\frac{n!}{\pi^n}\det(V_{j\bar k})\,\mathrm{d}\Lambda
\end{equation}
for smooth $V$. There is a smooth strictly plurisubharmonic $V$ on $\overline B$ with determinant $G$ and boundary value $\varphi$. For $n\geq2$, this is the classical Dirichlet input \cite[Theorem~1.1, p.~1186]{GL10}, applied with zero background and flat Hermitian metric. Indeed, if $\Phi$ smoothly extends $\varphi$, then $\Phi+A(|z-z_0|^2-R^2)$ is a strict subsolution when $A$ is sufficiently large. Here the cited equation has background zero, metric $(\mathrm{i}/2)\partial\bar\partial|z|^2$, and right-hand side $G$. For $n=1$, extend $G$ smoothly with compact support and convolve with $\pi^{-1}\log|z|^2$; adding the harmonic extension of the remaining smooth boundary values gives $V$.

The difference $v-V$ tends to zero at every boundary point. Compactness of $\partial B$ makes this convergence uniform near the boundary: otherwise a sequence approaching the boundary with difference bounded away from zero would have a subsequence approaching one fixed boundary point. Bounded Dirichlet uniqueness \cite[Corollary~4.4, p.~17]{BT82} gives $v=V$. Therefore $u=V-h$ is smooth and its background form is positive on $B$. The balls cover $\Omega$.
\end{proof}

The proof uses minimality only to obtain the local lower bound from $\rho$. Consequently the same conclusion holds for any solution already known to be locally bounded on $\Omega$ and satisfying the same trace condition. Local closedness supplies the potentials after the coordinate balls are made small enough; the background may be nonclosed elsewhere.

\begin{example}\label{ex:logarithmic-continuous-density}
For every $n\geq1$, the trace criterion admits a bounded continuous solution whose global density belongs to $L^1$ and to no $L^p$ with $p>1$. The solution need not be H\"older at the complement of $\Omega$.

Let $E=\mathbb C/(\mathbb Z+\mathrm{i}\mathbb Z)$, $X=E^n$, and
\[
 \omega_j=\frac{\mathrm{i}}{2\pi}\,\mathrm{d}z_j\wedge\mathrm{d}\bar z_j,
 \qquad \theta=\sum_{j=1}^n\omega_j,
 \qquad D=\{0\}\times E^{n-1}.
\]
Take $\mathrm{d}V=\theta^n$ and $\rho=0$. Choose a nonnegative smooth cutoff $\chi$ supported in a coordinate disc $|z|<r_0<1$ of $E$ and equal to one near zero. For $\beta>0$, define
\[
 \psi(z)=\chi(z)(-\log|z|^2)^{-\beta}\quad(z\ne0),
 \qquad \psi(0)=0,
\]
and extend by zero. On the smaller punctured disc put $t=|z|^2$ and $L=-\log t$. Differentiation gives
\begin{equation}\label{eq:continuous-logarithmic-profile}
 \psi_{z\bar z}=\frac{\beta(\beta+1)}{tL^{\beta+2}},
 \qquad t\frac{\mathrm{d}}{\mathrm{d}t}L^{-\beta}
       =\beta L^{-\beta-1}\longrightarrow0.
\end{equation}
The first expression is integrable and the second is the radial flux. Integration by parts on punctured discs, using also $\psi\to0$, shows that the distributional Hessian has no atom at zero. On the cutoff annulus its negative part is bounded. Hence, for a sufficiently small fixed $\eta>0$, $u=\eta\psi(z_1)$ satisfies $\theta+\ddc u\geq\theta/2$.

Only one coordinate contributes to $\ddc u$. Locally approximate the bounded subharmonic potential in that coordinate from above and use \cref{thm:bt-input}; the smooth diagonal determinant identity passes to the limit and yields
\[
 \mathrm{MA}_\theta(u)=f\theta^n,
 \qquad f=1+\eta\psi_{z\bar z}(z_1).
\]
Give $f$ any positive value on the null set $D$. It is smooth and positive on $\Omega=X\setminus D$, where $u$ is smooth and has smooth traces on all relatively compact coordinate balls. Near $D$ its singular term is comparable to $|z_1|^{-2}(-\log|z_1|^2)^{-\beta-2}$. Its radial $L^1$ integral converges, whereas
\[
 \int_0^{r_0}s^{1-2p}(-\log s^2)^{-p(\beta+2)}\,\mathrm{d}s=+\infty
 \qquad(p>1).
\]
Also $(-\log|z|^2)^{-\beta}/|z|^\alpha\to+\infty$ for every $\alpha>0$, so $u$ is not H\"older at $D$. Finally, boundedness gives minimal ordinary singularities, since every global $\theta$-plurisubharmonic competitor is bounded above.
\end{example}

\subsection{A bounded discontinuous potential}

Let $U=W=\{|z|<1/4\}\subset\mathbb C$, with coordinates $z$ and $w$ on the two factors. Write $\mathrm{d}A$ for planar Lebesgue measure and set
\[
 a_k=\alpha_k=2^{-k-4},\qquad
 A=\sum_{k\geq1}\alpha_k=\frac1{16},\qquad
 S=\{0\}\cup\{a_k:k\geq1\}.
\]
Define on $U$
\begin{equation}\label{eq:accumulating-polar-weight}
 h(z)=\sum_{k\geq1}\alpha_k\log|z-a_k|^2,
 \qquad b(z)=\mathrm{e}^{h(z)},
\end{equation}
with $b=0$ where $h=-\infty$.

\begin{lemma}\label{lem:discontinuous-constant-measure}
The function
\[
 a(z)=\frac1\pi\int_U\log|z-\zeta|^2\,b(\zeta)^{-1}\,\mathrm{d}A(\zeta)
\]
is bounded and continuous on $U$. The bounded function
\[
 P(z,w)=a(z)+b(z)|w|^2
\]
is plurisubharmonic on $U\times W$, is smooth off $S\times W$, and satisfies
\begin{equation}\label{eq:discontinuous-constant-measure}
 (\ddc P)^2=\frac2{\pi^2}\,\mathrm{d}A(z)\,\mathrm{d}A(w).
\end{equation}
It is discontinuous at every $(0,w_0)$ with $w_0\ne0$.
\end{lemma}

\begin{proof}
The summands defining $h$ are negative on $U$, so their partial sums decrease. Since $h(0)=2\sum_k\alpha_k\log a_k> -\infty$, the limit is subharmonic. On compact subsets of $U\setminus S$, the series and its derivatives converge uniformly, making $h$ harmonic there. Moreover $0\leq b\leq1$, $b(a_k)=0$, and $b(0)>0$.

Weighted arithmetic--geometric mean gives, almost everywhere,
\begin{equation}\label{eq:polar-reciprocal-integrability}
 b(z)^{-2}\leq\sum_{k\geq1}\frac{\alpha_k}{A}|z-a_k|^{-4A}.
\end{equation}
For the countable inequality, first use the first $N$ weights and assign the remaining weight to the factor $1$; then let $N$ increase. Since $4A=1/4<2$, the integrals of the terms $|z-a_k|^{-4A}$ on $U$ are bounded uniformly in $k$. Thus $b^{-1}\in L^2(U)$.

The logarithmic kernel has uniformly bounded $L^2$ norm on the translates of $U$ that occur in the definition of $a$. Cauchy--Schwarz bounds $a$ uniformly. Choose a compactly supported $L^2(\mathbb C)$ kernel agreeing with $\log|z|^2$ on $U-U$. Translation continuity in $L^2$, paired with $b^{-1}\mathbf1_U\in L^2$, proves continuity of $a$. The distribution identity $\partial_z\partial_{\bar z}\log|z|^2=\pi\delta_0$ gives $a_{z\bar z}=b^{-1}$. In particular $a$ is subharmonic and, by local elliptic regularity, smooth off $S$.

The second summand of $P$ is $\exp(h(z)+\log|w|^2)$, a convex increasing composition of a plurisubharmonic function, extended by zero at its poles. Thus $P$ is bounded and plurisubharmonic. At smooth points the harmonicity of $h=\log b$ gives $bb_{z\bar z}=|b_z|^2$, and hence
\[
 \det(P_{j\bar k})
 =ba_{z\bar z}+|w|^2\bigl(bb_{z\bar z}-|b_z|^2\bigr)=1.
\]
The set $S\times W$ is contained in the pole set of $h(z)+\log|z|$, a plurisubharmonic function not identically $-\infty$. It is also Lebesgue-null. Pluripolar nullity in \cref{thm:bt-input} and \eqref{eq:local-density-determinant} prove \eqref{eq:discontinuous-constant-measure} throughout the bidisc. Finally,
\[
 P(a_k,w_0)=a(a_k)\longrightarrow a(0)
 <a(0)+b(0)|w_0|^2=P(0,w_0)
\]
when $w_0\ne0$.
\end{proof}

\subsection{Compact gluing}

Let $P$, $a$, $b$, $S$, $U$ and $W$ be as in \cref{lem:discontinuous-constant-measure}. We work on the flat complex torus $X=(\mathbb C/(\mathbb Z+\mathrm{i}\mathbb Z))^2$, identifying a small coordinate ball with a ball in $U\times W$. For a smooth function $v$ on the ball write $H_v=(v_{j\bar k})$ for its complex Hessian and $I$ for the two-dimensional identity matrix.

\begin{theorem}\label{thm:compact-local-density-discontinuity}
There are a positive flat K\"ahler form $\theta$ on $X$, a bounded $u\in\PSH(X,\theta)$ with minimal ordinary singularities, and a density $F>0$ almost everywhere such that
\[
 \mathrm{MA}_\theta(u)=F\theta^2,\qquad
 F\in L^1(X,\theta^2),\qquad
 F\notin L^p(X,\theta^2)\quad(p>1).
\]
On an open coordinate ball $\Omega$, $F$ is a positive constant and $u$ is discontinuous. The constant function $\rho=0$ is a smooth strict barrier on $X$.
\end{theorem}

\begin{proof}
Set $R=1/8$ and choose $M$ so that $P_0=P+M$ satisfies $0\leq P_0\leq L$ on $B_R\Subset U\times W$ for some finite $L$. Choose $C>0$ with
\begin{equation}\label{eq:compact-gluing-collar}
 \frac{3CR^2}{16}\geq2,\qquad
 \frac{5CR^2}{16}>L+2,
\end{equation}
and put $Q(Z)=C(|Z|^2-R^2/4)$ for $Z=(z,w)$. Let $\chi$ be the convolution of $|t|$ with an even nonnegative smooth probability kernel supported in $(-1,1)$. Then $\chi$ is smooth and convex, $|\chi'|\leq1$, $\chi'(0)=0$, and $\chi(t)=|t|$ for $|t|\geq1$. The function
\[
 m(s,t)=\frac{s+t+\chi(s-t)}2,\qquad v=m(P_0,Q)
\]
is bounded and plurisubharmonic, because $m$ is convex and nondecreasing in each variable. The collar inequalities imply
\[
 v=P_0\quad\text{on }B_{R/4},\qquad
 v=Q\quad\text{when }3R/4\leq|Z|<R.
\]
Indeed, $Q\leq-2\leq P_0-2$ on the first ball and $Q>L+2\geq P_0+2$ on the collar. Define
\[
 \theta=\frac{\mathrm{i}C}{2\pi}
  (\mathrm{d}z\wedge\mathrm{d}\bar z+\mathrm{d}w\wedge\mathrm{d}\bar w),
 \qquad
 u=\begin{cases}v-Q&\text{on }B_R,\\0&\text{outside }B_R.\end{cases}
\]
The collar identity makes $u$ identically zero near $\partial B_R$ and gives a global bounded potential with $\theta+\ddc u=\ddc v$ on the ball and $\theta+\ddc u=\theta$ outside. Thus $u$ is admissible and has minimal ordinary singularities.

Outside the pluripolar set $(S\times W)\cap B_R$, set $\lambda=(1+\chi'(P_0-Q))/2$. Differentiating the regularized maximum gives
\begin{equation}\label{eq:regularized-max-hessian-density}
 H_v=\lambda H_P+(1-\lambda)CI
   +\frac{\chi''(P_0-Q)}2
        \partial(P_0-Q)\otimes\overline{\partial(P_0-Q)}.
\end{equation}
Both $H_P$ and $CI$ are positive definite there; the final term is positive semidefinite. Hence the classical density is positive. Bounded products charge no pluripolar sets, so there is no further measure on the exceptional set. Their mass on compact $X$ is finite by \cref{prop:bounded-products}. This proves $\mathrm{MA}_\theta(u)=F\theta^2$ for an almost-everywhere positive $F\in L^1$. On $\Omega=B_{R/4}$, \eqref{eq:discontinuous-constant-measure} gives the representative $F=C^{-2}$. There $u=P+M-Q$ is discontinuous at $(0,w_0)$ for $0<|w_0|<R/4$.

We prove failure of every global $L^p$ directly in the transition region. For large $k$, set
\[
 m_k=a(a_k)+M\in[0,L],\qquad
 w_k=\bigl(R^2/4+m_k/C-|a_k|^2\bigr)^{1/2},
 \qquad Z_k=(a_k,w_k).
\]
Then $w_k>0$, $P_0(Z_k)=Q(Z_k)=m_k$, and \eqref{eq:compact-gluing-collar} gives $|Z_k|^2<9R^2/16$. Near the isolated pole $a_k$,
\[
 h(z)=\alpha_k\log|z-a_k|^2+H_k(z),
\]
where $H_k$ is smooth and harmonic. Consequently $b$ and $P_0$ are continuous near $Z_k$. There is a product neighborhood $V_k\Subset B_R$ of $Z_k$, with $|w|$ bounded below by a positive number, on which $1/4\leq\lambda\leq3/4$.

At smooth points of $V_k$, positivity and the two-dimensional determinant expansion give
\[
 \begin{split}
 \det H_v&\geq\det\bigl(\lambda H_P+(1-\lambda)CI\bigr)\\
 &\geq\lambda(1-\lambda)C\operatorname{tr}H_P
 \geq c_k b_{z\bar z},
 \end{split}
\]
where $c_k>0$ may depend on $k$ and $C$. Here $P_{z\bar z}=b^{-1}+|w|^2b_{z\bar z}$. With $\zeta=z-a_k$, shrink the first factor of $V_k$ so that
\[
 b_{z\bar z}=|\zeta|^{2\alpha_k}\mathrm{e}^{H_k}
       \left|\frac{\alpha_k}{\zeta}+\partial_zH_k\right|^2
       \geq c'_k|\zeta|^{2\alpha_k-2}
\]
for some $c'_k>0$. Since $F=C^{-2}\det H_v$, this lower bound implies $F\notin L^p$ whenever $p(1-\alpha_k)\geq1$: the $w$-disc has positive area and the radial integral in $z$ contains
\[
 \int_0^{\delta_k}r^{1-2p(1-\alpha_k)}\,\mathrm{d}r=+\infty.
\]
For each $p>1$ such a $k$ exists because $\alpha_k\downarrow0$.
\end{proof}

Capacity domination of first order cannot replace the trace hypothesis. Indeed, for a bounded plurisubharmonic function $v$ on a bounded domain $D\subset\mathbb C^n$, choose $M\geq\sup_Dv-\inf_Dv>0$. The function $q=(v-\inf_Dv)/M$ is an admissible test for the relative capacity
\[
 \operatorname{Cap}_D(E)
 =\sup_{q\in\PSH(D),\,0\leq q\leq1}\int_E(\ddc q)^n.
\]
Thus $(\ddc v)^n(E)\leq M^n\operatorname{Cap}_D(E)$ for every Borel $E\subset D$. Constant functions have zero measure. The discontinuous $P$ satisfies this bound and even has constant positive volume measure.

\begin{example}\label{ex:local-lipschitz-constant-measure}
On the bidisc in $\mathbb C^2$ with coordinates $(z,w)$, the function
\[
 H(z,w)=2|z|(1+|w|^2)
\]
is locally Lipschitz, plurisubharmonic and not $C^1$ along $\{z=0\}$. Plurisubharmonicity follows by writing its two summands as convex exponential compositions of $\log|z|$ and $\log|z|+\log|w|^2$. At $z\ne0$,
\[
 H_{z\bar z}=\frac{1+|w|^2}{2|z|},\qquad
 H_{w\bar w}=2|z|,\qquad
 |H_{z\bar w}|^2=|w|^2.
\]
Hence $\det(H_{j\bar k})=1$, and pluripolar nullity gives $(\ddc H)^2=2\pi^{-2}\,\mathrm{d}\Lambda$ on the entire bidisc.
\end{example}

This last example rules out a purely local implication from H\"older continuity and a smooth positive density to smoothness. It does not supply the global hypotheses on the regular locus in \cref{thm:nonclosed-interior-smoothness}. In \cref{thm:compact-local-density-discontinuity}, the solution lacks local H\"older continuity on the indicated ball and the density is rough elsewhere on $X$, which is the regular locus of the chosen barrier $\rho=0$. The trace criterion instead identifies the given bounded solution with a smooth Dirichlet solution on each ball.


\Needspace{12\baselineskip}
\section{Equations in a prescribed ordinary type}\label{sec:ordinary-type-equations}

Return to a semipositive big $\theta$. We now require the solution to have the ordinary type of an envelope fixed point $\varphi$ satisfying \eqref{eq:als-model-condition}. Fix $\varepsilon\in(0,1)$ for which $P_{\varepsilon\theta}(\varphi)$ is genuine. The smaller-background envelope yields a uniform bound for the difference from $\varphi$. A supersolution and the addition theorem then prove existence in this type.

\Needspace{20\baselineskip}
\begin{theorem}
\label{thm:als-prescribed-lp}
Let $\mu=f\,\mathrm{d}V$, where $0\leq f\in L^p(\mathrm{d}V)$, $p>1$, and $\mu(X)>0$. Then:
\begin{enumerate}
\item For every $\lambda>0$ there is a unique $u\in\PSH(X,\theta)$ with $u\sim\varphi$ such that
\[
 \mathrm{MA}_\theta(u)=\mathrm{e}^{\lambda u}\mu.
\]
\item There are $u\in\PSH(X,\theta)$ with $u\sim\varphi$ and a uniquely determined constant $c>0$ such that
\[
 \mathrm{MA}_\theta(u)=c\mu.
\]
\end{enumerate}

\end{theorem}

Only $c$ is asserted to be unique in the second part. Uniqueness of the potential up to a constant is false under these hypotheses: \cref{thm:degenerate-hermitian-nonuniqueness} already gives a counterexample in the minimal singularity type of a Hermitian background. Positive-exponent uniqueness does not imply zero-exponent uniqueness.

We first bound the density norm from below in the prescribed ordinary type.

\begin{lemma}\label{lem:relative-density-lower-bound}
Fix $q>1$. There is $b_q>0$ such that every $u\in\PSH(X,\theta)$ with $u\sim\varphi$ and
\[
 \mathrm{MA}^{\mathrm{ext}}_\theta(u)=F\,\mathrm{d}V,
 \qquad 0\leq F\in L^q(\mathrm{d}V),
\]
satisfies $\|F\|_q\geq b_q$. For fixed $\theta$ and $\varepsilon$, the constant $b_q$ is independent of $\varphi$.
\end{lemma}
\begin{proof}
Since $u\sim\varphi$, the smaller-background envelope $P_{\varepsilon\theta}(u)$ is not identically $-\infty$. If $F=0$, then \cref{thm:singular-strict-domination}, applied to $u$ and $u+1$ with coefficient zero, would give $u\geq u+1$; hence $N=\|F\|_q>0$.

Apply \cref{lem:big-lp-subsolution} in the semipositive big background $(1-\varepsilon)\theta$ to $F/N$, whose $L^q$ norm is one. It gives a bounded $s$ with $-1\leq s\leq0$ and a constant $b_q>0$, depending only on this fixed background and on $q$, such that
\[
 \mathrm{MA}_{(1-\varepsilon)\theta}(s)\geq(b_q/N)F\,\mathrm{d}V.
\]
Put $h=P_{\varepsilon\theta}(\varphi)+s$. This is a $\theta$-plurisubharmonic function with $h\preceq u$. Since $\theta_h$ is the sum of the positive currents $\varepsilon\theta+\ddc P_{\varepsilon\theta}(\varphi)$ and $(1-\varepsilon)\theta+\ddc s$, the same density bound holds for $\mathrm{MA}^{\mathrm{ext}}_\theta(h)$. To justify this, we truncate $P_{\varepsilon\theta}(\varphi)$ at finite levels, use positivity of bounded mixed products, and exhaust the set where $P_{\varepsilon\theta}(\varphi)$ is finite; its pole set has zero volume.

If $N<b_q$, we would have
\[
 \mathrm{MA}^{\mathrm{ext}}_\theta(u)
       \leq(N/b_q)\mathrm{MA}^{\mathrm{ext}}_\theta(h+t)
 \quad\text{for every }t\in\mathbb R.
\]
Since $N/b_q<1$, \cref{thm:singular-strict-domination} would then imply $u\geq h+t$ for every $t$, which is impossible at a point where both potentials are finite. Thus $N\geq b_q$. The construction of $s$ used no property of $\varphi$, which proves the last assertion.
\end{proof}

The lemma bounds the $L^q$ norm of the density from below; it gives no positive lower bound for the total mass $\int_XF\,\mathrm{d}V$.


The next estimate explains how the same singularity type is retained in the limit.

\Needspace{18\baselineskip}
\begin{theorem}\label{thm:relative-linfty}
Under the hypotheses on $(X,\theta,\varphi)$ above, suppose $u\in\PSH(X,\theta)$, $u\sim\varphi$, $\sup_Xu=0$, and
\[
 \mathrm{MA}_\theta(u)\leq f\,\mathrm{d}V,
 \qquad0\leq f\in L^p(\mathrm{d}V),\quad p>1.
\]
There is a constant $C$, depending on the fixed geometry, $\varphi$, $p$, and an upper bound for $\|f\|_{L^p}$, such that
\[
 \varphi-C\leq u\leq\varphi.
\]
\end{theorem}

\begin{proof}
The upper bound follows directly from the fixed-point property: $u\leq0$ and $u\preceq\varphi$, so $u$ is a competitor for $P_\theta[\varphi]$.

Put $\varphi_\varepsilon=P_{\varepsilon\theta}(\varphi)$. Since $\varphi_\varepsilon\leq\varphi\leq0$, uniform exponential integrability for normalized $u$ and exponential integrability for this fixed $\varphi_\varepsilon$ allow a choice of $\gamma>0$ and $1<q<p$ such that
\[
 g=\mathrm{e}^{-\gamma(\varphi_\varepsilon+u)}f
\]
has a uniform $L^q$ bound. More explicitly, choose $q<p$ first, apply H\"older with exponents $p/q$ and $p/(p-q)$, and then choose $\gamma$ small enough to control the exponential factor by Cauchy--Schwarz and the two exponential integrability estimates.

Apply \cref{thm:big-bounded-normalized} in background $(1-\varepsilon)\theta$ to the density $1+g$. We obtain a bounded $v$, with $\sup_Xv=0$, and $c>0$ such that
\[
 ((1-\varepsilon)\theta+\ddc v)^n=c(1+g)\,\mathrm{d}V.
\]
Both $\|v\|_\infty$ and $c^{-1}$ are uniformly bounded. The added constant one ensures a uniform positive root integral, even if no lower bound for the integral of $f$ was specified.

Put $h=\varphi_\varepsilon+v\leq\varphi$ and define
\[
 \chi(t)=\frac1{1-t}\quad(t\leq0),\qquad
 \psi=\varphi+\chi(h-\varphi).
\]
Then $0\leq\chi\leq1$, $0<\chi'\leq1$, and $\chi''>0$. Interpret the formula off the pole sets. For bounded ordered potentials, \cref{lem:bounded-convex-pair}, with $G(x,y)=x+\chi(y-x)$, gives
\[
 \theta+\ddc G(\varphi,h)
 \geq(1-\chi'(h-\varphi))\theta_\varphi
       +\chi'(h-\varphi)\theta_h\geq0.
\]
For unbounded potentials apply the bounded assertion to
$\varphi_k=\max(\varphi,-k)$ and $h_k=\max(h,-k)\leq\varphi_k$. Monotonicity of $G$ makes $G(\varphi_k,h_k)$ decrease. Its limit is $\theta$-plurisubharmonic and lies between $\varphi$ and $\varphi+1$, so it is genuine and equals $\psi$ as defined above. Locality on common finite levels and bounded product positivity give
\[
 \mathrm{MA}^{\mathrm{ext}}_\theta(\psi)
 \geq(\varphi-h+1)^{-2n}\mathrm{MA}^{\mathrm{ext}}_\theta(h)
 \geq c(\varphi-h+1)^{-2n}\mathrm{e}^{-\gamma\varphi_\varepsilon}
             \mathrm{e}^{-\gamma u}f\,\mathrm{d}V.
\]
Pole sets have zero smooth volume, so the last inequality extends by exhaustion to all Borel sets.

Write $T=-\varphi_\varepsilon\geq0$. If $C\geq1+\|v\|_\infty$, then
$\varphi-h+1\leq T+C$. Since
\[
 A:=\inf_{T\geq0}\frac{\mathrm{e}^{\gamma T}}{(T+C)^{2n}}>0,
\]
we obtain, with a uniform $d>0$,
\[
 \mathrm{MA}^{\mathrm{ext}}_\theta(\psi)
 \geq d \mathrm{e}^{-\gamma u}\mathrm{MA}^{\mathrm{ext}}_\theta(u).
\]
Choose $K=2+\gamma^{-1}|\log d|$. On $\{u<\psi-K\}$ we have $u<1-K$, hence
\[
 \mathrm{MA}^{\mathrm{ext}}_\theta(u)
 \leq d^{-1}\mathrm{e}^{\gamma(1-K)}
            \mathrm{MA}^{\mathrm{ext}}_\theta(\psi-K)
 \leq \mathrm{e}^{-\gamma}\mathrm{MA}^{\mathrm{ext}}_\theta(\psi-K).
\]
The two potentials have the same ordinary singularity type. Moreover $P_{\varepsilon\theta}(u)$ is genuine, since $u\geq\varphi-C_u$ for some constant $C_u$. Singular strict domination yields $u\geq\psi-K\geq\varphi-K$. All constants used to choose $K$ are uniform in the asserted family. The possibly nonuniform $C_u$ is used only to provide a genuine envelope, not in any estimate.
\end{proof}

The estimate controls $u-\varphi$, rather than $u$ itself. The smaller-background condition supplies $\varphi_\varepsilon$ and enables the final domination argument; the bounded auxiliary equation is now supplied by the preceding sections.

\begin{lemma}\label{lem:relative-lp-supersolution}
For nonzero $0\leq f\in L^p(\mathrm{d}V)$ and $\lambda>0$, there exists $w\in\PSH(X,\theta)$ with $w\sim\varphi$ such that
\[
 \mathrm{MA}^{\mathrm{ext}}_\theta(w)\leq \mathrm{e}^{\lambda w}f\,\mathrm{d}V.
\]
\end{lemma}
\begin{proof}
Let $v$ be the bounded solution of $\mathrm{MA}_\theta(v)=\mathrm{e}^{\lambda v}f\,\mathrm{d}V$ supplied by \cref{thm:big-bounded-exponential}, and write $M=\sup_Xv$, $a=v-M\leq0$. The rooftop $p=P_\theta(a,\varphi)$ satisfies
$\varphi-\|a\|_\infty\leq p\leq\varphi$, so $p\sim\varphi$. By the rooftop inequality,
\[
 \mathrm{MA}^{\mathrm{ext}}_\theta(p)
 \leq\mathbf1_{\{p=a\leq\varphi\}}\mathrm{MA}_\theta(a)
   +\mathbf1_{\{p=\varphi<a\}}\mathrm{MA}^{\mathrm{ext}}_\theta(\varphi).
\]
The second term vanishes: \cref{prop:fixed-point-finite} carries the measure of $\varphi$ on $\{\varphi=0\}$, while $\varphi<a\leq0$ excludes that set. On the first contact set, $a=p$ and
$\mathrm{MA}_\theta(a)=\mathrm{e}^{\lambda(a+M)}f\,\mathrm{d}V$. Thus
$\mathrm{MA}^{\mathrm{ext}}_\theta(p)\leq \mathrm{e}^{\lambda(p+M)}f\,\mathrm{d}V$. Take $w=p+M$.
\end{proof}

\subsection{Existence for a positive exponent}

We prove part (1) of \cref{thm:als-prescribed-lp} by maximizing a weighted mass. Among the potentials $z\sim\varphi$ whose weighted measure $\mathrm{e}^{-\lambda z}\mathrm{MA}_\theta(z)$ lies below $\mu=f\,\mathrm{d}V$, we find one of maximal weighted mass; \cref{thm:relative-linfty} provides the compactness needed for this. If its weighted measure were strictly smaller than $\mu$, the addition theorem would produce a potential of larger weighted mass.

\Needspace{12\baselineskip}
\begin{proof}[of \cref{thm:als-prescribed-lp}, part (1)]
\emph{Step 1: a maximizing sequence.} Fix $\lambda>0$. By \cref{lem:relative-lp-supersolution}, there is $w\sim\varphi$ such that
\[
 \nu_0:=\mathrm{e}^{-\lambda w}\mathrm{MA}^{\mathrm{ext}}_\theta(w)
 \leq\mu=f\,\mathrm{d}V.
\]
Consider all potentials $z\sim\varphi$ for which
\[
 \nu_0\leq\nu_z:=\mathrm{e}^{-\lambda z}
                   \mathrm{MA}^{\mathrm{ext}}_\theta(z)\leq\mu,
\]
and let $B$ be the supremum of their weighted masses $\nu_z(X)$. This family is nonempty, since it contains $w$, and $B\leq\mu(X)<+\infty$. Choose $u_j$ in the family with $\nu_{u_j}(X)\to B$. Since $\nu_0\leq\nu_{u_j}$, \cref{cor:singular-exponential-order} gives $u_j\leq w$.

\emph{Step 2: uniform bounds.} We need a lower bound for the suprema of the $u_j$. Put $\varphi_\varepsilon=P_{\varepsilon\theta}(\varphi)\leq0$, and let $\psi$ be the bounded solution of
\[
 ((1-\varepsilon)\theta+\ddc\psi)^n=\mathrm{e}^{\lambda\psi}f\,\mathrm{d}V
\]
given by \cref{thm:big-bounded-exponential}. For $h=\psi+\varphi_\varepsilon$, positivity of the splitting $\theta_h=(\varepsilon\theta+\ddc\varphi_\varepsilon)+((1-\varepsilon)\theta+\ddc\psi)$ gives
\[
 \mathrm{e}^{-\lambda h}\mathrm{MA}^{\mathrm{ext}}_\theta(h)
 \geq \mathrm{e}^{-\lambda\varphi_\varepsilon}f\,\mathrm{d}V
 \geq\mu\geq\nu_{u_j}.
\]
Here $h\preceq u_j$, because $h\preceq\varphi$ and $u_j\sim\varphi$, and each $u_j$ admits a smaller-background envelope because $u_j\sim\varphi$. \Cref{cor:singular-exponential-order} therefore gives $u_j\geq h$. Consequently
\[
 \sup_Xh\leq M_j:=\sup_Xu_j\leq\sup_Xw,
\]
and the lower bound is finite.

We now apply \cref{thm:relative-linfty} to $v_j=u_j-M_j$. Their measures satisfy
\[
 \mathrm{MA}^{\mathrm{ext}}_\theta(v_j)
 =\mathrm{MA}^{\mathrm{ext}}_\theta(u_j)
 \leq \mathrm{e}^{\lambda\sup_Xw}f\,\mathrm{d}V,
\]
so $\varphi-C\leq v_j\leq\varphi$. Together with the bounds for $M_j$, this shows that all $u_j$ lie within a fixed bounded distance of $\varphi$.

\emph{Step 3: the limit.} Extract an $L^1$ limit $u$. By \cref{cor:common-barrier-continuity}, applied with the dominating measure $\mathrm{e}^{\lambda\sup_Xw}\mu$, the convergence holds in capacity, $u\sim\varphi$, and
\[
 \mathrm{MA}^{\mathrm{ext}}_\theta(u_j)
 \rightharpoonup\mathrm{MA}^{\mathrm{ext}}_\theta(u).
\]
After a further extraction we may also assume $\nu_{u_j}\rightharpoonup\nu$. Since $\nu_0\leq\nu_{u_j}\leq\mu$, the limit satisfies the same inequalities. We claim that
\[
 \mathrm{e}^{\lambda u_j}\nu_{u_j}\rightharpoonup \mathrm{e}^{\lambda u}\nu.
\]
Indeed, $\mathrm{e}^{\lambda u_j}\to \mathrm{e}^{\lambda u}$ in $L^1(\mu)$ by \cref{lem:fixed-npp-integrals}, and all the measures $\nu_{u_j}$ are at most $\mu$; this controls the error caused by changing the weight. The fixed bounded weight $\mathrm{e}^{\lambda u}$ can be approximated in $L^1(\mu)$ by continuous functions, and since both $\nu_{u_j}$ and $\nu$ are dominated by $\mu$, this approximation is uniform in the integrals. The weak convergence $\nu_{u_j}\rightharpoonup\nu$ then gives the claim, and the same argument applies after multiplication by any continuous test function.

\Needspace{7\baselineskip}
Since $\mathrm{e}^{\lambda u_j}\nu_{u_j}=\mathrm{MA}^{\mathrm{ext}}_\theta(u_j)$, it follows that
\[
 \mathrm{MA}^{\mathrm{ext}}_\theta(u)=\mathrm{e}^{\lambda u}\nu,
 \qquad \mathrm{e}^{-\lambda u}\mathrm{MA}^{\mathrm{ext}}_\theta(u)=\nu,
 \qquad\nu(X)=B.
\]
The second identity holds off the pole set of $u$, which has zero $\mu$-measure. Hence $u$ belongs to the family and attains the supremum $B$. Write $\nu=g\,\mathrm{d}V$, with $0\leq g\leq f$.

\emph{Step 4: no deficit.} Suppose that $g\neq f$ in $L^1(\mathrm{d}V)$. Applying \cref{lem:relative-lp-supersolution} to $f-g$, we find $v\sim\varphi$ with
\[
 \nu_v:=\mathrm{e}^{-\lambda v}\mathrm{MA}^{\mathrm{ext}}_\theta(v)
 \leq(f-g)\,\mathrm{d}V.
\]
This measure is nonzero. In fact, no potential $v$ admitting a smaller-background envelope has zero Monge--Amp\`ere measure, since \cref{thm:singular-strict-domination}, applied with coefficient zero to $v$ and $v+1$, would give $v\geq v+1$. Thus $\nu_v(X)>0$.

The sum $\nu+\nu_v$ is finite, so \cref{thm:weighted-measure-addition} gives $a\sim\varphi$ with
\[
 \mathrm{e}^{-\lambda a}\mathrm{MA}^{\mathrm{ext}}_\theta(a)=\nu+\nu_v.
\]
This measure still lies between $\nu_0$ and $\mu$, but its mass is $B+\nu_v(X)>B$, a contradiction. Therefore $g=f$, and $u$ solves the required equation. Uniqueness in the prescribed ordinary type follows from \cref{cor:singular-exponential-order}.
\end{proof}

\subsection{The normalized equation}

For the normalized equation we let the positive exponent tend to zero. Bounded comparison controls the normalizing constants from above, while the lower density estimate of \cref{lem:relative-density-lower-bound} keeps them away from zero. The common barriers then give convergence of the Monge--Amp\`ere measures and keep the limit in the prescribed type.

\Needspace{8\baselineskip}
\begin{proof}[of \cref{thm:als-prescribed-lp}, part (2)]
For each $j\geq1$, part (1) gives $u_j\sim\varphi$ with
\[
 \mathrm{MA}^{\mathrm{ext}}_\theta(u_j)=\mathrm{e}^{u_j/j}f\,\mathrm{d}V.
\]
Write $M_j=\sup_Xu_j$, $v_j=u_j-M_j$ and $c_j=\mathrm{e}^{M_j/j}$. Let $h$ be a bounded normalized solution of $\mathrm{MA}_\theta(h)=\alpha f\,\mathrm{d}V$, with $\alpha>0$, given by \cref{thm:big-bounded-normalized}.

\Needspace{5\baselineskip}
For $H_j=h+\|h\|_\infty+j\log\alpha$ we have
\[
 \mathrm{e}^{-H_j/j}\mathrm{MA}_\theta(H_j)
 =\mathrm{e}^{-(h+\|h\|_\infty)/j}f\,\mathrm{d}V
 \leq f\,\mathrm{d}V
 =\mathrm{e}^{-u_j/j}\mathrm{MA}^{\mathrm{ext}}_\theta(u_j).
\]
The bounded function $H_j$ admits a smaller-background envelope, and $u_j\preceq H_j$, so \cref{cor:singular-exponential-order} gives $u_j\leq H_j$. In particular,
\[
 0<c_j\leq\alpha\exp(2\|h\|_\infty/j).
\]

The lower density bound of \cref{lem:relative-density-lower-bound}, applied to $v_j$ with $q=p$, gives
\[
 b_p\leq\|c_j\mathrm{e}^{v_j/j}f\|_p\leq c_j\|f\|_p,
\]
so the $c_j$ are also bounded away from zero.

Since $\sup_Xv_j=0$ and $\mathrm{MA}^{\mathrm{ext}}_\theta(v_j)=c_j\mathrm{e}^{v_j/j}f\,\mathrm{d}V\leq C f\,\mathrm{d}V$, \cref{thm:relative-linfty} gives $\varphi-C'\leq v_j\leq\varphi$. Extract a subsequence with $v_j\to v$ in $L^1(X)$ and $c_j\to c>0$. By \cref{cor:common-barrier-continuity}, the convergence holds in capacity, $v\sim\varphi$, and the Monge--Amp\`ere measures converge weakly. Meanwhile, \cref{lem:normalized-exponential-lp} gives
\[
 \int_X|\mathrm{e}^{v_j/j}-1|f\,\mathrm{d}V\longrightarrow0.
\]
Thus the right-hand sides converge in total variation to $cf\,\mathrm{d}V$, and $\mathrm{MA}^{\mathrm{ext}}_\theta(v)=cf\,\mathrm{d}V$.

For the uniqueness of the constant, suppose that $a,b\sim\varphi$ solve the normalized equation with constants $c_a<c_b$. For every real $t$,
\[
 \mathrm{MA}^{\mathrm{ext}}_\theta(a)
 =\frac{c_a}{c_b}\mathrm{MA}^{\mathrm{ext}}_\theta(b+t),
\]
so \cref{thm:singular-strict-domination} gives $a\geq b+t$. This is impossible as $t\to+\infty$, at any point where both potentials are finite. Interchanging the roles of the two constants proves that $c$ is unique.
\end{proof}


\begin{lemma}\label{lem:bounded-measure-carrier}
Let $\theta\geq0$ be smooth and big, and let $\mu$ be a finite non-pluripolar measure on $X$. There are a bounded $\psi\in\PSH(X,\theta)$ and a nonnegative
$f\in L^1(\mathrm{MA}_\theta(\psi))$ such that
\[
 \mu=f\,\mathrm{MA}_\theta(\psi).
\]

\end{lemma}
\begin{proof}
Apply \cref{prop:relative-measure-carrier} with $\varphi=0$. Semipositivity gives $P_{\epsilon\theta}(0)=0$ for every $0<\epsilon<1$, and $\psi\sim0$ means that $\psi$ is bounded.
\end{proof}

These representations give an $L^1$ density, whereas \cref{thm:bounded-subsolution} requires a bounded density against the measure of a bounded potential. For an unbounded prescribed type, the representing potential is itself unbounded. The general-measure existence argument in \cref{sec:general-measure-relative} uses a bounded Hermitian reference and relative compactness to pass to the equation.


\section{Representing measures in an ordinary type}
\label{sec:relative-measure-carrier}\label{sec:bounded-measure-carrier}

We return to the general setting of the chapter, with the Hermitian form $\omega$ fixed at its beginning. Locally, every finite non-pluripolar measure is absolutely continuous with respect to the Monge--Amp\`ere measure of a bounded plurisubharmonic function (\cref{lem:local-bounded-measure-carrier}). We globalize this statement in a prescribed ordinary type. The local potentials are glued to a strictly $\theta$-plurisubharmonic potential lying below the reference and are then truncated relative to the reference. A decreasing convex sum of the resulting countable family has a Monge--Amp\`ere measure that dominates every local piece in the sense of absolute continuity. Its potential has the prescribed ordinary type, and the given measure has an integrable density with respect to its Monge--Amp\`ere measure.

\Needspace{13\baselineskip}
\begin{proposition}\label{prop:relative-measure-carrier}\index{Non-pluripolar measure!ordinary-type representation}
Let $\theta$ be a smooth big real $(1,1)$-form, and let $\varphi\in\PSH(X,\theta)$ be such that the envelope $P_{\epsilon\theta}(\varphi)$ of \cref{sec:obstacle-envelopes} is not identically $-\infty$ for some $0<\epsilon<1$. For every finite non-pluripolar measure $\mu$ there are
\[
 \psi\in\PSH(X,\theta),\qquad\psi\sim\varphi,
 \qquad 0\leq f\in L^1\bigl(\mathrm{MA}^{\mathrm{ext}}_\theta(\psi)\bigr)
\]
such that
\[
 \mu=f\,\mathrm{MA}^{\mathrm{ext}}_\theta(\psi).
\]
\end{proposition}
\begin{proof}
Translate $\varphi$ below zero. Put $q=P_{\epsilon\theta}(\varphi)\leq\varphi$, and choose a strict barrier $\rho\leq0$ with $\theta+\ddc\rho\geq a\omega$, where $a>0$. Then
\[
 s=q+(1-\epsilon)\rho\leq\varphi,\qquad
 \theta+\ddc s\geq(1-\epsilon)a\omega.
\]
In particular $s$ is $\theta$-plurisubharmonic and not identically $-\infty$. Its pole set is pluripolar, as is that of $\varphi$, so both are negligible for $\mu$.

Cover $X$ by finitely many closed coordinate balls $K_i$, each contained in a larger coordinate ball $B_i$. Apply \cref{lem:local-bounded-measure-carrier} to $\mu_i=\mathbf1_{K_i}\mu$. This gives bounded local plurisubharmonic $v_i$, normalized by $1\leq v_i\leq2$, with
\[
 \mu_i\ll\nu_i:=(\ddc v_i)^n.
\]
Take a smooth cutoff $\chi_i$ which is zero near $K_i$ and $-3$ on a boundary collar of $B_i$. Choose $\delta_i>0$ so small that
\[
 \theta+\ddc\bigl(s+\delta_i(v_i+\chi_i)\bigr)
 \geq\tfrac12(1-\epsilon)a\omega+\delta_i\ddc v_i.
\]
Let $S_i$ equal $\max(s,s+\delta_i(v_i+\chi_i))$ on $B_i$ and $s$ elsewhere. The boundary collar makes this a global $\theta$-plurisubharmonic function. Since $\chi_i\leq0$ may be arranged,
\[
 s\leq S_i\leq s+2\delta_i\leq\varphi+2\delta_i.
\]
\Needspace{7\baselineskip}
For every positive integer $m$, set
\[
 q_{i,m}=\max(\varphi-m,S_i)-2\delta_i.
\]
These potentials have the same ordinary type as $\varphi$, and
\begin{equation}\label{eq:relative-carrier-bounds}
 \varphi-m-2\delta_i\leq q_{i,m}\leq\varphi.
\end{equation}

On a neighborhood of $K_i$, $S_i=s+\delta_i v_i$ away from the pole set of $s$. Let
\[
 D_{i,m}=\{s+\delta_i v_i>\varphi-m\}
\]
in that neighborhood, restricting to points where the potentials are finite. This is a plurifinely open set there. Locality and positive mixed products give
\begin{equation}\label{eq:relative-carrier-local-domination}
 \mathbf1_{D_{i,m}}\mathrm{MA}^{\mathrm{ext}}_\theta(q_{i,m})
 \geq\delta_i^n\mathbf1_{D_{i,m}}\nu_i.
\end{equation}
Here $s$ need not be bounded. After a local smooth correction, the positive product inequality holds at every common finite truncation level of $s$ and the bounded potential $v_i$. We then exhaust the complement of the pole set of $s$; the non-pluripolar measure $\nu_i$ loses no mass in this exhaustion. Restricting the resulting inequality by plurifine locality proves \eqref{eq:relative-carrier-local-domination} for the extended measures.

For $\mu_i$-almost every point of $K_i$, both $s$ and $\varphi$ are finite, and hence that point belongs to $D_{i,m}$ for some $m$. Consequently, a Borel set which is null for every $\mathrm{MA}^{\mathrm{ext}}_\theta(q_{i,m})$ is null for every $\mu_i$, and therefore for $\mu$.

Enumerate the countable family $q_{i,m}$ as $q_k$, and write its bounds as
$\varphi-C_k\leq q_k\leq\varphi$ with $C_k<\infty$. Choose positive $t_k$ such that
\[
 \sum_k t_k<1,\qquad B:=\sum_k t_kC_k<\infty.
\]
The finite convex combinations
\[
 \psi_N=\sum_{k=1}^N t_kq_k
               +\left(1-\sum_{k=1}^N t_k\right)\varphi
\]
decrease and lie between $\varphi-B$ and $\varphi$. Their limit $\psi$ is $\theta$-plurisubharmonic, satisfies $\psi\geq\varphi-B$, and hence $\psi\sim\varphi$. For any fixed $k$, separate its summand to write
\[
 \psi=t_kq_k+(1-t_k)w_k,
 \qquad w_k\in\PSH(X,\theta),\qquad w_k\sim\varphi.
\]
The remainder is obtained by the same decreasing convex-sum construction. Positive non-pluripolar mixed products give
\[
 \mathrm{MA}^{\mathrm{ext}}_\theta(\psi)
       \geq t_k^n\mathrm{MA}^{\mathrm{ext}}_\theta(q_k).
\]
This is valid also for infinite measures: first restrict away from the common pole set, which is the pole set of $\varphi$ by the ordinary-type bounds, and then exhaust by finite truncations. All measures ignore that pole set. The preceding null-set conclusion now gives
$\mu\ll\mathrm{MA}^{\mathrm{ext}}_\theta(\psi)$.

\Needspace{10\baselineskip}
By \cref{thm:arbitrary-extended-ma}, the measure $\mathrm{MA}^{\mathrm{ext}}_\theta(\psi)$, which may have infinite mass, is $\sigma$-finite. The Radon--Nikodym theorem applies to this $\sigma$-finite measure and the finite $\mu$. It gives a nonnegative density $f$ with
\[
 \int_X f\,\mathrm{MA}^{\mathrm{ext}}_\theta(\psi)=\mu(X)<\infty,
\]
as claimed.
\end{proof}


\section{Lower mass and solvability for all measures}
\label{sec:lower-mass-solvability}

Throughout the rest of the chapter, $h$ denotes a Hermitian form on $X$, used as the background; the sections that use another background say so explicitly. No singularity type is prescribed: the unknown potentials range over all of $\PSH(X,h)$. Constants belong to $\PSH(X,h)$, and every $u\in\PSH(X,h)$ is bounded above, so $u\preceq v$ for every bounded $v\in\PSH(X,h)$.

The masses of bounded potentials are measured by the lower and upper masses $v_\pm(h)$ of \cref{def:bmp}, the infimum and the supremum of $\int_X(h+\ddc v)^n$ over the bounded $v\in\PSH(X,h)$; recall that $h$ has positive lower mass if $v_-(h)>0$. If $h$ is K\"ahler, Stokes' theorem gives $v_-(h)=v_+(h)=\int_Xh^n$; on a non-K\"ahler manifold the lower mass can vanish (\cref{thm:vanishing-lower-mass}).

For a nonzero nonnegative $q\in L^p(\mathrm{d}V)$ with $p>1$, \cref{thm:kn-continuous} gives a constant $c>0$ and a continuous $u\in\PSH(X,h)$ with $\mathrm{MA}_h(u)=cq\,\mathrm{d}V$, and \cref{prop:vanishing-exponent} shows that $c$ is determined by $q$. We call $c_h(q)=c$ the \emph{normalization scalar}\index{Normalization scalar!definition} of $q$, and for the measure $\nu=q\,\mathrm{d}V$ we also write $c_h(\nu)=c_h(q)$. We say that a finite measure $\nu$ admits a \emph{positive scaled subsolution} if there are $u\in\PSH(X,h)$ and a constant $a>0$ such that $\mathrm{MA}^{\mathrm{ext}}_h(u)\geq a\nu$.

If $\nu$ admits a solution of $\mathrm{MA}^{\mathrm{ext}}_h(u)=c\nu$ with $c>0$, then it admits in particular a positive scaled subsolution. We show that, when $h$ has zero lower mass, some probability density admits no positive scaled subsolution. Conversely, if $v_-(h)>0$, every nonzero finite non-pluripolar measure admits such a solution (\cref{cor:positive-lower-mass-solvability} in \cref{chap:relative-finite-measures}). The key estimate is a consequence of strict domination: if $\mathrm{MA}^{\mathrm{ext}}_h(u)\geq aq\,\mathrm{d}V$, then $a\leq c_h(q)$. When $v_-(h)=0$, there are smooth probability densities with arbitrarily small normalization scalars, and a weighted sum of such densities, with scalars decreasing fast enough, is a single probability density admitting no positive scaled subsolution.

\begin{lemma}\label{lem:bounded-normalizer-subsolution}\label{lem:smooth-normalizer-subsolution-bound}
Let $0\leq q\in L^p(\mathrm{d}V)$, $p>1$, be nonzero. If $u\in\PSH(X,h)$ and $a>0$ satisfy
\[
 \mathrm{MA}^{\mathrm{ext}}_h(u)\geq aq\,\mathrm{d}V,
\]
then $a\leq c_h(q)$. Consequently, for nonzero nonnegative $q,r\in L^p(\mathrm{d}V)$,
\[
 q\leq r\ \Longrightarrow\ c_h(q)\geq c_h(r),
 \qquad c_h(tq)=t^{-1}c_h(q)\quad(t>0).
\]
\end{lemma}
\begin{proof}
Let $v$ be a continuous solution of $\mathrm{MA}_h(v)=c_h(q)q\,\mathrm{d}V$, and suppose, to reach a contradiction, that $a>c_h(q)$. Then $\tau=c_h(q)/a<1$ and, for every real $A$,
\[
 \mathrm{MA}_h(v+A)=c_h(q)q\,\mathrm{d}V
       \leq\tau\,\mathrm{MA}^{\mathrm{ext}}_h(u).
\]
The translate $v+A$ is bounded, so $P_{h/2}(v+A)\geq\inf_Xv+A$ is not identically $-\infty$, and $u\preceq v+A$. \Cref{thm:singular-strict-domination}, applied with $\delta=1/2$ and $c=\tau$ to $v+A$ in place of $u$ and to $u$ in place of $v$, therefore gives $v+A\geq u$ for every real $A$. Letting $A\to-\infty$ at a point where $u$ is finite gives a contradiction. Hence $a\leq c_h(q)$.

If $q\leq r$, a continuous solution of $\mathrm{MA}_h(w)=c_h(r)r\,\mathrm{d}V$ satisfies $\mathrm{MA}_h(w)\geq c_h(r)q\,\mathrm{d}V$, so the first assertion gives $c_h(r)\leq c_h(q)$. Finally, the equation $\mathrm{MA}_h(v)=c_h(q)q\,\mathrm{d}V$ can be written as $\mathrm{MA}_h(v)=t^{-1}c_h(q)\,(tq)\,\mathrm{d}V$, and uniqueness of the scalar gives $c_h(tq)=t^{-1}c_h(q)$.
\end{proof}

\begin{lemma}\label{lem:lp-normalizer-continuity}
Fix $p>1$. If nonzero nonnegative densities $q_j,q\in L^p(\mathrm{d}V)$ satisfy $q_j\to q$ in $L^p$, then $c_h(q_j)\to c_h(q)$.
\end{lemma}
\begin{proof}
We first show that the scalars $c_h(q_j)$ lie in a compact subset of $(0,+\infty)$. Fix a Gauduchon metric $\gamma$ on $X$ (\cref{def:gauduchon-metric}; in dimension one, any Hermitian metric). Pairing the equation for $q_j$ with $\gamma^{n-1}$, as in \cref{prop:big-singular-gauduchon-trace} with $\theta=h$ and barrier $\rho=0$, gives
\[
 c_h(q_j)^{1/n}\int_X\left(\frac{q_j\,\mathrm{d}V}{\gamma^n}\right)^{1/n}\gamma^n
 \leq\int_Xh\wedge\gamma^{n-1}.
\]
Since $q_j^{1/n}\to q^{1/n}$ in $L^n$, the integrals on the left converge to the corresponding integral for $q$, which is positive. Hence the scalars are bounded above. The lower-density estimate of \cref{thm:lp-potential-estimate} gives $c_h(q_j)\|q_j\|_{L^p}\geq c_p>0$, so they are also bounded away from zero.

Consequently, the normalized continuous solutions $u_j$ of $\mathrm{MA}_h(u_j)=c_h(q_j)q_j\,\mathrm{d}V$ have uniformly bounded $L^p$ densities. By \cref{thm:lp-uniform-compactness}, every subsequence has a further subsequence along which $u_j$ converges uniformly to some $u$ and $c_h(q_j)$ converges to some $c>0$. Continuity of bounded Monge--Amp\`ere products under uniform convergence gives $\mathrm{MA}_h(u)=cq\,\mathrm{d}V$, and uniqueness of the scalar gives $c=c_h(q)$. Since every subsequence of $(c_h(q_j))$ has a further subsequence converging to $c_h(q)$, the whole sequence converges to $c_h(q)$.
\end{proof}

\begin{theorem}\label{thm:lower-mass-obstructing-measure}\index{Non-pluripolar measure!solvability obstruction}
Suppose $v_-(h)=0$. There is a probability measure $\mu=f\,\mathrm{d}V$, with $f\in L^1(\mathrm{d}V)$ and $f\geq m$ almost everywhere for a constant $m>0$, such that no $u\in\PSH(X,h)$ and $a>0$ satisfy
\begin{equation}\label{eq:lower-mass-impossible-subsolution}
 \mathrm{MA}^{\mathrm{ext}}_h(u)\geq a\mu.
\end{equation}
In particular, the equation $\mathrm{MA}^{\mathrm{ext}}_h(u)=c\mu$ has no solution $u\in\PSH(X,h)$ for any $c>0$; neither a singularity type nor finiteness of the mass of $u$ is imposed.
\end{theorem}
\begin{proof}
By \cref{thm:gaud-volume-mass-extrema,thm:gaud-volume-dual-formula}, $v_-(h)$ is the infimum of the normalization scalars $c_h(\nu)$ over smooth strictly positive probability measures $\nu$. Since $v_-(h)=0$, we can choose such measures $\nu_j=f_j\,\mathrm{d}V$ with
\[
 c_h(\nu_j)\leq 2^{-2j}\qquad(j\geq1).
\]
Set
\[
 \mu=\sum_{j=1}^{\infty}2^{-j}\nu_j,
 \qquad f=\sum_{j=1}^{\infty}2^{-j}f_j.
\]
By monotone convergence, $\int_Xf\,\mathrm{d}V=1$, so $f\in L^1$. Moreover $f\geq\tfrac12\min_X f_1=:m>0$, and $\mu$ is non-pluripolar, being absolutely continuous with respect to $\mathrm{d}V$.

Suppose that \eqref{eq:lower-mass-impossible-subsolution} held for some $u$ and $a>0$. Since $\mu\geq2^{-j}\nu_j$, we would have
\[
 \mathrm{MA}^{\mathrm{ext}}_h(u)\geq a2^{-j}\nu_j
\]
for every $j$, and \cref{lem:bounded-normalizer-subsolution} would give $a2^{-j}\leq c_h(\nu_j)\leq2^{-2j}$. Thus $a\leq2^{-j}$ for every $j$, contradicting $a>0$.
\end{proof}

\Needspace{12\baselineskip}
\begin{corollary}\label{cor:lower-mass-lone-instability}
Assume $v_-(h)=0$, and let $\mu=f\,\mathrm{d}V$ be the measure of \cref{thm:lower-mass-obstructing-measure}. Then $f\notin L^p(\mathrm{d}V)$ for every $p>1$. Moreover, for every probability density $g\in L^1(\mathrm{d}V)$ and every $0<t<1$, the probability density
\[
 g_t=(1-t)g+tf
\]
admits no positive scaled subsolution, although $\|g_t-g\|_{L^1}\leq2t$. In particular, the probability densities admitting no positive scaled subsolution are dense among all probability densities in $L^1(\mathrm{d}V)$.
\end{corollary}
\begin{proof}
If $f\in L^p$ for some $p>1$, \cref{thm:kn-continuous} would give a continuous solution of $\mathrm{MA}_h(u)=cf\,\mathrm{d}V$ with $c>0$, contradicting the theorem. Since $g_t\,\mathrm{d}V\geq t\mu$, a positive scaled subsolution for $g_t\,\mathrm{d}V$ would also give one for $\mu$. Finally,
$\|g_t-g\|_{L^1}=t\|f-g\|_{L^1}\leq2t$.
\end{proof}

Smooth strictly positive probability densities are also dense in $L^1(\mathrm{d}V)$, and each of them admits a smooth solution of the normalized equation. Thus, when the lower mass vanishes, solvability is not preserved under $L^1$ approximation. The obstruction comes from the way the density concentrates, not from its size: the density in \cref{thm:lower-mass-obstructing-measure} is bounded below by a positive constant. It occurs, for instance, on the compact threefold of \cref{thm:vanishing-lower-mass}, where every Hermitian form has zero lower mass.

\Cref{thm:lower-mass-obstructing-measure} does not exclude unbounded solutions of the exponential equation $\mathrm{MA}^{\mathrm{ext}}_h(u)=\mathrm{e}^{\lambda u}\mu$ with $\lambda>0$, because the factor $\mathrm{e}^{\lambda u}$ need not have a positive lower bound. It does exclude bounded solutions, since a bounded solution would satisfy $\mathrm{MA}^{\mathrm{ext}}_h(u)\geq\mathrm{e}^{\lambda\inf_Xu}\mu$. On suitable products, \cref{thm:zero-threshold-solvable-concentration} chooses the concentrating components of such a density so that its exponential equation has a solution for every positive exponent. We do not know whether this holds for every density admitting no positive scaled subsolution, nor which individual measures admit a solution of the normalized equation when $v_-(h)=0$.


\section{Subsolutions for an individual density}
\label{sec:individual-density-subsolutions}

For one fixed density, the normalizers of its bounded truncations give an exact test for the existence of a positive scaled subsolution. A cluster limit attains the threshold because density lower bounds survive $L^1$ convergence of potentials through upper envelopes.

Let $0\leq f\in L^1(\mathrm{d}V)$ have positive integral, and put
\begin{equation}\label{eq:individual-truncated-normalizer}
 f_k=\min(f,k),\qquad
 b_h(f)=\lim_{k\to\infty}c_h(f_k).
\end{equation}
The limit exists in $[0,+\infty)$ because the constants decrease and $c_h(f_1)$ is finite. The order lemma also gives
\begin{equation}\label{eq:individual-bounded-tests}
 b_h(f)=\inf\{c_h(q):0\leq q\leq f,\ q\in L^\infty,
                                      \ \textstyle\int_Xq\,\mathrm{d}V>0\}.
\end{equation}
In fact every such $q$ is bounded above by $f_k$ for sufficiently large $k$, and the $f_k$ themselves are admissible tests.

A lower bound for a Monge--Amp\`ere density survives $L^1$ convergence of potentials. This uses upper envelopes and damped continuity, rather than continuity of the original Monge--Amp\`ere measures.

\Needspace{8\baselineskip}
\begin{lemma}\label{lem:density-subsolution-lone-closure}
Let $s_j,s\in\PSH(X,h)$ have a common upper bound and satisfy $s_j\to s$ in $L^1(X)$. Suppose $g_j,g\geq0$, $g_j\to g$ in $L^1(\mathrm{d}V)$, and
\[
 \mathrm{MA}^{\mathrm{ext}}_h(s_j)\geq g_j\,\mathrm{d}V.
\]
Then $\mathrm{MA}^{\mathrm{ext}}_h(s)\geq g\,\mathrm{d}V$.
\end{lemma}
\begin{proof}
Translate the potentials to be nonpositive and extract a subsequence for which both convergences hold almost everywhere. Set
\[
 a_j=\inf_{\ell\geq j}g_\ell,\qquad
 w_j=\left(\sup_{\ell\geq j}s_\ell\right)^*.
\]
Then $0\leq a_j\uparrow g$ almost everywhere and $w_j\downarrow s$. For fixed $j$, the finite maxima $\max_{j\leq\ell\leq N}s_\ell$ have measures at least $a_j\,\mathrm{d}V$ by ordered contact. They increase to $w_j$ after regularization, with the genuine lower bound $s_j$ and upper bound zero. The convergence holds in capacity by \cref{lem:monotone-psh-capacity}. Damped continuity therefore gives
\[
 \mathrm{e}^{w_j}\mathrm{MA}^{\mathrm{ext}}_h(w_j)
       \geq \mathrm{e}^{w_j}a_j\,\mathrm{d}V.
\]
Indeed the right sides for the finite maxima converge by dominated convergence, since $a_j\leq g$ almost everywhere. Divide by the positive weight off the pole set to obtain $\mathrm{MA}^{\mathrm{ext}}_h(w_j)\geq a_j\,\mathrm{d}V$.

Now $w_j\downarrow s$ is again convergence in capacity with a genuine lower bound. Apply damped continuity once more. Since $0\leq\mathrm{e}^{w_j}a_j\leq g$, dominated convergence gives
\[
 \mathrm{e}^s\mathrm{MA}^{\mathrm{ext}}_h(s)
       \geq\mathrm{e}^s g\,\mathrm{d}V.
\]
Divide on $\{s>-N\}$ and let $N\to\infty$. Both sides of the desired inequality give zero mass to the pole set. Restore the common translation.
\end{proof}

\begin{theorem}\label{thm:individual-density-subsolution}\index{Subsolution!individual density threshold}
For every such $f$,
\[
 b_h(f)=\max\{a\geq0:\text{there exists }u\in\PSH(X,h)
       \text{ with }\mathrm{MA}^{\mathrm{ext}}_h(u)
                         \geq af\,\mathrm{d}V\}.
\]
In particular, a positive scaled subsolution exists if and only if $b_h(f)>0$.
\end{theorem}
\begin{proof}
If $\mathrm{MA}^{\mathrm{ext}}_h(u)\geq af\,\mathrm{d}V$ with $a>0$, then it dominates $af_k\,\mathrm{d}V$. \Cref{lem:bounded-normalizer-subsolution} gives $a\leq c_h(f_k)$ for every $k$, hence $a\leq b_h(f)$. If $b_h(f)=0$, the zero potential realizes the coefficient zero.

\Needspace{7\baselineskip}
Suppose $b=b_h(f)>0$. Choose continuous normalized solutions
\[
 \mathrm{MA}_h(u_k)=c_h(f_k)f_k\,\mathrm{d}V,
 \qquad \sup_Xu_k=0.
\]
After extraction, $u_{k_j}\to u$ in $L^1$, where $u$ is a genuine normalized $h$-plurisubharmonic function. The densities $c_h(f_{k_j})f_{k_j}$ converge to $bf$ in $L^1$: they are dominated by $c_h(f_1)f$ and converge pointwise almost everywhere. \Cref{lem:density-subsolution-lone-closure} gives $\mathrm{MA}^{\mathrm{ext}}_h(u)\geq bf\,\mathrm{d}V$, attaining the maximum.
\end{proof}

The argument preserves a density lower bound. It does not claim that the original normalized solutions converge in capacity, that their limit has finite mass, or that the limiting inequality is an equality. These are additional requirements in the individual equation problem. The definition is unchanged if truncation is replaced by any increasing sequence of nonzero bounded densities tending to $f$: necessity bounds its limiting constants below by $b_h(f)$, and the same upper-envelope construction bounds them above by $b_h(f)$.

\subsection{The greatest normalized threshold subsolution}\label{subsec:threshold-envelope}

Let $X$ be a connected compact complex manifold of dimension $n\geq1$, and let $h$ be a Hermitian form on $X$. Write $\mathrm{MA}^{\mathrm{ext}}_h(u)$ for the extended non-pluripolar measure of $u\in\PSH(X,h)$, and $\mathrm{MA}_h(u)$ when $u$ is bounded. The quantities $v_-(h)$ and $v_+(h)$ are respectively the infimum and supremum of the masses of bounded $h$-plurisubharmonic potentials. Fix a smooth strictly positive probability volume form $\mathrm{d}V$; all function norms below use this volume. Fix a nonzero nonnegative $f\in L^1(\mathrm{d}V)$ and use its subsolution threshold $b_h(f)$ from \eqref{eq:individual-truncated-normalizer}.

\begin{proposition}\label{prop:greatest-normalized-threshold}
For $b=b_h(f)>0$, the family
\[
 \mathcal S_f=\{s\in\PSH(X,h):\sup_Xs=0,\
             \mathrm{MA}^{\mathrm{ext}}_h(s)\geq bf\,\mathrm{d}V\}
\]
is nonempty and compact in $L^1(X)$. It has a greatest member
\[
 \Phi_f=\left(\sup_{s\in\mathcal S_f}s\right)^*.
\]
Every normalized $L^1$ cluster limit of the greatest exponential solutions from \cref{cor:individual-positive-threshold-exponential} belongs to $\mathcal S_f$ and hence lies below $\Phi_f$.
\end{proposition}
\begin{proof}
The family is nonempty by \cref{thm:individual-density-subsolution}. Its compactness follows from normalized compactness and \cref{lem:density-subsolution-lone-closure}. Ordered contact makes it stable under finite maxima. Choose a countable Choquet family for its supremum and take successive finite maxima, including one fixed member of $\mathcal S_f$. This increasing sequence has a genuine lower bound and upper bound zero. Its canonical limit is $\Phi_f$ and its convergence holds in $L^1$. The same lemma gives the required measure inequality. Its supremum is zero because it lies between a normalized member and zero. Thus $\Phi_f\in\mathcal S_f$ and it dominates every member.

For an exponential cluster limit, the densities of the measures converge in $L^1$ to $bf$ by \cref{cor:individual-positive-threshold-exponential}. Apply the lemma to these densities and the normalized potentials.
\end{proof}

This proposition asserts a greatest normalized subsolution, not equality of its measure with $bf\,\mathrm{d}V$. In particular it does not assert that $\Phi_f$ has finite mass.


\Needspace{12\baselineskip}
\begin{proposition}\label{prop:individual-normalizer-lp}
If $f\in L^p(\mathrm{d}V)$ for some $p>1$, then $b_h(f)$ equals the normalization constant of the continuous equation for $f$. In general,
\[
 b_h(f)\geq\frac{v_-(h)}{\int_Xf\,\mathrm{d}V}.
\]
If $h$ is K\"ahler, then $b_h(f)=\int_Xh^n/\int_Xf\,\mathrm{d}V$ for every nonzero nonnegative $L^1$ density.
\end{proposition}
\begin{proof}
Let $v$ be a continuous solution for $f\in L^p$, with scalar $c>0$. It is a subsolution for $cf\,\mathrm{d}V$, so the maximal-subsolution characterization gives $c\leq b_h(f)$. Conversely, the bounded-solution comparison in \cref{lem:bounded-normalizer-subsolution}, applied to the attained threshold subsolution and $v$, gives $b_h(f)\leq c$.

Each bounded solution has total mass at least $v_-(h)$, so $c_h(f_k)\int_Xf_k\,\mathrm{d}V\geq v_-(h)$. Take the limit. In the K\"ahler case, bounded Stokes instead gives the exact identity $c_h(f_k)\int_Xf_k\,\mathrm{d}V=\int_Xh^n$, which proves the last assertion.
\end{proof}

\subsection{Smooth comparison tests}\label{subsec:individual-smooth-tests}

Bounded measurable tests in \eqref{eq:individual-bounded-tests} detect concentrations that smooth tests can miss. To compare them, let $f$ be a probability density bounded below by a positive constant and define
\[
 C_h^{\mathrm{sm}}(f)=\inf_g c_h(g)\|g/f\|_{L^\infty},
\]
where $g$ ranges over smooth strictly positive probability densities. The subsolution bound gives $b_h(f)\leq C_h^{\mathrm{sm}}(f)$, with equality for smooth $f$. The next construction separates even their positivity when lower mass vanishes.

\Needspace{10\baselineskip}
\begin{theorem}\label{thm:smooth-tests-insufficient}
If $v_-(h)=0$, there is a probability density $f\in L^1(\mathrm{d}V)$ bounded below by a positive constant such that
\[
 b_h(f)=0<C_h^{\mathrm{sm}}(f).
\]
Thus no positive multiple of $f\,\mathrm{d}V$ is dominated by an extended Monge--Amp\`ere measure, although every smooth comparison test gives a common positive bound.
\end{theorem}
\begin{proof}
Let $(U_i)_{i\geq1}$ be a countable basis of nonempty open subsets of $X$. We construct numbers $\alpha_j>0$, bounded densities $q_j\not\equiv0$, and closed nowhere dense sets $K_j$ such that
\[
 \int_Xq_j\,\mathrm{d}V\leq1,\qquad
 \operatorname{supp}q_j\subset K_j,\qquad
 c_h(q_j)\leq\alpha_j2^{-j}.
\]
Start with $\alpha_1=1/4$. Once $\alpha_j$ is chosen, the identity between lower mass and the infimum of smooth normalization constants gives a smooth strictly positive probability density $g_j$ with $c_h(g_j)<\alpha_j2^{-j}$. Choose a closed nowhere dense $K_j$ so that $q_j=g_j\mathbf1_{K_j}$ is sufficiently close to $g_j$ in $L^2$ for the displayed scalar bound to hold, using \cref{lem:lp-normalizer-continuity}. Such sets exist: remove an open dense set of arbitrarily small measure from $X$; boundedness of the fixed $g_j$ makes the $L^2$ error arbitrarily small. Choose it small enough also that $q_j$ is nonzero.

After choosing $K_1,\ldots,K_N$, put
\[
 m_N=\min_{1\leq i\leq N}
    \mathrm{d}V\left(U_i\setminus\bigcup_{j=1}^NK_j\right)>0.
\]
Each set inside the minimum is nonempty and open. Choose
\[
 0<\alpha_{N+1}\leq
       \min\{\alpha_N/2,\ 2^{-N-2}m_N\},
\]
and continue. In particular,
\[
 \sum_{j>N}\alpha_j\leq2\alpha_{N+1}
             \leq2^{-N-1}m_N.
\]
The function $Q=\sum_j\alpha_jq_j$ is nonnegative and integrable. Its essential infimum on every nonempty open set is zero. To see this on $U_i$, take $N\geq i$ and remove the first $N$ sets $K_j$. On the remaining set the first $N$ summands vanish. Markov's inequality bounds the measure where the tail exceeds $2^{-N}$ by
\[
 2^N\sum_{j>N}\alpha_j\int_Xq_j\,\mathrm{d}V
       \leq m_N/2.
\]
The remaining set originally has measure at least $m_N$. A positive-measure subset therefore satisfies $Q\leq2^{-N}$. Let $N\to\infty$ and use the basis property.

\Needspace{8\baselineskip}
Set $M=\int_X(1+Q)\,\mathrm{d}V$ and $f=(1+Q)/M$. If $\mathrm{MA}^{\mathrm{ext}}_h(u)\geq af\,\mathrm{d}V$ with $a>0$, then the bounded comparison lemma gives
\[
 \frac{a\alpha_j}{M}\leq c_h(q_j)\leq\alpha_j2^{-j}
 \quad\text{for every }j.
\]
This is impossible. Hence $b_h(f)=0$.

Finally let $g$ be any smooth strictly positive probability density and put $R=\|g/f\|_\infty$. The inequality $g\leq R(1+Q)/M$ almost everywhere forces $g\leq R/M$ everywhere. Otherwise continuity would give a nonempty open set on which $g>R(1+\varepsilon)/M$ for some $\varepsilon>0$, contradicting the essential infimum property of $Q$ on that set. Order and homogeneity of the scalar now give
\[
 R c_h(g)\geq M c_h(1).
\]
The constant test $g=1$ attains this value, since $\operatorname*{ess\,inf}_XQ=0$. Thus
$C_h^{\mathrm{sm}}(f)=M c_h(1)>0$, as required.
\end{proof}

The distinction is caused by the order restriction on a smooth test: it cannot use a measurable increase of a density whose essential lower value is unchanged in every open set. Bounded measurable tests retain this information. The theorem concerns subsolution existence for individual $L^1$ densities. A characterization of exact normalized equations at $b_h(f)$, including their natural singularity class and finite-mass compactness, remains a separate problem.


\section{Gauduchon tests for integrable densities}
\label{sec:integrable-gauduchon-tests}

The normalization constants of bounded truncations have an exact description by Gauduchon metrics. Fix a Hermitian form $h$ and a smooth positive probability volume $\mathrm{d}V$. For a nonnegative function $p$, set
\[
 L_{h,\gamma}(p)=
 \frac{\displaystyle\int_Xp
              (\mathrm{d}V/\gamma^n)^{1/n}\gamma^n}
      {\displaystyle\int_Xh\wedge\gamma^{n-1}},
\]
where $\gamma$ is Gauduchon. In dimension one every Hermitian metric is permitted and this quotient is $\int_Xp\,\mathrm{d}V/\int_Xh$.

For a real $p\in L^n(\mathrm{d}V)$, define
\[
 \rho_h(p)=b_h(|p|^n)^{-1/n}\quad(p\ne0),\qquad \rho_h(0)=0,
\]
with $0^{-1/n}=+\infty$. Each individual Gauduchon test is finite on $L^n$, since its coefficient relative to $\mathrm{d}V$ is smooth. Their supremum need not be finite.

\Needspace{18\baselineskip}
\begin{theorem}\label{thm:integrable-gauduchon-formula}
For every real $p\in L^n(\mathrm{d}V)$,
\begin{equation}\label{eq:integrable-gauduchon-root}
 \rho_h(p)=\sup_{\gamma\text{ Gauduchon}}L_{h,\gamma}(|p|).
\end{equation}
Equivalently, for every nonzero nonnegative $f\in L^1(\mathrm{d}V)$,
\begin{equation}\label{eq:integrable-gauduchon-threshold}
 b_h(f)^{1/n}=
 \inf_{\gamma\text{ Gauduchon}}
 \frac{\displaystyle\int_Xh\wedge\gamma^{n-1}}
      {\displaystyle\int_X(f\,\mathrm{d}V/\gamma^n)^{1/n}\gamma^n}.
\end{equation}
The formulas allow $\rho_h(p)=+\infty$ and $b_h(f)=0$. No attainment of the infimum for a nonsmooth density is asserted.
\end{theorem}
\begin{proof}
It suffices to prove \eqref{eq:integrable-gauduchon-root} for $p\geq0$, since both sides depend only on $|p|$; the case $p=0$ is immediate. If $p$ is smooth and strictly positive, \cref{thm:gaud-volume-dual-formula}, with the density $p^n$ divided by its integral, gives
\[
 c_h(p^n)^{-1/n}=\sup_\gamma L_{h,\gamma}(p).
\]
The integral normalization cancels by homogeneity. Also $b_h(p^n)=c_h(p^n)$.

We need uniformity in $\gamma$ before approximating a measurable density. The lower-density estimate in \cref{thm:lp-potential-estimate}, with exponent two, gives a constant $a_h>0$ such that
\[
 c_h(p^n)\|p^n\|_{L^2}\geq a_h
\]
for smooth strictly positive $p$. Therefore
\begin{equation}\label{eq:gauduchon-test-uniform-bound}
 L_{h,\gamma}(p)\leq a_h^{-1/n}\|p\|_{L^{2n}}
 \quad\text{for every }\gamma.
\end{equation}
Positive smooth approximation and duality show that the smooth coefficient of each $L_{h,\gamma}$ has $L^{2n/(2n-1)}$ norm at most $a_h^{-1/n}$. In particular, the supremum of these functionals is Lipschitz on the positive cone of $L^{2n}$, with that same bound.

Now let $p\geq0$ be bounded and nonzero. Approximate it in $L^{2n}$ by smooth strictly positive $p_j$. Then $p_j^n\to p^n$ in $L^2$. The normalization constants converge:
\[
 c_h(p_j^n)\longrightarrow c_h(p^n)>0.
\]
Here the scalar continuity is the one proved at the start of \cref{thm:smooth-tests-insufficient}: the Gauduchon upper bound and the lower-density estimate bound the scalars above and away from zero, and uniform compactness followed by scalar uniqueness identifies their limit. The uniform bound \eqref{eq:gauduchon-test-uniform-bound} also gives
\[
 \sup_\gamma L_{h,\gamma}(p_j)
       \longrightarrow\sup_\gamma L_{h,\gamma}(p).
\]
The formula follows for bounded $p$.

For general nonnegative $p\in L^n$, put $p_k=\min(p,k)$. The definition of $b_h$ gives
\[
 \rho_h(p)=\lim_{k\to\infty}c_h(p_k^n)^{-1/n}
          =\sup_k\sup_\gamma L_{h,\gamma}(p_k).
\]
For each fixed $\gamma$, monotone convergence gives
$\sup_kL_{h,\gamma}(p_k)=L_{h,\gamma}(p)$. The two suprema commute, proving \eqref{eq:integrable-gauduchon-root}. Every denominator in \eqref{eq:integrable-gauduchon-threshold} is positive; taking reciprocals proves that formula, including the zero endpoint.
\end{proof}

For $n\geq2$, write $L_{h,\gamma}(p)=\int_Xp\,a_\gamma\,\mathrm{d}V$. Its coefficient satisfies
\[
 \|a_\gamma\|_{L^{n/(n-1)}}
 =\frac{(\int_X\gamma^n)^{(n-1)/n}}
        {\int_Xh\wedge\gamma^{n-1}}.
\]
Thus \cref{thm:gaud-volume-mass-extrema} identifies positive lower mass with a uniform bound for these endpoint norms. The weaker bound used above is available on every compact Hermitian manifold. It suffices to pass from smooth tests to bounded densities; monotone truncation then treats all integrable densities.

\begin{proposition}\label{prop:individual-threshold-backgrounds}
If $h_1,h_2$ are Hermitian and $m h_1\leq h_2\leq M h_1$ for positive constants $m,M$, then
\[
 m^n b_{h_1}(f)\leq b_{h_2}(f)\leq M^n b_{h_1}(f),
 \qquad
 M^{-1}\rho_{h_1}(p)\leq\rho_{h_2}(p)\leq m^{-1}\rho_{h_1}(p).
\]
If $h_2-h_1=\ddc a$ for a smooth real function $a$, their thresholds agree.
\end{proposition}
\begin{proof}
For every Gauduchon metric the numerators in \eqref{eq:integrable-gauduchon-threshold} satisfy the same bounds with factors $m,M$; the denominators agree. Taking the infimum and then the $n$th power proves the threshold bounds. Taking reciprocals gives the bounds for $\rho$, with the stated conventions at zero and infinity. In the last assertion, distributional Stokes makes the numerators equal.
\end{proof}

All Hermitian metrics on compact $X$ are comparable. Thus positivity of the individual threshold and the set of admissible density roots are independent of the chosen Hermitian metric. If the reference volume is replaced by $\mathrm{d}V'=a\,\mathrm{d}V$ with $a>0$ smooth, the same measure has density $f'=f/a$. The maximal-subsolution characterization gives $b'_h(f')=b_h(f)$, and $p\mapsto a^{-1/n}p$ preserves both the root functional and the $L^n$ norm. This identifies the construction under a change of reference volume.


\section{Stability of the subsolution threshold}
\label{sec:individual-density-stability}

We now study how the threshold $b_h(f)$ depends on $f$. By \eqref{eq:integrable-gauduchon-threshold}, $b_h^{1/n}$ is an infimum of functions that are continuous on $L^1$, and is therefore upper semicontinuous. It is not continuous in general: when $v_-(h)=0$, every density with positive threshold is an $L^1$ limit of densities with zero threshold. For density roots, \eqref{eq:integrable-gauduchon-root} expresses $\rho_h$ as a supremum of positive linear functionals. Hence $\rho_h$ is a norm on the space where it is finite, this space is complete for a suitable norm, and the threshold depends continuously on the root in that norm.

\begin{theorem}\label{thm:individual-threshold-semicontinuity}
Suppose $f_j,f$ are nonzero nonnegative $L^1$ densities and $f_j\to f$ in $L^1(\mathrm{d}V)$. Then
\[
 \limsup_{j\to\infty}b_h(f_j)\leq b_h(f).
\]
Moreover,
\begin{equation}\label{eq:individual-threshold-relaxation}
 b_h(f)=\lim_{\delta\downarrow0}\sup\{c_h(q):q\in L^\infty,\ q\geq0,\ q\not\equiv0,
                                      \ \|q-f\|_{L^1}<\delta\},
\end{equation}
and the suprema are finite for sufficiently small $\delta$.
\end{theorem}
\begin{proof}
For a Gauduchon form $\gamma$ on $X$, write $A_\gamma(f)=\int_X(f\,\mathrm{d}V/\gamma^n)^{1/n}\gamma^n$. Since $|s^{1/n}-t^{1/n}|\leq|s-t|^{1/n}$ for $s,t\geq0$, H\"older's inequality shows that $A_\gamma$ is continuous on $L^1$; it is positive at every nonzero nonnegative density. By \cref{thm:integrable-gauduchon-formula}, $b_h^{1/n}$ is the infimum over $\gamma$ of the continuous positive functions $f\mapsto\int_Xh\wedge\gamma^{n-1}/A_\gamma(f)$. It is therefore upper semicontinuous, and testing a single $\gamma$ bounds it above on an $L^1$ neighborhood of $f$.

For the second assertion, recall that $b_h(q)=c_h(q)$ for bounded $q$ (\cref{prop:individual-normalizer-lp}). By the first assertion, $\limsup_jc_h(q_j)\leq b_h(f)$ for every sequence of nonzero nonnegative bounded densities $q_j\to f$ in $L^1$. Hence the limit in \eqref{eq:individual-threshold-relaxation} is at most $b_h(f)$, and the local upper bound makes the suprema finite for small $\delta$. Conversely, the truncations $f_k$ converge to $f$ in $L^1$ and satisfy $c_h(f_k)\geq b_h(f)$, so every supremum in \eqref{eq:individual-threshold-relaxation} is at least $b_h(f)$.
\end{proof}

Since $b_h\geq0$, upper semicontinuity makes $b_h$ continuous at every density where it vanishes. If $v_-(h)=0$, these are its only points of continuity. Indeed, let $b_h(f)>0$ and $g=f/\int_Xf\,\mathrm{d}V$. By \cref{cor:lower-mass-lone-instability} and \cref{thm:individual-density-subsolution}, the densities $(1-t)g+tf_0$, where $f_0$ is the density of \cref{thm:lower-mass-obstructing-measure}, have zero threshold and converge to $g$ in $L^1$ as $t\downarrow0$; multiplying them by $\int_Xf\,\mathrm{d}V$ gives densities with zero threshold converging to $f$. These statements concern only the scalar $b_h$; they say nothing about capacity convergence of the corresponding potentials or about the limits of their Monge--Amp\`ere measures.

\subsection{Finite sums of densities}\label{subsec:finite-density-sums}

Since \eqref{eq:integrable-gauduchon-root} expresses $\rho_h(p)$ as a supremum of positive linear functionals evaluated at $|p|$, the root functional is homogeneous, monotone and subadditive. Subadditivity shows that the sum of two densities with positive threshold again has positive threshold. In the statements below, $f$ and $g$ are nonzero nonnegative $L^1$ densities, and $p$ and $q$ are real functions in $L^n$.

\Needspace{12\baselineskip}
\begin{proposition}\label{prop:individual-root-triangle}
The functional $\rho_h$ satisfies $\rho_h(tp)=|t|\rho_h(p)$ for real $t\ne0$,
$\rho_h(p)\leq\rho_h(q)$ whenever $|p|\leq|q|$, and
\[
 \rho_h(p+q)\leq\rho_h(p)+\rho_h(q).
\]
In particular, if $b_h(f),b_h(g)>0$, then
\[
 b_h(f+g)\geq
       \bigl(b_h(f)^{-1/n}+b_h(g)^{-1/n}\bigr)^{-n}>0.
\]
\end{proposition}
\begin{proof}
By \eqref{eq:integrable-gauduchon-root}, $\rho_h(p)=\sup_\gamma L_{h,\gamma}(|p|)$ is a supremum of positive linear functionals evaluated at $|p|$. This gives homogeneity and monotonicity, and, together with $|p+q|\leq|p|+|q|$, the triangle inequality. Finally, $(f+g)^{1/n}\leq f^{1/n}+g^{1/n}$ gives
\[
 b_h(f+g)^{-1/n}=\rho_h\bigl((f+g)^{1/n}\bigr)\leq\rho_h(f^{1/n})+\rho_h(g^{1/n})
 =b_h(f)^{-1/n}+b_h(g)^{-1/n}.
\]
\end{proof}

\begin{proposition}\label{prop:individual-density-banach}\index{Density space!completeness}
The real vector space
\[
 \mathcal D_h=\{p\in L^n(\mathrm{d}V):\rho_h(p)<+\infty\}
\]
is complete for the norm $\|p\|_{\mathcal D_h}=\|p\|_{L^n}+\rho_h(p)$. It is an order ideal in $L^n$: if $q\in\mathcal D_h$ and $|p|\leq|q|$, then $p\in\mathcal D_h$. It contains $L^\infty$ and is therefore dense in $L^n$ for the $L^n$ norm.

Moreover, $\mathcal D_h=L^n$ if and only if $v_-(h)>0$, and in that case the norm $\|\cdot\|_{\mathcal D_h}$ is equivalent to the $L^n$ norm. If $h$ is K\"ahler, then
\[
 \rho_h(p)=\left(\int_Xh^n\right)^{-1/n}\|p\|_{L^n}.
\]
\end{proposition}
\begin{proof}
By \cref{prop:individual-root-triangle}, $\mathcal D_h$ is a vector space and an order ideal, and $\|\cdot\|_{\mathcal D_h}$ is a norm on it. Bounded functions belong to $\mathcal D_h$, because $b_h(|p|^n)=c_h(|p|^n)>0$ for bounded nonzero $p$.

\Needspace{5\baselineskip}
We next show that $\rho_h$ is lower semicontinuous on $L^n$. If $p_j\to p$ in $L^n$, then $|p_j|^n\to|p|^n$ in $L^1$, and the upper semicontinuity of $b_h$ in \cref{thm:individual-threshold-semicontinuity} gives
\[
 \rho_h(p)\leq\liminf_j\rho_h(p_j).
\]
This includes the case $b_h(|p|^n)=0$, in which the right side must be infinite; for $p=0$ the inequality is trivial. Now let $(p_j)$ be a Cauchy sequence in $\mathcal D_h$. It converges in $L^n$ to some $p$, and lower semicontinuity gives $\rho_h(p)\leq\liminf_j\rho_h(p_j)<+\infty$, so $p\in\mathcal D_h$. Applied to $p_j-p_k$ as $k\to\infty$, lower semicontinuity also gives
\[
 \rho_h(p_j-p)\leq\liminf_{k\to\infty}\rho_h(p_j-p_k),
\]
which tends to zero as $j\to\infty$. Together with $L^n$ convergence, this proves completeness.

If $v_-(h)>0$, \cref{prop:individual-normalizer-lp} gives
$\rho_h(p)\leq v_-(h)^{-1/n}\|p\|_{L^n}$, so $\mathcal D_h=L^n$ with an equivalent norm. If $v_-(h)=0$, the density $f$ of \cref{thm:lower-mass-obstructing-measure} has $b_h(f)=0$ by \cref{thm:individual-density-subsolution}; thus $f^{1/n}$ belongs to $L^n$ but not to $\mathcal D_h$. The K\"ahler formula follows from the same proposition.
\end{proof}

Thus, when the lower mass vanishes, $\mathcal D_h$ is a proper dense subspace of $L^n$, complete for the stronger norm $\|\cdot\|_{\mathcal D_h}$; its nonnegative elements are the $n$th roots of the densities admitting a positive scaled subsolution. A finite sum of such densities again admits a positive scaled subsolution, but a convergent infinite sum need not: the smooth summands in \cref{thm:lower-mass-obstructing-measure} give an example. By \cref{prop:individual-threshold-backgrounds}, the space $\mathcal D_h$ does not depend on the Hermitian form $h$, and the norms $\|\cdot\|_{\mathcal D_h}$ for different Hermitian forms are equivalent. The closure of the bounded functions for this norm is studied in \cref{subsec:bounded-root-approximation}.

The restriction of $\rho_h$ to $\mathcal D_h$ is itself a norm: if $p\ne0$, then $\rho_h(p)\geq L_{h,\gamma}(|p|)>0$ for any Gauduchon form $\gamma$. Whether this weaker norm is complete depends on the upper mass.

\Needspace{15\baselineskip}
\begin{proposition}\label{prop:individual-root-completeness}\index{Density space!root norm}
The norm $\rho_h$ is complete on $\mathcal D_h$ if and only if $v_+(h)<+\infty$. When the upper mass is finite,
\[
 \|p\|_{L^n}\leq v_+(h)^{1/n}\rho_h(p)\qquad(p\in\mathcal D_h),
\]
so $\rho_h$ and $\|\cdot\|_{\mathcal D_h}$ are equivalent norms.
\end{proposition}
\begin{proof}
Suppose $v_+(h)<+\infty$. A continuous solution for a bounded density $q$ has mass $c_h(q)\int_Xq\,\mathrm{d}V\leq v_+(h)$. Applying this to the truncations of $|p|^n$ and letting the truncation level increase gives $b_h(|p|^n)\|p\|_{L^n}^n\leq v_+(h)$, which is the asserted inequality. Hence $\rho_h$ and $\|\cdot\|_{\mathcal D_h}$ are equivalent, and completeness follows from \cref{prop:individual-density-banach}.

Suppose instead that $v_+(h)=+\infty$. By \cref{thm:gaud-volume-mass-extrema,thm:gaud-volume-dual-formula}, the normalization scalars of smooth strictly positive probability densities are unbounded above. For each $j$, choose such a density $f_j$ with
\[
 c_h(f_j)\geq(j2^j)^n,
 \qquad\text{and put}\qquad p_j=jf_j^{1/n}.
\]
Then $\|p_j\|_{L^n}=j$ and $\rho_h(p_j)=jc_h(f_j)^{-1/n}\leq2^{-j}$, so the partial sums $P_N=\sum_{j\leq N}p_j$ form a Cauchy sequence for $\rho_h$.

Fix a Gauduchon form $\gamma$. Its coefficient $a_\gamma$ is bounded below by a positive constant, so \eqref{eq:integrable-gauduchon-root} gives $\rho_h(p)\geq c\|p\|_{L^1}$ with a constant $c>0$. Hence $\sum_jp_j$ converges in $L^1$ to a nonnegative function $P$. If $P_N$ converged for $\rho_h$ to an element of $\mathcal D_h$, that element would also be the $L^1$ limit, hence equal to $P$ almost everywhere. But $P\geq p_j$ gives $\|P\|_{L^n}\geq j$ for every $j$, so $P\notin L^n$. Therefore $\rho_h$ is not complete.
\end{proof}

The two mass properties play separate roles: positive lower mass makes every $L^n$ function admissible, whereas finite upper mass makes the root norm control the $L^n$ norm. The norm $\|\cdot\|_{\mathcal D_h}$ is complete without either assumption.

All four combinations occur. On a compact K\"ahler manifold both mass properties hold. The Hopf threefold has positive lower mass and infinite upper mass (\cref{cor:hopf-lower-mass}). The elliptic threefold of \cref{thm:vanishing-lower-mass} has zero lower mass, and finite upper mass for its balanced metric (\cref{prop:lower-mass-balanced-metric}). The product manifolds of \cref{cor:both-mass-properties-fail} have zero lower mass and infinite upper mass. Since neither $\mathcal D_h$ nor the completeness of $\rho_h$ depends on the choice of $h$ (\cref{prop:individual-threshold-backgrounds,prop:bmp-metric}), we obtain the following table.
\begin{center}
\begin{tabular}{@{}l@{\qquad}l@{\qquad}l@{}}
Manifold & Domain $\mathcal D_h$ & Is $\rho_h$ complete?\\
\hline
K\"ahler manifold & $L^n$ & Yes\\
Hopf threefold & $L^n$ & No\\
Elliptic threefold & $\mathcal D_h\subsetneq L^n$ & Yes\\
Product example & $\mathcal D_h\subsetneq L^n$ & No
\end{tabular}
\end{center}

\Needspace{14\baselineskip}
\subsection{Approximation by bounded density roots}\label{subsec:bounded-root-approximation}

Bounded functions are dense in $\mathcal D_h$ for the $L^n$ norm, but not obviously for the norm $\|\cdot\|_{\mathcal D_h}$. Let $\mathcal D_h^0$ denote the closure of $L^\infty$ in $\mathcal D_h$ for this norm. We show that a function belongs to $\mathcal D_h^0$ exactly when its tails are small uniformly over all Gauduchon tests. The basic tool is a subcritical estimate. For every $r>n$,
\begin{equation}\label{eq:root-subcritical-bound}
 \rho_h(p)\leq C_r\|p\|_{L^r}\qquad(p\in L^r),
\end{equation}
where the constant $C_r$ depends on $h$, the reference volume and $r$, but not on $p$. Indeed, for nonzero $p\in L^r$ the density $|p|^n$ belongs to $L^{r/n}$, so $b_h(|p|^n)=c_h(|p|^n)$ by \cref{prop:individual-normalizer-lp}, and the lower-density estimate of \cref{thm:lp-potential-estimate}, with exponent $r/n$, gives $c_h(|p|^n)\||p|^n\|_{L^{r/n}}\geq c_{r/n}>0$.

\begin{proposition}\label{prop:root-uniform-tails}\index{Uniform tails!root density}
For $p\in\mathcal D_h$, the following conditions are equivalent:
\begin{enumerate}
\item $p\in\mathcal D_h^0$;
\item $\rho_h(|p|\mathbf1_{\{|p|>K\}})\to0$ as $K\to\infty$;
\item whenever $|q_j|\leq|p|$ and $q_j\to0$ in measure, one has $\|q_j\|_{\mathcal D_h}\to0$.
\end{enumerate}
By \eqref{eq:integrable-gauduchon-root}, the second condition can be written as
\[
 \sup_{\gamma\text{ Gauduchon}}
       L_{h,\gamma}(|p|\mathbf1_{\{|p|>K\}})\longrightarrow0
 \qquad(K\to\infty).
\]
The space $\mathcal D_h^0$ is a closed order ideal in $\mathcal D_h$. It contains $L^r$ for every $r>n$, and smooth functions are dense in it. If $v_-(h)>0$, then $\mathcal D_h^0=\mathcal D_h=L^n$.
\end{proposition}
\begin{proof}
For $K>0$, set $p_K=\operatorname{sgn}(p)\min(|p|,K)$. Then $|p-p_K|=(|p|-K)_+$ and
\[
 (|p|-K)_+\leq |p|\mathbf1_{\{|p|>K\}},\qquad
 |p|\mathbf1_{\{|p|>2K\}}\leq2(|p|-K)_+.
\]
Since the $L^n$ norms of the tails tend to zero, these inequalities and the monotonicity of $\rho_h$ show that the second condition is equivalent to $p_K\to p$ in $\mathcal D_h$; in particular, it implies the first. Conversely, if $a$ is bounded and $K\geq\|a\|_{L^\infty}$, then $(|p|-K)_+\leq|p-a|$, so $\|p-p_K\|_{\mathcal D_h}\leq\|p-a\|_{\mathcal D_h}$. Approximating $p$ by bounded functions $a$ shows that the first condition implies $p_K\to p$, hence the second.

Assume the second condition, and let $q_j$ be as in the third. Fix $K$. On $\{|p|\leq K\}$, the functions $q_j$ are bounded by $K$ and tend to zero in measure, hence in $L^{2n}$; by \eqref{eq:root-subcritical-bound} with $r=2n$, the root functional of $q_j\mathbf1_{\{|p|\leq K\}}$ tends to zero. On the complementary set, $\rho_h(q_j\mathbf1_{\{|p|>K\}})\leq\rho_h(|p|\mathbf1_{\{|p|>K\}})$. Letting first $j$ and then $K$ tend to infinity gives $\rho_h(q_j)\to0$, and dominated convergence gives $\|q_j\|_{L^n}\to0$. Conversely, the third condition applied to $q_j=|p|\mathbf1_{\{|p|>j\}}$ gives the second.

Since the second condition passes from $p$ to every function dominated by $|p|$, the space $\mathcal D_h^0$ is an order ideal; it is closed by definition. For $p\in L^r$ with $r>n$, the truncations $p_K$ converge to $p$ in $L^r$, hence in $\mathcal D_h$ by \eqref{eq:root-subcritical-bound}. Bounded functions can be approximated in $L^{2n}$ by smooth functions, which then converge in $\mathcal D_h$ by the same estimate; hence smooth functions are dense in $\mathcal D_h^0$. Under positive lower mass, \cref{prop:individual-density-banach} shows that $\mathcal D_h=L^n$ with an equivalent norm, and bounded functions are dense in $L^n$; hence $\mathcal D_h^0=\mathcal D_h=L^n$.
\end{proof}

By \cref{prop:individual-threshold-backgrounds}, the space $\mathcal D_h^0$ does not depend on the choice of the Hermitian form $h$. Convergence in $\mathcal D_h$ controls the threshold: if $p_j\to p\ne0$ in $\mathcal D_h$, then
\[
 |\rho_h(p_j)-\rho_h(p)|\leq\rho_h(p_j-p)\longrightarrow0,
 \qquad b_h(|p_j|^n)\longrightarrow b_h(|p|^n)>0.
\]
For $p\in\mathcal D_h^0$, the uniform tail condition also gives a uniform modulus for exponential damping and error bounds for the scalars of the exponential equations; see \cref{prop:exponential-root-tail-modulus}. It does not by itself give capacity convergence of the solutions. We do not know whether $\mathcal D_h^0=\mathcal D_h$ when $v_-(h)=0$.



\section{The zero-exponent limit and its topology}
\label{sec:zero-exponent-topology}

For an unbounded prescribed type, the zero-exponent limit must be formulated through normalized potentials and their measures. Even multiplying an individual unbounded potential by a positive exponent leaves its uniform distance from every finite constant infinite. The useful convergence comes instead from integration against the fixed right-hand-side measure.

\Cref{prop:singular-no-uniform-limit} makes this obstruction precise for every unbounded ordinary type.

For a fixed finite non-pluripolar measure $\mu$ and solutions of $\mathrm{MA}_\theta(u_\lambda)=\mathrm{e}^{\lambda u_\lambda}\mu$, write
\[
 v_\lambda=u_\lambda-\sup_Xu_\lambda,\qquad
 c_\lambda=\mathrm{e}^{\lambda\sup_Xu_\lambda}.
\]
By \cref{lem:normalized-exponential-npp}, if $c_\lambda\to c<+\infty$, then the right-hand sides converge in total variation to $c\mu$. For an $L^p$ density, \cref{lem:normalized-exponential-lp} gives a linear error in the exponential factor. Identifying the measure of a limit of $v_\lambda$ additionally requires a continuity theorem for the Monge--Amp\`ere operator.

Here is one precise continuity input that supplies this step in a common regime.

\begin{theorem}\label{thm:relative-dominated-continuity}
Assume now that $\theta$ is semipositive and big, that the reference $\varphi$ satisfies \eqref{eq:als-model-condition}, and that $v_-(\theta)>0$. These are the background hypotheses of \cref{thm:als-general-measure}. Use the relative class $\mathcal E_{\mathrm{env}}(X,\theta,\varphi)$ defined in \cref{def:relative-envelope-class}. Suppose $u_j\in\mathcal E_{\mathrm{env}}(X,\theta,\varphi)$, $u_j\to u\in\PSH(X,\theta)$ in $L^1(X)$, and
\[
 \mathrm{MA}_\theta(u_j)\leq\nu
\]
for one finite non-pluripolar measure $\nu$. Then $u\in\mathcal E_{\mathrm{env}}(X,\theta,\varphi)$, the convergence holds in capacity, and
$\mathrm{MA}_\theta(u_j)\rightharpoonup\mathrm{MA}_\theta(u)$.
\end{theorem}

A complete reconstruction is given in \cref{thm:relative-fixed-measure-continuity}. Its proof first constructs a common rooftop envelope after extraction, establishes finite mass of the limiting lower roofs, and then applies domination. Mere $L^1$ convergence, even for uniformly bounded potentials, is insufficient for the Monge--Amp\`ere passage to the limit. The common dominating measure and the positive lower mass assumption are the substantive hypotheses here.

\section{Limits of positive-exponent equations}
\label{sec:exponential-closedness}

At a strictly positive limiting exponent, exponential decay controls the mass on deep sublevels uniformly. Once the potentials converge in capacity, the tail criterion identifies their limiting Monge--Amp\`ere measure. Convergence of the right-hand-side densities then identifies the equation.

\Needspace{15\baselineskip}
\begin{theorem}\label{thm:exponential-closedness}
Fix a Hermitian metric on $X$, let $\theta$ be any smooth real $(1,1)$-form, and let $\mu$ be a finite non-pluripolar measure. Suppose $f_j\geq0$, $f_j\to f$ in $L^1(\mu)$, and $\lambda_j>0$ with $\lambda_j\to\lambda>0$. Let $u_j,u\in\PSH(X,\theta)$ be nonpositive, suppose $u_j\to u$ in capacity, and assume
\[
 \mathrm{MA}^{\mathrm{ext}}_\theta(u_j)
     =\mathrm{e}^{\lambda_j u_j}f_j\,\mu.
\]
Then $\mathrm{MA}^{\mathrm{ext}}_\theta(u)$ is finite and
\[
 \mathrm{MA}^{\mathrm{ext}}_\theta(u)=\mathrm{e}^{\lambda u}f\,\mu.
\]
The measures on the left converge weakly; the right-hand sides converge in total variation.
\end{theorem}
\begin{proof}
Discard finitely many indices so that $\lambda_j\geq\lambda/2$. Their masses are at most $\|f_j\|_{L^1(\mu)}$, hence uniformly bounded. On a deep sublevel,
\[
 \mathrm{MA}^{\mathrm{ext}}_\theta(u_j)(\{u_j<-T\})
 \leq \mathrm{e}^{-\lambda T/2}\|f_j\|_{L^1(\mu)}.
\]
The tail criterion \cref{thm:undamped-tail-continuity}, in its general-background form \cref{thm:general-background-damping}, shows that the limiting Monge--Amp\`ere measure is finite and that the left-hand sides converge weakly to it.

It remains to identify this weak limit. By \cref{lem:npp-uniform-capacity}, $u_j\to u$ in $\mu$-measure. The functions $\mathrm{e}^{\lambda u_j}$ are bounded by one and converge in $\mu$-measure to $\mathrm{e}^{\lambda u}$. The change of exponent is uniform for nonpositive arguments, because the mean-value theorem gives
\[
 \sup_{t\leq0}|\mathrm{e}^{\lambda_jt}-\mathrm{e}^{\lambda t}|
 \leq |\lambda_j-\lambda|
       \sup_{s\geq0}s \mathrm{e}^{-(\lambda/2)s}\longrightarrow0.
\]
Thus $\mathrm{e}^{\lambda_j u_j}\to \mathrm{e}^{\lambda u}$ in $\mu$-measure and, by absolute continuity, in $f\mu$-measure as well. Boundedness by one and finiteness of $f\mu$ give convergence in $L^1(f\mu)$. Therefore
\[
 \begin{split}
 \int_X|\mathrm{e}^{\lambda_j u_j}f_j-\mathrm{e}^{\lambda u}f|\,\mathrm{d}\mu
 &\leq\|f_j-f\|_{L^1(\mu)}\\
 &\quad+\int_X|\mathrm{e}^{\lambda_j u_j}-\mathrm{e}^{\lambda u}|f\,\mathrm{d}\mu
 \longrightarrow0.
 \end{split}
\]
This proves total-variation convergence and identifies the equation for $u$.
\end{proof}

Normalizing constants can be absorbed into $f_j$ whenever the resulting densities converge in $L^1(\mu)$. The theorem does not assert that an arbitrary $L^1$-convergent sequence converges in capacity or stays in a prescribed envelope class.

This argument also pinpoints the difficulty as $\lambda_j\downarrow0$: the tail bound $\mathrm{e}^{-\lambda_jT}\|f_j\|_1$ no longer tends to zero uniformly in $j$ as $T\to\infty$. A zero-exponent approximation needs another source of uniform tail control, such as the integrability condition in \cref{cor:integrability-continuity}. Global bounded mass would not by itself supply that missing tail estimate.

\section{Exponential equations for an individual density}
\label{sec:individual-exponential-equations}

Fix a nonzero nonnegative $f\in L^1(\mathrm{d}V)$. An exponential subsolution is a common lower bound for the continuous solutions of the truncated equations. Those solutions decrease, so \cref{thm:exponential-closedness} gives an exact equation at the limit. The limiting solution dominates every subsolution. A positive threshold $b_h(f)$ supplies the required lower bound for each exponent.

\begin{theorem}\label{thm:individual-exponential-subsolution}\index{Monge--Ampere equation@Monge--Amp\`ere equation!greatest exponential solution}
Fix $\lambda>0$. Suppose there is $s\in\PSH(X,h)$ such that
\[
 \mathrm{MA}^{\mathrm{ext}}_h(s)\geq\mathrm{e}^{\lambda s}f\,\mathrm{d}V.
\]
Then the equation
\[
 \mathrm{MA}^{\mathrm{ext}}_h(u_\lambda)
       =\mathrm{e}^{\lambda u_\lambda}f\,\mathrm{d}V
\]
has a greatest solution. More precisely, $u_\lambda$ is the greatest genuine $h$-plurisubharmonic subsolution of this equation, and its Monge--Amp\`ere measure is finite. It is the decreasing limit of the continuous solutions for $f_k=\min(f,k)$.
\end{theorem}
\begin{proof}
Let $u_k$ solve the continuous exponential equation for $f_k$, using \cref{thm:nguyen-exp}. By \cref{lem:exp-order}, $u_{k+1}\leq u_k$.

Every genuine subsolution $v$ for $f$ lies below each $u_k$. Indeed, for $t>0$, on $\{u_k<v-t\}$ we have
\[
 \mathrm{MA}_h(u_k)=\mathrm{e}^{\lambda u_k}f_k\,\mathrm{d}V
 \leq\mathrm{e}^{-\lambda t}\mathrm{e}^{\lambda v}f\,\mathrm{d}V
 \leq\mathrm{e}^{-\lambda t}\mathrm{MA}^{\mathrm{ext}}_h(v-t).
\]
The first potential $u_k$ is bounded, so $P_{h/2}(u_k)$ is genuine and $v-t\preceq u_k$. Strict domination, \cref{thm:singular-strict-domination}, gives $u_k\geq v-t$. Let $t\downarrow0$.

In particular, $s\leq u_k\leq u_1$. Thus the decreasing limit $u_\lambda$ is genuine. By \cref{lem:monotone-psh-capacity}, the convergence holds in capacity. Choose a common constant $A\geq\sup_Xu_1$. Apply \cref{thm:exponential-closedness} to $u_k-A$, with densities $\mathrm{e}^{\lambda A}f_k\to\mathrm{e}^{\lambda A}f$ in $L^1$. It proves the equation and finiteness of its left side. Finally every subsolution $v$ lies below all $u_k$, hence below $u_\lambda$, which is itself a subsolution.
\end{proof}

Every subsolution lies below this solution, so its extremal property is independent of the truncation used to construct it. Smaller solutions with other singularities are not excluded. For $f\in L^p$, $p>1$, the continuous solution gives a bounded lower barrier. The greatest solution is then bounded and agrees with it by bounded exponential comparison.

\Needspace{20\baselineskip}
\begin{corollary}\label{cor:individual-positive-threshold-exponential}
If $b=b_h(f)>0$, the preceding greatest solution exists for every $\lambda>0$. If $v\in\PSH(X,h)$ has $\sup_Xv=0$ and
$\mathrm{MA}^{\mathrm{ext}}_h(v)\geq bf\,\mathrm{d}V$, then
\[
 u_\lambda\geq v+\lambda^{-1}\log b.
\]
Writing
\[
 M_\lambda=\sup_Xu_\lambda,\qquad
 v_\lambda=u_\lambda-M_\lambda,\qquad
 c_\lambda=\mathrm{e}^{\lambda M_\lambda},
\]
we have
\[
 c_\lambda\longrightarrow b_h(f),\qquad
 \mathrm{e}^{\lambda u_\lambda}f\longrightarrow b_h(f)f
           \quad\text{in }L^1(\mathrm{d}V)
 \quad(\lambda\downarrow0).
\]
\end{corollary}
\begin{proof}
The attained-threshold theorem supplies $v$. Since $v\leq0$, the function $s=v+\lambda^{-1}\log b$ satisfies
\[
 \mathrm{MA}^{\mathrm{ext}}_h(s)\geq bf\,\mathrm{d}V
       \geq\mathrm{e}^{\lambda s}f\,\mathrm{d}V.
\]
The theorem gives $u_\lambda\geq s$, hence $c_\lambda\geq b$.

Set $g_\lambda=\mathrm{e}^{\lambda v_\lambda}f$. By \cref{lem:normalized-exponential-npp}, applied to the fixed finite measure $f\,\mathrm{d}V$, we have $g_\lambda\to f$ in $L^1$. Each $g_\lambda$ is nonzero, since a genuine quasi-plurisubharmonic function is finite almost everywhere. The equation for $u_\lambda$ implies
\[
 c_\lambda\leq b_h(g_\lambda)
\]
by the maximal subsolution characterization. Upper semicontinuity therefore gives
\[
 b\leq\liminf_{\lambda\downarrow0}c_\lambda
 \leq\limsup_{\lambda\downarrow0}c_\lambda
 \leq b_h(f)=b.
\]
This proves the scalar limit. Since $g_\lambda\to f$ in $L^1$, it also proves $c_\lambda g_\lambda\to bf$ in $L^1$.
\end{proof}

The right-hand-side measures therefore converge in total variation. This answers the scalar and measure questions in \cref{sec:zero-exponent-topology}. Under positive lower mass, \cref{cor:positive-lower-mass-exact-threshold} also gives capacity convergence along every $L^1$-convergent normalized subsequence and identifies an exact solution at $b_h(f)$. For $L^p$ data, $p>1$, \cref{prop:symmetric-equation-clusters} gives uniform precompactness. In either case, convergence of the whole normalized family requires identification of all cluster limits. Smooth strictly positive dual attainment gives whole-family uniform convergence in \cref{cor:smooth-dual-exponential-limit}. If the smooth maximizing form is only semipositive, but strictly positive outside a proper analytic set, \cref{prop:semipositive-dual-exponential-limit} gives whole-family $L^1$ convergence and a fixed lower barrier. For $L^p$ data, $p>1$, \cref{cor:semipositive-dual-lp-limit} strengthens this to uniform convergence.

\Needspace{7\baselineskip}
For a general positive individual threshold, the corollary gives
\[
 v_\lambda\geq v-\lambda^{-1}\log(c_\lambda/b_h(f)).
\]
The next estimate quantifies this scalar error when the density root is approximable by bounded roots. Whether an exponential subsolution exists when $b_h(f)=0$ is a separate question; \cref{thm:individual-exponential-subsolution} applies whenever such a subsolution is available.

\subsection{An error bound for the scalar limit}

Suppose that $f^{1/n}$ belongs to the space $\mathcal D_h^0$ of \cref{subsec:bounded-root-approximation}. The uniform tail condition of \cref{prop:root-uniform-tails} then gives a rate for the convergence $c_\lambda\to b_h(f)$ in \cref{cor:individual-positive-threshold-exponential}. Split the root $p=f^{1/n}$ at a height $K$. On $\{p\leq K\}$, the error $p(1-\mathrm{e}^{\lambda v/n})$ of the damped root has norm $O(K\lambda)$ in $\mathcal D_h$, uniformly over normalized potentials $v$, while on $\{p>K\}$ it is bounded by the tail of $p$, which is small in that norm. Optimizing over $K$ gives a modulus valid for every normalized potential.

\begin{proposition}\label{prop:exponential-root-tail-modulus}
Let $0\leq p\in\mathcal D_h^0$ be nonzero, and put $f=p^n$, $R=\rho_h(p)$ and $b=R^{-n}=b_h(f)$. There is a constant $C_h$, depending only on $h$ and $\mathrm{d}V$, such that the function
\[
 \eta_p(\lambda)=\inf_{K>0}
 \left(\|p\mathbf1_{\{p>K\}}\|_{\mathcal D_h}+C_hK\lambda\right)
\]
tends to zero as $\lambda\downarrow0$ and satisfies
\begin{equation}\label{eq:uniform-root-damping}
 \big\|p(1-\mathrm{e}^{\lambda v/n})\big\|_{\mathcal D_h}
 \leq\eta_p(\lambda)
 \qquad\bigl(v\in\PSH(X,h),\ \sup_Xv=0\bigr).
\end{equation}
Let $u_\lambda$, $v_\lambda$ and $c_\lambda$ be the greatest exponential solutions and their normalizations from \cref{cor:individual-positive-threshold-exponential}. If $\eta_p(\lambda)<R$, then
\[
 b\leq c_\lambda\leq\bigl(R-\eta_p(\lambda)\bigr)^{-n}
\]
and
\begin{align*}
 \big\|\mathrm{MA}^{\mathrm{ext}}_h(v_\lambda)
              -bf\,\mathrm{d}V\big\|_{\mathrm{TV}}
 &\leq(c_\lambda-b)\|p\|_{L^n}^{n}
       +nb\|p\|_{L^n}^{n-1}\eta_p(\lambda).
\end{align*}
If $p\in L^r$ for some finite $r>n$, the same bounds hold with $\eta_p(\lambda)$ replaced by $C_{h,r}\|p\|_{L^r}\lambda$, where $C_{h,r}$ depends only on $h$, $\mathrm{d}V$ and $r$. In particular, for each such $p$ the scalar and total-variation errors are $O(\lambda)$.
\end{proposition}
\Needspace{10\baselineskip}
\begin{proof}
By \cref{lem:uniform-exponential}, normalized $h$-plurisubharmonic functions are uniformly bounded in every $L^s$ with $s<+\infty$. By \eqref{eq:root-subcritical-bound} with $r=2n$, their norms $\|v\|_{\mathcal D_h}$ are therefore uniformly bounded; let $C_h$ be $1/n$ times this bound. Let $v$ be normalized and $K>0$. On $\{p\leq K\}$ we have
\[
 0\leq p(1-\mathrm{e}^{\lambda v/n})
     \leq K\lambda(-v)/n,
\]
since $1-\mathrm{e}^{-x}\leq x$ for $x\geq0$, while on $\{p>K\}$ we use the bound $p(1-\mathrm{e}^{\lambda v/n})\leq p$. The monotonicity of $\|\cdot\|_{\mathcal D_h}$ and the triangle inequality give
\[
 \big\|p(1-\mathrm{e}^{\lambda v/n})\big\|_{\mathcal D_h}
 \leq\|p\mathbf1_{\{p>K\}}\|_{\mathcal D_h}+C_hK\lambda,
\]
and taking the infimum over $K$ proves \eqref{eq:uniform-root-damping}. By \cref{prop:root-uniform-tails}, and since the $L^n$ tails of $p$ tend to zero, $\|p\mathbf1_{\{p>K\}}\|_{\mathcal D_h}\to0$ as $K\to\infty$. Choosing $K$ first and then letting $\lambda\downarrow0$ shows that $\eta_p(\lambda)\to0$.

Now set $q_\lambda=p\,\mathrm{e}^{\lambda v_\lambda/n}$, so that $q_\lambda^n=g_\lambda=\mathrm{e}^{\lambda v_\lambda}f$, with $g_\lambda$ as in the proof of \cref{cor:individual-positive-threshold-exponential}. Since $0\leq q_\lambda\leq p$ and $\rho_h(p-q_\lambda)\leq\eta_p(\lambda)$ by \eqref{eq:uniform-root-damping}, the triangle inequality and monotonicity give
\[
 R-\eta_p(\lambda)\leq\rho_h(q_\lambda)\leq R.
\]
As in the proof of \cref{cor:individual-positive-threshold-exponential}, $b\leq c_\lambda\leq b_h(g_\lambda)=\rho_h(q_\lambda)^{-n}$, which yields the scalar bound. For the measures, the factorization of $p^n-q_\lambda^n$ gives $0\leq p^n-q_\lambda^n\leq np^{n-1}(p-q_\lambda)$, and H\"older's inequality implies
\[
 \|f-g_\lambda\|_{L^1}
 \leq n\|p\|_{L^n}^{n-1}\|p-q_\lambda\|_{L^n}
 \leq n\|p\|_{L^n}^{n-1}\eta_p(\lambda).
\]
Since $\mathrm{MA}^{\mathrm{ext}}_h(v_\lambda)=c_\lambda g_\lambda\,\mathrm{d}V$, the identity $c_\lambda g_\lambda-bf=(c_\lambda-b)g_\lambda-b(f-g_\lambda)$, together with $\|g_\lambda\|_{L^1}\leq\|f\|_{L^1}=\|p\|_{L^n}^n$, proves the total-variation bound.

Finally, suppose $p\in L^r$ with $n<r<+\infty$, and choose $n<s<r$. By \eqref{eq:root-subcritical-bound} with exponent $s$, and since $\|\cdot\|_{L^n}\leq\|\cdot\|_{L^s}$, we have $\|\cdot\|_{\mathcal D_h}\leq(1+C_s)\|\cdot\|_{L^s}$. H\"older's inequality and \cref{lem:uniform-exponential} then give, for every normalized $v$,
\[
 \|p(-v)\|_{\mathcal D_h}
 \leq(1+C_s)\|p(-v)\|_{L^s}
 \leq(1+C_s)\|p\|_{L^r}\|v\|_{L^{rs/(r-s)}}
 \leq nC_{h,r}\|p\|_{L^r}.
\]
Since $0\leq1-\mathrm{e}^{\lambda v/n}\leq\lambda(-v)/n$, the left side of \eqref{eq:uniform-root-damping} is therefore at most $C_{h,r}\|p\|_{L^r}\lambda$, and the scalar and total-variation estimates can be repeated with this bound in place of $\eta_p(\lambda)$. For $\lambda$ so small that $C_{h,r}\|p\|_{L^r}\lambda\leq R/2$, the function $t\mapsto(R-t)^{-n}$ has bounded derivative on $[0,R/2]$, so the scalar error $c_\lambda-b$, and hence the total-variation error, is $O(\lambda)$.
\end{proof}

For a general $p\in\mathcal D_h^0$ the modulus $\eta_p(\lambda)$ tends to zero, but it need not be $O(\lambda)$. The rescaled scalar error $\lambda^{-1}\log(c_\lambda/b)$ is the loss in the lower bound $v_\lambda\geq v-\lambda^{-1}\log(c_\lambda/b)$ for normalized threshold subsolutions $v$, and the scalar bound of the proposition gives
\[
 0\leq\lambda^{-1}\log(c_\lambda/b)
 \leq-\frac n\lambda\log\bigl(1-\eta_p(\lambda)/R\bigr),
\]
which stays bounded as $\lambda\downarrow0$ when $\eta_p(\lambda)=O(\lambda)$, in particular for $p\in L^r$ with $r>n$; root tails alone do not give such a bound. Indeed, on the products of \cref{sec:zero-lower-mass-product-equations} there are densities whose roots lie in $\mathcal D_h^0$ and whose rescaled errors $\lambda^{-1}\log(c_\lambda/b_h(f))$ tend to $+\infty$ (\cref{sec:nonlp-selection}); for their roots $\eta_p(\lambda)/\lambda$ is unbounded. Even when a uniform lower barrier for $v_\lambda$ is available, passing to the limit in capacity requires the singularity-type or lower-mass hypotheses of a continuity theorem. The proposition estimates only scalars and measures; when the lower mass vanishes, it does not give an exact equation at the threshold.


\Needspace{12\baselineskip}
\subsection{Exponential equations at zero threshold}
\label{subsec:zero-threshold-product-exponential}

Let $(X_0,h_0)$ be a connected compact Hermitian manifold of dimension $m\geq1$ with $v_-(h_0)=0$. Put $Y=\mathbf P^k$, $k\geq1$, with normalized Fubini--Study form $\kappa$, and set $X=X_0\times Y$, $H=h_0+\kappa$, $N=m+k$ and
\[
 C=\binom Nm,\qquad V_0=\int_{X_0}h_0^m,\qquad
 \mu=V_0^{-1}h_0^m\wedge\kappa^k,\qquad R=\mu/\mathrm dV.
\]
Here $\mathrm dV$ is any smooth strictly positive probability volume form on $X$, and factor pullbacks are understood. The potential $u_\alpha=1-(1-\psi)^\alpha$ and its density $q_\alpha$ are those of \cref{lem:projective-sublogarithmic-density}, for fixed $0<\alpha<1$ and pole $o=[1:0:\cdots:0]$. Its exponential weight can absorb a concentration. Choosing the projective components accordingly gives one zero-threshold density with genuine exponential solutions for every positive exponent. On $Y\setminus\{o\}$ put $L=1-\psi\geq1$, $t=\mathrm e^{1-L}$ and
\[
 a=\alpha L^{\alpha-1},\qquad Q=t+(1-t)a,\qquad
 D=t(1-a)+(1-t)\alpha(1-\alpha)L^{\alpha-2},
\]
\[
 W=kQ^{k-1}D,\qquad
 B_\tau=q_\alpha\mathrm e^{\tau(L^\alpha-1)}\quad(\tau>0),\qquad
 b_0=(1-\alpha)^k.
\]

\begin{lemma}\label{lem:projective-exponential-barrier-growth}
The radial expressions are
\[
 q_\alpha=Q^{k-1}D\,\mathrm e^{k(L-1)},\qquad q_\alpha\geq b_0,
 \qquad W(L)\sim k\alpha^k(1-\alpha)L^{-1-k(1-\alpha)}.
\]
For every $\tau>0$, $\mathrm e^{\tau(L^\alpha-1)}W(L)\to+\infty$ as $L\to+\infty$. In particular, $B_\tau\kappa^k$ has infinite mass on every punctured neighbourhood of $o$.
\end{lemma}
\begin{proof}
In the radial calculation of \cref{lem:projective-sublogarithmic-density}, $t=r/(1+r)$ and $-dQ/dL=D$. The Fubini--Study radial probability is $k\mathrm e^{-k(L-1)}\mathrm dL$, whereas the radial mass of $q_\alpha\kappa^k$ is $W(L)\mathrm dL$. Their ratio gives the formula for $q_\alpha$. The lower bound follows from $\kappa+\ddc u_\alpha\geq(1-\alpha)\kappa$ away from the pole. Since $Q\sim\alpha L^{\alpha-1}$ and $D\sim\alpha(1-\alpha)L^{\alpha-2}$, the displayed asymptotic follows. The logarithm of the weighted radial density is $\tau L^\alpha-[1+k(1-\alpha)]\log L+O(1)$, which tends to infinity. Its integral over every tail is therefore infinite.
\end{proof}

\begin{theorem}\label{thm:product-zero-threshold-exponential-criterion}\index{Monge--Ampere equation@Monge--Amp\`ere equation!zero threshold criterion}
Let $f\geq0$ be nonzero and integrable on $X$, and put $F=f/R$. If $\lambda>0$ and $F\leq AB_\lambda$ almost everywhere for some finite $A>0$, then
\[
 s_\lambda=u_\alpha+\lambda^{-1}\log(CV_0/A)
\]
is a genuine exponential subsolution of actual mass $CV_0$. The exponential equation has a greatest genuine solution of finite positive mass. This conclusion holds for every $\lambda>0$ if $F$ is locally essentially bounded off $X_0\times\{o\}$ and
\[
 \lim_{T\to+\infty}\operatorname*{ess\,sup}_{x\in X_0,\ L(y)>T}
 \frac{\log^+(F(x,y)/q_\alpha(y))}{L(y)^\alpha}=0,
 \qquad \log^+t=\log\max(1,t).
\]
\end{theorem}
\begin{proof}
The product measure in \cref{thm:zero-lower-mass-product-equation} is $\mathrm{MA}^{\mathrm{ext}}_H(u_\alpha)=CV_0q_\alpha\mu$. Since $\mathrm e^{\lambda u_\alpha}B_\lambda=q_\alpha$ off the pole,
\[
 \mathrm e^{\lambda s_\lambda}f\,\mathrm dV
 \leq(CV_0/A)\mathrm e^{\lambda u_\alpha}AB_\lambda\mu
 =\mathrm{MA}^{\mathrm{ext}}_H(s_\lambda).
\]
Both measures vanish on the polar fibre. Adding the constant preserves admissibility and mass, and \cref{thm:individual-exponential-subsolution} gives the greatest exact solution. Its mass is finite because its supremum is finite and $f\in L^1$; it is positive because a genuine quasi-plurisubharmonic function is finite almost everywhere and $f$ has positive integral.

Under the limit condition, $F/q_\alpha\leq\mathrm e^{(\lambda/2)L^\alpha}$ almost everywhere sufficiently near the pole. On the remaining compact set the ratio is bounded. Increasing a finite constant $A$ therefore gives $F\leq AB_\lambda$ on all of $X$.
\end{proof}

\begin{theorem}\label{thm:zero-threshold-solvable-concentration}\index{Monge--Ampere density@Monge--Amp\`ere density!solvable concentration}\index{Subsolution!zero threshold exponential equation}
Let $F_j$ be any smooth strictly positive probability densities relative to $\mathrm dV_0=V_0^{-1}h_0^m$ satisfying $c_{h_0}(F_j)\leq2^{-2j}/C$. There are smooth strictly positive probability densities $g_j$ relative to $\kappa^k$ such that
\[
 f_j=RF_jg_j,\qquad f=\sum_{j\geq1}2^{-j}f_j,
 \qquad c_H(f_j)\leq2^{-2j},\qquad b_H(f)=0.
\]
The function $f$ is finite everywhere, smooth off the polar fibre, has integral one and is bounded below by a positive constant. This single density has a genuine exponential subsolution and a greatest finite-mass exact solution for every $\lambda>0$. Every genuine exact solution is unbounded.
\end{theorem}
\begin{proof}
Put $w_j=2^{-j}$ and $M_j=\max F_j$. Using an increasing sequence of smooth coordinate seminorms on $X_0$, choose $0<\varepsilon_j\leq1/2$ so that
\[
 w_j\varepsilon_jM_j\leq b_0 2^{-j-2},\qquad
 w_j\varepsilon_j\|F_j\|_{C^j}\leq2^{-j}.
\]
Then $E=\sum_jw_j\varepsilon_jF_j$ is smooth and satisfies $0<E\leq b_0/4$.

Choose a smooth function $0\leq\rho\leq1$ on $\mathbb R$, supported in $[0,1]$, positive inside and flat at the endpoints. Starting with $A_0=0$, inductively choose $A_j\geq2$, $A_j>A_{j-1}+2$, with
\[
 Z_j=\int_{A_j}^{A_j+1}\rho(L-A_j)
       \mathrm e^{(L^\alpha-1)/j}W(L)\,\mathrm dL
       \geq2w_jM_j.
\]
For each fixed $j$, \cref{lem:projective-exponential-barrier-growth} makes the integrand's weighted radial factor uniformly large on sufficiently distant intervals. Thus such a choice is possible. On that annulus put $b_j=\rho(L-A_j)B_{1/j}/Z_j$, extend it by zero elsewhere and set $g_j=\varepsilon_j+(1-\varepsilon_j)b_j$. Flatness makes $b_j$ smooth across the annulus boundaries; it vanishes near the pole and the infinity hyperplane. The integral of $b_j$ is one, so $g_j$ is a smooth strictly positive probability density. Moreover,
\[
 w_jM_j(1-\varepsilon_j)b_j\leq B_{1/j}/2.
\]

Apply \cref{thm:tw-smooth} on $Y$ to solve $(\kappa+\ddc\phi_j)^k=g_j\kappa^k$ with a smooth positive form; closed integration fixes its scalar as one. On $X_0$ the same theorem gives a smooth potential $v_j$ with $(h_0+\ddc v_j)^m=c_{h_0}(F_j)F_j\,\mathrm dV_0$. The product potential $v_j+\phi_j$ satisfies
\[
 (H+\ddc(v_j+\phi_j))^N
 =C c_{h_0}(F_j)F_jg_j\mu
 =C c_{h_0}(F_j)f_j\,\mathrm dV.
\]
The factor dimensions leave only the displayed binomial term. Scalar uniqueness gives $c_H(f_j)=C c_{h_0}(F_j)\leq2^{-2j}$.

Tonelli's theorem gives $\int_Xf\,\mathrm dV=\sum_jw_j=1$. The annuli are disjoint and escape to the pole, so only finitely many bumps meet any compact set outside the polar fibre. Hence $F=E+\sum_jw_j(1-\varepsilon_j)F_jb_j$ is smooth there. At the fibre the pointwise sum equals $E$, so $f$ is finite everywhere. Also $f\geq w_1f_1\geq\frac12\min_Xf_1>0$. If a positive coefficient $a$ were admitted by the maximal-subsolution characterization in \cref{thm:individual-density-subsolution}, then comparison with each smooth component would give
\[
 aw_j\leq c_H(f_j)\leq2^{-2j},\qquad a\leq2^{-j}
 \quad\text{for every }j.
\]
This is impossible; thus $b_H(f)=0$.

Fix $\lambda>0$ and choose $J$ with $1/j\leq\lambda$ for $j\geq J$. Since $u_\alpha\leq0$, one has $B_{1/j}\leq B_\lambda$ for these indices. Disjointness bounds the sum of their contributions by $B_\lambda/2$, and the baseline is at most $B_\lambda/4$. The remaining finitely many bumps have a finite bound
\[
 K_\lambda=\max_Y\sum_{j<J}w_j(1-\varepsilon_j)M_jb_j,
 \qquad D_\lambda=3/4+K_\lambda/b_0.
\]
Thus $F\leq D_\lambda B_\lambda$. The preceding criterion gives the explicit subsolution $u_\alpha+\lambda^{-1}\log(CV_0/D_\lambda)$ and the greatest exact solution. A bounded exponential subsolution would give a positive constant coefficient for $f$, contradicting its zero threshold. Every genuine exact solution is therefore unbounded.
\end{proof}

The constructed density also satisfies the logarithmic criterion: on the $j$th annulus, $F/q_\alpha\leq1/4+\frac12\mathrm e^{(L^\alpha-1)/j}$, and away from all annuli the ratio is at most $1/4$. Approaching the pole through annuli forces $j\to\infty$, giving the required uniform limit. Since $b_H(f)=0$, \cref{prop:individual-normalizer-lp} also shows that $f$ belongs to no $L^p$ for $p>1$.

For the elliptic threefold, \cref{thm:vanishing-lower-mass} supplies smooth positive forms with masses $\delta_\ell\to0$. Normalize their volume densities by $\delta_\ell$ and choose a subsequence with $\delta_{\ell_j}\leq2^{-2j}/C$. Scalar uniqueness gives the required $F_j$ with $c_{h_0}(F_j)=\delta_{\ell_j}$. Thus the construction applies to this threefold times $\mathbf P^k$ for every $k\geq1$. By \cref{prop:lower-mass-balanced-metric}, $V_0=24\pi$, so the explicit subsolutions have mass $96\pi$ on the fourfold and $240\pi$ on the fivefold. These numerical values concern the subsolutions; the greatest exact solutions have their own finite positive masses.

The projective components are chosen after the base sequence. The theorem does not determine solvability for an arbitrary previously prescribed concentration family or give a necessary criterion at zero threshold.

\Needspace{12\baselineskip}
\subsection{Solvable exponents and collapse at zero threshold}

Fix a nonzero nonnegative $f\in L^1(\mathrm dV)$ on a connected compact Hermitian manifold $(X,h)$. Let $\Lambda_h(f)$ consist of the positive exponents for which
\[
 \mathrm{MA}^{\mathrm{ext}}_h(u)=\mathrm e^{\lambda u}f\,\mathrm dV
\]
has a genuine $h$-plurisubharmonic solution. Existence of an exponential subsolution is equivalent to membership by \cref{thm:individual-exponential-subsolution}.

\Needspace{9\baselineskip}
\begin{proposition}\label{prop:solvable-exponents-upper-interval}\index{Monge--Ampere equation@Monge--Amp\`ere equation!solvable exponents}
If $\lambda_0\in\Lambda_h(f)$, then $[\lambda_0,+\infty)\subset\Lambda_h(f)$. Consequently this set is empty, equals $(a,+\infty)$ for some $a\geq0$, or equals $[a,+\infty)$ for some $a>0$. If $b_h(f)>0$, then $\Lambda_h(f)=(0,+\infty)$.
\end{proposition}
\begin{proof}
Let $u$ solve the equation at $\lambda_0$, and put $s=u-\max(0,\sup_Xu)$. Then $s\leq0$ and
\begin{align*}
 \mathrm{MA}^{\mathrm{ext}}_h(s)
 &=\mathrm e^{\lambda_0u}f\,\mathrm dV\\
 &\geq\mathrm e^{\lambda_0s}f\,\mathrm dV
 \geq\mathrm e^{\lambda s}f\,\mathrm dV
 \qquad(\lambda\geq\lambda_0).
\end{align*}
The greatest-solution theorem supplies an exact solution at each such $\lambda$. Every nonempty upper set of positive real numbers has the stated form, determined by its infimum and whether that infimum is attained. The positive-threshold case follows from \cref{cor:individual-positive-threshold-exponential}.
\end{proof}

\Needspace{15\baselineskip}
\begin{proposition}\label{prop:zero-threshold-exponential-mass-collapse}\index{Normalization scalar!zero threshold collapse}\index{Monge--Ampere mass@Monge--Amp\`ere mass!zero threshold exponential limit}
Suppose $b_h(f)=0$. Let $\lambda_j>0$ decrease to zero, and let $u_j$ be any genuine exact solution at $\lambda_j$. Put
\[
 M_j=\sup_Xu_j,\qquad v_j=u_j-M_j,\qquad
 c_j=\mathrm e^{\lambda_jM_j}.
\]
Then
\begin{gather*}
 c_j\longrightarrow0,\qquad \lambda_jM_j\longrightarrow-\infty,\\
 \|\mathrm{MA}^{\mathrm{ext}}_h(u_j)\|_{\mathrm{TV}}
 =\|\mathrm{MA}^{\mathrm{ext}}_h(v_j)\|_{\mathrm{TV}}
 \longrightarrow0.
\end{gather*}
Moreover $\mathrm e^{\lambda_ju_j}\to0$ uniformly on $X$, with value zero at poles, and $u_j\to-\infty$ uniformly in the sense that $\sup_Xu_j\to-\infty$.
\end{proposition}
\begin{proof}
For $g_j=\mathrm e^{\lambda_jv_j}f$, \cref{lem:normalized-exponential-npp} applied to the fixed finite non-pluripolar measure $f\,\mathrm dV$ gives $g_j\to f$ in $L^1$. Each $g_j$ is nonzero, because $v_j$ is finite almost everywhere. Constant invariance gives
\[
 \mathrm{MA}^{\mathrm{ext}}_h(v_j)=c_jg_j\,\mathrm dV,
 \qquad 0<c_j\leq b_h(g_j)
\]
by \cref{thm:individual-density-subsolution}. Upper semicontinuity, \cref{thm:individual-threshold-semicontinuity}, now gives $\limsup_jc_j\leq b_h(f)=0$. Since $0\leq g_j\leq f$, the total mass is at most $c_j\|f\|_1$ and tends to zero. Also $\lambda_jM_j=\log c_j\to-\infty$ and $0\leq\mathrm e^{\lambda_ju_j}\leq c_j$. For all sufficiently large $j$, $\lambda_j<1$ and $\log c_j<0$, so $M_j=(\log c_j)/\lambda_j\leq\log c_j\to-\infty$.
\end{proof}

For the tailored density of \cref{thm:zero-threshold-solvable-concentration}, solutions exist for every positive exponent, so the collapse applies as $\lambda\downarrow0$ to any choice of genuine exact solutions, including the greatest ones. The subsolutions in \cref{thm:product-zero-threshold-exponential-criterion} retain actual mass $CV_0$ for every exponent, while the exact-solution masses tend to zero. Order of potentials does not give order of their total Monge--Amp\`ere masses on a nonclosed background. Convergence of normalized potentials in capacity and an exact equation for a normalized cluster limit remain separate questions at exponent zero. Positive exponent endpoints are governed instead by the compactness criterion below.

\Needspace{12\baselineskip}
\subsection{Positive endpoints and right limits}

For $\lambda\in\Lambda_h(f)$, write $U_\lambda$ for the greatest solution of \cref{thm:individual-exponential-subsolution}. A bounded truncated-density equation gives a common upper barrier when the exponents stay away from zero. Greatest-subsolution maximality then supplies an order estimate, which controls positive endpoints and right limits.

\Needspace{12\baselineskip}
\begin{lemma}\label{lem:positive-exponent-uniform-upper-barrier}\index{Subsolution!common upper barrier}\index{Monge--Ampere equation@Monge--Amp\`ere equation!positive exponent order}
For every $\tau>0$, there is a nonnegative $t_\tau\in\PSH(X,h)\cap C(X)$ such that every genuine exponential subsolution at an exponent $\lambda\geq\tau$ lies below $t_\tau$. Put $K_\tau=\sup_Xt_\tau$. For $\tau\leq\lambda_1\leq\lambda_2$ in $\Lambda_h(f)$,
\[
 U_{\lambda_2}\geq U_{\lambda_1}
             -K_\tau\frac{\lambda_2-\lambda_1}{\lambda_2}.
\]
Consequently $U_\lambda+K_\tau\log(\lambda/\tau)$ is increasing on $\Lambda_h(f)\cap[\tau,+\infty)$.
\end{lemma}
\begin{proof}
Set $q=\min(f,1)$. This is a bounded nonzero density, so \cref{thm:nguyen-exp} gives a continuous solution $w$ of $\mathrm{MA}_h(w)=\mathrm e^{\tau w}q\,\mathrm dV$. Put $t_\tau=w+\max(0,-\inf_Xw)$. Then $t_\tau\geq0$ and, for $\lambda\geq\tau$,
\[
 \mathrm{MA}_h(t_\tau)=\mathrm e^{\tau w}q\,\mathrm dV
       \leq\mathrm e^{\lambda t_\tau}f\,\mathrm dV.
\]
If $u$ is a genuine subsolution at $\lambda$, then on $\{t_\tau<u-\epsilon\}$,
\[
 \mathrm{MA}_h(t_\tau)
       \leq\mathrm e^{-\lambda\epsilon}
                  \mathrm{MA}^{\mathrm{ext}}_h(u-\epsilon).
\]
The bounded first potential admits a genuine $h/2$-envelope, and $u-\epsilon\preceq t_\tau$. Thus \cref{thm:singular-strict-domination} gives $t_\tau\geq u-\epsilon$. Let $\epsilon\downarrow0$.

For the order estimate, put $\delta=K_\tau(\lambda_2-\lambda_1)/\lambda_2$. Since $U_{\lambda_1}\leq K_\tau$, constant invariance gives
\begin{align*}
 \mathrm{MA}^{\mathrm{ext}}_h(U_{\lambda_1}-\delta)
   &=\mathrm e^{\lambda_1U_{\lambda_1}}f\,\mathrm dV\\
   &\geq\mathrm e^{\lambda_2(U_{\lambda_1}-\delta)}f\,\mathrm dV.
\end{align*}
Greatest-subsolution maximality yields $U_{\lambda_2}\geq U_{\lambda_1}-\delta$. Finally $\log x\geq1-x^{-1}$ for $x\geq1$, which proves the logarithmically corrected order.
\end{proof}

\Needspace{16\baselineskip}
\begin{proposition}\label{prop:positive-exponent-endpoint-attainment}\index{Monge--Ampere equation@Monge--Amp\`ere equation!positive endpoint}
Suppose $\Lambda_h(f)\ne\varnothing$ and $a=\inf\Lambda_h(f)>0$. Then $a\in\Lambda_h(f)$ if and only if there are exponents $\lambda_j>a$ decreasing to $a$ and genuine exact solutions $u_j$ at $\lambda_j$ with $\inf_j\sup_Xu_j> -\infty$.

If $a\notin\Lambda_h(f)$, every sequence of genuine exact solutions at exponents $\lambda_j>a$ decreasing to $a$ satisfies
\begin{gather*}
 \sup_Xu_j\longrightarrow-\infty,\qquad
 \mathrm e^{\lambda_ju_j}\longrightarrow0\ \text{uniformly on }X,\\
 \|\mathrm{MA}^{\mathrm{ext}}_h(u_j)\|_{\mathrm{TV}}
                \longrightarrow0.
\end{gather*}
The exponential has value zero at poles.
\end{proposition}
\begin{proof}
If $u_a$ is an exact solution at $a$, put $s=u_a-\max(0,\sup_Xu_a)$. Then $s\leq0$ and $\mathrm{MA}^{\mathrm{ext}}_h(s)\geq\mathrm e^{\lambda s}f\,\mathrm dV$ for every $\lambda\geq a$. Thus $U_\lambda\geq s$, which supplies the required sequence.

Conversely, the lemma with $\tau=a$ gives a common upper bound $K_a$ for $u_j$. Write $M_j=\sup_Xu_j$. Extract a subsequence for which $M_j\to M\in\mathbb R$ and, by \cref{thm:normalized-compactness}, $u_j-M_j\to v\in\PSH(X,h)$ in $L^1$. Then $u_j\to u=v+M$ in $L^1$, with a genuine limit and a common upper bound. After a further extraction the convergence holds almost everywhere. Since $\lambda_j\leq\lambda_1$ and $K_a\geq0$, dominated convergence gives
\[
 \mathrm e^{\lambda_ju_j}f\longrightarrow\mathrm e^{au}f
              \quad\text{in }L^1(\mathrm dV).
\]
Indeed all these densities are at most $\mathrm e^{\lambda_1K_a}f$. The lower-density closure \cref{lem:density-subsolution-lone-closure} gives $\mathrm{MA}^{\mathrm{ext}}_h(u)\geq\mathrm e^{au}f\,\mathrm dV$. The greatest-solution theorem now supplies an exact solution at $a$.

If $a$ is not admitted and the suprema of an exact solution sequence failed to tend to $-\infty$, a subsequence would satisfy the lower-bound condition just proved. Thus $M_j\to-\infty$. Eventually $M_j<0$; since $\lambda_j\geq a>0$,
\[
 0\leq\mathrm e^{\lambda_ju_j}\leq\mathrm e^{aM_j}
                 \longrightarrow0.
\]
The exact mass is at most $\mathrm e^{aM_j}\|f\|_1$, proving the last assertion.
\end{proof}

\Needspace{10\baselineskip}
\begin{proposition}\label{prop:greatest-exponential-right-continuity}\index{Monge--Ampere equation@Monge--Amp\`ere equation!right continuity}\index{Convergence!greatest exponential solutions}
For every $a\in\Lambda_h(f)$, the greatest solutions satisfy
\[
 U_\lambda\longrightarrow U_a
       \quad\text{in }L^1\text{ and in capacity}\quad(\lambda\downarrow a).
\]
Their Monge--Amp\`ere measures converge in total variation to $\mathrm e^{aU_a}f\,\mathrm dV$.
\end{proposition}
\begin{proof}
Put $K=K_a$ from the lemma. For any $\lambda_j>a$ decreasing to $a$, the potentials
\[
 V_j=U_{\lambda_j}+K\log(\lambda_j/a)
\]
decrease and lie above $U_a$. Their limit $V$ is genuine. Dominated convergence between the integrable bounds $U_a$ and $V_1$ gives $L^1$ convergence, and \cref{lem:monotone-psh-capacity} gives capacity convergence. Since the added constants tend to zero, $U_{\lambda_j}\to V$ in both topologies.

Apply \cref{thm:exponential-closedness} to the nonpositive potentials $U_{\lambda_j}-K$, the exponents $\lambda_j\to a>0$ and the densities $\mathrm e^{\lambda_jK}f\to\mathrm e^{aK}f$ in $L^1$. It gives $\mathrm{MA}^{\mathrm{ext}}_h(V)=\mathrm e^{aV}f\,\mathrm dV$, with total-variation convergence of the equation measures. Hence $V\leq U_a$ by greatest-solution maximality, while the lower bound gives $V\geq U_a$. Thus $V=U_a$. The argument applies to every decreasing sequence of exponents, giving the stated right limit.
\end{proof}

The endpoint test is expressed in terms of existing solutions. An intrinsic criterion on $f$, the existence or nonexistence of a positive excluded endpoint, and continuity from the left remain separate questions. The right limit selects greatest solutions; it does not identify limits of arbitrary solutions with other singularities.


The scalar estimates lead to the variational description in \cref{chap:normalization}. There, logarithmic variation of the density gives supporting probabilities for the normalization functional. Identifying these probabilities with adjoint measures at a singular solution requires the regularity established in the relevant equation regime.
